% Réactions nucléaires provoquées et radioprotection — Physique Tle C — Cours signé OnBuch+
% Programme officiel MINESEC. Compiler avec : tectonic P20-reactions-nucleaires-radioprotection.tex
\newcommand{\DOCMATIERE}{Physique}
\newcommand{\DOCNIVEAU}{Tle C}
\newcommand{\DOCMODULE}{Module 4 : Onde, matière et transformations nucléaires}
\newcommand{\DOCLECON}{P20}
\newcommand{\DOCTITRE}{Réactions nucléaires provoquées et radioprotection}
% OnBuch+ — préambule « cours signés » ENRICHI (compilé avec tectonic / XeLaTeX)
% Système visuel « Pop » d'OnBuch+ : fond crème (jamais blanc), bordures encre
% épaisses, ombres « dures » décalées, titres Archivo Black en capitales.
% Version MATHÉMATIQUES : graphes (pgfplots/tikz), boîtes pédagogiques
% (définition, propriété, méthode, attention, le savais-tu, exercice résolu,
% exercice, TP, bilan), page de garde et signature OnBuch+ ancrée en bas.
%
% Macros attendues dans main.tex AVANT \input{preamble} :
%   \DOCMATIERE  \DOCNIVEAU  \DOCMODULE  \DOCLECON (numéro)  \DOCTITRE
\documentclass[11pt]{article}

\usepackage{fontspec}
\usepackage[french]{babel}
\usepackage[a4paper,margin=2cm,top=3.4cm,bottom=2.8cm,headheight=1.4cm,footskip=1.3cm]{geometry}
\usepackage{amsmath,amssymb,amsfonts}
\usepackage{mathtools} % \xmapsto, \coloneqq... souvent utilisés par les modèles
\usepackage{cancel} % simplifications barrées (bilans de forces)
% \abs{z} (module, valeur absolue) et \norm{u} (norme), utilisés par les modèles
\providecommand{\abs}{}\let\abs\relax\DeclarePairedDelimiter{\abs}{\lvert}{\rvert}
\providecommand{\norm}{}\let\norm\relax\DeclarePairedDelimiter{\norm}{\lVert}{\rVert}
\usepackage[table]{xcolor} % [table] : fournit aussi \rowcolors (alternance de lignes)
\usepackage{pagecolor}
\usepackage{tikz}
\usetikzlibrary{arrows.meta,positioning,calc,shapes.geometric,shapes.misc,shapes.arrows,decorations.pathmorphing,decorations.markings,patterns,shadows,mindmap,trees,fit,backgrounds,matrix,intersections,angles,quotes}
\usepackage{pgfplots}
\usepackage[version=4]{mhchem} % équations nucléaires \ce{^{235}_{92}U}
\usepackage[european,siunitx]{circuitikz} % schémas électriques
\pgfplotsset{compat=1.18}
\usepgfplotslibrary{fillbetween} % aires entre courbes (calcul intégral)
\usepackage{enumitem}
\usepackage{fancyhdr}
% [most] charge toutes les bibliothèques tcolorbox (listings, minted, raster,
% documentation...) et gonfle inutilement la mémoire TeX sur un document long
% (~300 boîtes) : on ne charge que ce qui est réellement utilisé.
\usepackage[skins,breakable]{tcolorbox}
\usepackage{graphicx}
\usepackage{colortbl}
\usepackage{booktabs}
\usepackage{tabularx}
\usepackage{diagbox} % case d'en-tête diagonale des tableaux à double entrée
% Colonnes extensibles centrée / gauche / droite (C, L, R), souvent utilisées par les modèles
\newcolumntype{C}{>{\centering\arraybackslash}X}
\newcolumntype{L}{>{\raggedright\arraybackslash}X}
\newcolumntype{R}{>{\raggedleft\arraybackslash}X}
\usepackage{array}
\usepackage{multicol}
\usepackage{siunitx}
% parse-numbers=false : accepte aussi \SI{2,5 \times 10^{-3}}{...}
\sisetup{locale=FR,output-decimal-marker={,},group-minimum-digits=4,parse-numbers=false}
\DeclareSIUnit\ohm{\ensuremath{\symup{\Omega}}} % U+2126 absent de la police
\DeclareSIUnit\division{div} % réglages d'oscilloscope (ms/div, V/div)
\usepackage{qrcode}
% \degree : commande fréquemment utilisée par les modèles pour un angle ou une
% température en dehors de \SI{}{} (non fournie par les paquets chargés).
\providecommand{\degree}{^{\circ}}
% \qed (fin de démonstration), utilisé par les modèles sans amsthm
\providecommand{\qed}{\hfill$\square$}
% \barre : double barre des valeurs interdites dans un tableau de variations (tkz-tab)
\providecommand{\barre}{\ensuremath{\|}}
% environnement proof (amsthm non chargé) utilisé par les modèles
\ifdefined\proof\else\newenvironment{proof}[1][Démonstration]{\par\noindent\textit{#1.}\ }{\qed\par}\fi
\usepackage{lastpage}
\usepackage{hyperref}
\hypersetup{hidelinks}

% ============================================================
% Palette « Pop » OnBuch+ — mêmes hex que la classe P (plus_ui.dart)
% ============================================================
\definecolor{poporange}{HTML}{F59321}
\definecolor{popdark}{HTML}{E07A0C}
\definecolor{popcream}{HTML}{FFF6E8}
\definecolor{popink}{HTML}{1C1714}
\definecolor{poppurple}{HTML}{7A5AE0}
\definecolor{popgreen}{HTML}{2E9E6A}
\definecolor{poppink}{HTML}{E0457B}
\definecolor{popblue}{HTML}{2F7DE1}
\definecolor{popgold}{HTML}{C98A12}
\definecolor{popmuted}{HTML}{8A7F76}
\definecolor{popline}{HTML}{EADFD2}
\definecolor{popgray}{HTML}{8A7F76}
\colorlet{popgrey}{popgray}
% Alias tolérants (le modèle nomme parfois les couleurs différemment)
\colorlet{popwhite}{white}
\colorlet{popblack}{popink}
\colorlet{popred}{poppink}
\colorlet{popyellow}{popgold}
% Teintes claires pour fonds de tableaux / zones
\colorlet{popblueL}{popblue!10!white}
\colorlet{poporangeL}{poporange!14!white}
\colorlet{popgreenL}{popgreen!12!white}
\colorlet{poppurpleL}{poppurple!12!white}
\colorlet{poppinkL}{poppink!12!white}
\colorlet{popgoldL}{popgold!14!white}

\pagecolor{popcream}

% ============================================================
% Typographie
% ============================================================
\defaultfontfeatures{Ligatures=TeX}
\setmainfont{PlusJakartaSans-Medium.ttf}[Path=../../../fonts/,BoldFont=PlusJakartaSans-Bold.ttf]
% Maths en sans-serif (Fira Math) avec chiffres et lettres droites de Plus
% Jakarta, pour que les formules se fondent dans le texte.
\usepackage[math-style=ISO,bold-style=ISO]{unicode-math}
\setmathfont{FiraMath-Regular.otf}[Path=../../../fonts/]
\setmathfont{PlusJakartaSans-Medium.ttf}[Path=../../../fonts/,range={up/num,up/{Latin,latin},\mathup}]
\usepackage{icomma} % virgule décimale sans espace en mode math
% Caractères mathématiques Unicode absents des polices de texte (Plus Jakarta,
% Archivo Black) : rendus par leur équivalent mathématique, en texte comme en
% formule — sinon ils s'affichent en carré vide (ex. « Arithmétique dans ℤ »).
\usepackage{newunicodechar}
\newunicodechar{ℝ}{\ensuremath{\mathbb{R}}}
\newunicodechar{ℤ}{\ensuremath{\mathbb{Z}}}
\newunicodechar{ℂ}{\ensuremath{\mathbb{C}}}
\newunicodechar{ℕ}{\ensuremath{\mathbb{N}}}
\newunicodechar{ℚ}{\ensuremath{\mathbb{Q}}}
\newunicodechar{𝔻}{\ensuremath{\mathbb{D}}}
\newunicodechar{α}{\ensuremath{\alpha}}
\newunicodechar{β}{\ensuremath{\beta}}
\newunicodechar{γ}{\ensuremath{\gamma}}
\newunicodechar{δ}{\ensuremath{\delta}}
\newunicodechar{ε}{\ensuremath{\varepsilon}}
\newunicodechar{θ}{\ensuremath{\theta}}
\newunicodechar{λ}{\ensuremath{\lambda}}
\newunicodechar{μ}{\ensuremath{\mu}}
\newunicodechar{σ}{\ensuremath{\sigma}}
\newunicodechar{τ}{\ensuremath{\tau}}
\newunicodechar{φ}{\ensuremath{\varphi}}
\newunicodechar{ψ}{\ensuremath{\psi}}
\newunicodechar{ω}{\ensuremath{\omega}}
\newunicodechar{Σ}{\ensuremath{\Sigma}}
% majuscules produites par \MakeUppercase dans les titres (α-aminés -> Α)
\newunicodechar{Α}{\ensuremath{\alpha}}
\newunicodechar{Β}{\ensuremath{\beta}}
\newunicodechar{Γ}{\ensuremath{\Gamma}}
\newunicodechar{Δ}{\ensuremath{\Delta}}
\newunicodechar{Ε}{\ensuremath{\varepsilon}}
\newunicodechar{Θ}{\ensuremath{\Theta}}
\newunicodechar{Λ}{\ensuremath{\Lambda}}
\newunicodechar{Μ}{\ensuremath{\mu}}
\newunicodechar{Τ}{\ensuremath{\tau}}
\newunicodechar{Φ}{\ensuremath{\Phi}}
\newunicodechar{Ψ}{\ensuremath{\Psi}}
\newunicodechar{Ω}{\ensuremath{\Omega}}
\newunicodechar{↦}{\ensuremath{\mapsto}}
\newunicodechar{∈}{\ensuremath{\in}}
\newunicodechar{∉}{\ensuremath{\notin}}
\newunicodechar{∧}{\ensuremath{\wedge}}
\newunicodechar{⇔}{\ensuremath{\Leftrightarrow}}
\newunicodechar{⟺}{\ensuremath{\Longleftrightarrow}}
\newunicodechar{⇒}{\ensuremath{\Rightarrow}}
\newunicodechar{⟹}{\ensuremath{\Longrightarrow}}
\newunicodechar{∘}{\ensuremath{\circ}}
\newunicodechar{∪}{\ensuremath{\cup}}
\newunicodechar{∩}{\ensuremath{\cap}}
\newunicodechar{≡}{\ensuremath{\equiv}}
\newunicodechar{⊂}{\ensuremath{\subset}}
\newunicodechar{⊥}{\ensuremath{\perp}}
\newunicodechar{∃}{\ensuremath{\exists}}
\newunicodechar{∀}{\ensuremath{\forall}}
\newunicodechar{∛}{\ensuremath{\sqrt[3]{\,}}}
\newunicodechar{∜}{\ensuremath{\sqrt[4]{\,}}}
\newunicodechar{⃗}{\ensuremath{{}^{\to}}}
\newunicodechar{ⁿ}{\ifmmode^{n}\else\textsuperscript{\ensuremath{n}}\fi}
\newunicodechar{ˣ}{\ifmmode^{x}\else\textsuperscript{\ensuremath{x}}\fi}
\newunicodechar{ᵇ}{\ifmmode^{b}\else\textsuperscript{\ensuremath{b}}\fi}
\newunicodechar{ʳ}{\ifmmode^{r}\else\textsuperscript{\ensuremath{r}}\fi}
\newunicodechar{ᵘ}{\ifmmode^{u}\else\textsuperscript{\ensuremath{u}}\fi}
\newunicodechar{ʸ}{\ifmmode^{y}\else\textsuperscript{\ensuremath{y}}\fi}
\newunicodechar{ᵃ}{\ifmmode^{a}\else\textsuperscript{\ensuremath{a}}\fi}
\newunicodechar{ᵉ}{\ifmmode^{e}\else\textsuperscript{\ensuremath{e}}\fi}
\newunicodechar{ᵏ}{\ifmmode^{k}\else\textsuperscript{\ensuremath{k}}\fi}
\newunicodechar{ᵖ}{\ifmmode^{p}\else\textsuperscript{\ensuremath{p}}\fi}
\newunicodechar{ᴺ}{\ifmmode^{N}\else\textsuperscript{\ensuremath{N}}\fi}
\newunicodechar{ᶜ}{\ifmmode^{c}\else\textsuperscript{\ensuremath{c}}\fi}
\newunicodechar{ʰ}{\ifmmode^{h}\else\textsuperscript{\ensuremath{h}}\fi}
\newunicodechar{ᵠ}{\ifmmode^{q}\else\textsuperscript{\ensuremath{q}}\fi}
\newunicodechar{ₙ}{\ifmmode_{n}\else\textsubscript{\ensuremath{n}}\fi}
\newunicodechar{ₐ}{\ifmmode_{a}\else\textsubscript{\ensuremath{a}}\fi}
\newunicodechar{ᵢ}{\ifmmode_{i}\else\textsubscript{\ensuremath{i}}\fi}
\newunicodechar{ᵣ}{\ifmmode_{r}\else\textsubscript{\ensuremath{r}}\fi}
\newunicodechar{ₖ}{\ifmmode_{k}\else\textsubscript{\ensuremath{k}}\fi}
\newunicodechar{ᵦ}{\ifmmode_{\beta}\else\textsubscript{\ensuremath{\beta}}\fi}
\newunicodechar{ₓ}{\ifmmode_{x}\else\textsubscript{\ensuremath{x}}\fi}
\newfontfamily\popdisplay{ArchivoBlack-Regular.ttf}[Path=../../../fonts/]
\newfontfamily\poplogo{SpaceGrotesk-Bold.ttf}[Path=../../../fonts/]
\newfontfamily\popsemi{PlusJakartaSans-SemiBold.ttf}[Path=../../../fonts/]
\newfontfamily\popxbold{PlusJakartaSans-ExtraBold.ttf}[Path=../../../fonts/]
\newfontfamily\popxboldls{PlusJakartaSans-ExtraBold.ttf}[Path=../../../fonts/,LetterSpace=3.0]

\setlist[enumerate]{leftmargin=1.7em,itemsep=5pt,topsep=4pt}
\setlist[itemize]{leftmargin=1.7em,itemsep=4pt,topsep=4pt,label={\color{poporange}\textbullet}}
\setlength{\parindent}{0pt}
\setlength{\parskip}{5pt}
\renewcommand{\arraystretch}{1.25}

% pgfplots : style Pop par défaut
\pgfplotsset{
  every axis/.append style={axis line style={popink,line width=1pt},tick style={popink},
    grid style={popline,line width=0.5pt},label style={font=\small},tick label style={font=\footnotesize}},
  popcurve/.style={poporange,line width=1.8pt,no markers,smooth},
}
% Même style disponible dans un \draw TikZ simple (hors pgfplots)
\tikzset{popcurve/.style={poporange,line width=1.8pt}}
\tikzset{popgrid/.style={popline,very thin}}
\colorlet{popcurve}{poporange} % le nom du style est parfois utilisé comme couleur

% ============================================================
% Pill / squiggle
% ============================================================
\newcommand{\poppill}[3]{%
  \tikz[baseline=-0.55ex]\node[draw=popink,line width=1.2pt,rounded corners=2.6pt,%
    fill=#2,inner xsep=6.5pt,inner ysep=3pt]{%
    \popxboldls\fontsize{7}{7}\selectfont\color{#3}\MakeUppercase{#1}};%
}
\newcommand{\popsquiggle}[1]{%
  \tikz[baseline=-0.4ex]{
    \draw[#1,line width=2.6pt,line cap=round]
      plot [smooth] coordinates {(0,0.09) (0.34,0.01) (0.68,0.21) (1.02,0.01) (1.36,0.10)};
  }%
}
% Mot-clé surligné (marqueur orange sous le texte)
\newcommand{\cle}[1]{{\popxbold\color{popdark}#1}}

% ============================================================
% En-tête / pied
% ============================================================
\pagestyle{fancy}
\fancyhf{}
\renewcommand{\headrulewidth}{0pt}
\renewcommand{\footrulewidth}{0pt}
\lhead{\raisebox{-0.15ex}{\poplogo\fontsize{13}{13}\selectfont\color{popink}OnBuch\color{poporange}+}}
\rhead{\poppill{\DOCMATIERE\ \textbullet\ \DOCNIVEAU\ \textbullet\ Leçon \DOCLECON}{white}{popink}}
\lfoot{\popxbold\fontsize{7.6}{7.6}\selectfont\color{popmuted}\MakeUppercase{Cours signé OnBuch+}\ \textbullet\ {\popsemi\DOCTITRE}}
\rfoot{\tikz[baseline=-0.6ex]\node[rounded corners=8.5pt,fill=popink,minimum height=17pt,minimum width=17pt,inner xsep=5pt,inner sep=0]{\popxbold\fontsize{8}{8}\selectfont\color{popcream}\thepage\,/\,\pageref*{LastPage}};}

% ============================================================
% Titres
% ============================================================
\newcommand{\chaptitre}[2]{%
  \begin{tcolorbox}[enhanced,colback=poporange,colframe=popink,boxrule=2.6pt,arc=11pt,
      drop shadow=popink,left=15pt,right=15pt,top=12pt,bottom=11pt,before skip=3pt,after skip=18pt]
    \poppill{#2}{popink}{popcream}\par\vspace{9pt}
    {\raggedright\hyphenpenalty=10000\exhyphenpenalty=10000\popdisplay\fontsize{22}{24}\selectfont\color{popcream}\MakeUppercase{#1}\par}
  \end{tcolorbox}
}
% Section numérotée : pastille orange avec le numéro + titre Archivo
\newcounter{popsec}
\newcommand{\coursec}[1]{%
  \stepcounter{popsec}\setcounter{popsub}{0}%
  \par\vspace{16pt}\noindent
  \begin{minipage}{\linewidth}
  \tikz[baseline=-0.7ex]\node[draw=popink,line width=1.6pt,fill=poporange,rounded corners=4pt,minimum size=22pt,inner sep=0]{\popdisplay\fontsize{12}{12}\selectfont\color{popcream}\thepopsec};%
  \hspace{8pt}{\popdisplay\fontsize{15}{17}\selectfont\color{popink}\MakeUppercase{#1}}\par
  \vspace{4pt}\noindent{\color{poporange}\rule{36pt}{3.6pt}}
  \end{minipage}\par\nopagebreak\vspace{8pt}
}
\newcounter{popsub}
\newcommand{\courssub}[1]{%
  \stepcounter{popsub}%
  \par\vspace{9pt}\noindent{\popxbold\fontsize{11.5}{13}\selectfont\color{popink}{\color{poporange}\thepopsec.\thepopsub}\hspace{6pt}#1}\par\nopagebreak\vspace{4pt}
}

% ============================================================
% Boîtes pédagogiques — même recette : fond blanc, bordure épaisse colorée,
% ombre dure encre, badge plein. Titre optionnel [#1].
% ============================================================
\tcbset{popbase/.style={breakable,enhanced,colback=white,boxrule=2.3pt,arc=9pt,
  shadow={3pt}{-3pt}{0pt}{popink},left=14pt,right=14pt,top=11pt,bottom=11pt,before skip=11pt,after skip=15pt,
  fontupper=\popsemi\fontsize{10.6}{14.5}\selectfont\color{popink},
  fontlower=\popsemi\fontsize{10.6}{14.5}\selectfont\color{popink}}}
% Coche dessinée (le glyphe ✓ n'existe pas dans les polices du document)
\newcommand{\popcheck}[1]{\tikz[baseline=-0.5ex]\draw[#1,line width=1.6pt,line cap=round,line join=round](-0.13,0.0)--(-0.04,-0.09)--(0.13,0.1);}
\providecommand{\coche}{\popcheck{popgreen}}
\providecommand{\resultat}[1]{\par\smallskip\noindent\fcolorbox{popgreen}{popgreenL}{\parbox{\dimexpr\linewidth-2\fboxsep-2\fboxrule\relax}{\textbf{Résultat.} #1}}\par\smallskip}
\newcommand{\popboxhead}[3]{\poppill{#1}{#2}{white}\ifx\relax#3\relax\else\hspace{6pt}{\popxbold\fontsize{10.6}{12}\selectfont\color{popink}#3}\fi\par\vspace{7pt}}

\newtcolorbox{aretenir}[1][]{popbase,colframe=popgreen,before upper={\popboxhead{À retenir}{popgreen}{#1}}}
\newtcolorbox{exemplebox}[1][]{popbase,colframe=poporange,before upper={\popboxhead{Exemple}{poporange}{#1}}}
\newtcolorbox{definition}[1][]{popbase,colframe=popblue,before upper={\popboxhead{Définition}{popblue}{#1}}}
\newtcolorbox{propriete}[1][]{popbase,colframe=poppurple,before upper={\popboxhead{Propriété}{poppurple}{#1}}}
\newtcolorbox{methode}[1][]{popbase,colframe=popgold,before upper={\popboxhead{Méthode}{popgold}{#1}}}
\newtcolorbox{attention}[1][]{popbase,colframe=poppink,before upper={\popboxhead{Attention}{poppink}{#1}}}
\newtcolorbox{experience}[1][]{popbase,colframe=popblue,colback=popblueL,before upper={\popboxhead{Expérience}{popblue}{#1}}}
% « Le savais-tu ? » : carte violette pleine (contexte, culture, Cameroun)
\newtcolorbox{savaistu}[1][]{popbase,colback=poppurple,colframe=popink,
  fontupper=\popsemi\fontsize{10.4}{14.2}\selectfont\color{white},
  before upper={\poppill{Le savais-tu\,?}{white}{poppurple}\ifx\relax#1\relax\else\hspace{6pt}{\popxbold\color{white}#1}\fi\par\vspace{7pt}}}
% Exercice résolu : énoncé en haut, \tcblower puis la solution sur fond crème
\newcounter{popexo}\newcounter{popexr}
\newtcolorbox{exoresolu}[1][]{popbase,colframe=poporange,colbacklower=popcream,
  segmentation style={popink,line width=1.2pt,dashed},
  before upper={\stepcounter{popexr}\popboxhead{Exercice résolu \thepopexr}{poporange}{#1}},
  before lower={\poppill{Solution}{popink}{popcream}\par\vspace{6pt}}}
% Exercice (non résolu en place) — la correction est dans la partie Corrigés
\newtcolorbox{exercice}[1][]{popbase,colframe=popink,
  before upper={\stepcounter{popexo}\popboxhead{Exercice \thepopexo}{popink}{#1}}}
% Corrigé
\newtcolorbox{corrige}[1][]{popbase,colframe=popgreen,colback=popgreenL,
  before upper={\popboxhead{Corrigé}{popgreen}{#1}}}
% Objectifs / prérequis / bilan
\newtcolorbox{objectifs}{popbase,colframe=popink,colback=poporangeL,
  before upper={\popboxhead{Objectifs de la leçon}{popink}{}}}
\newtcolorbox{prerequis}{popbase,colframe=popink,colback=popblueL,
  before upper={\popboxhead{Prérequis}{popblue}{}}}
\newtcolorbox{bilan}{popbase,colframe=popink,colback=white,boxrule=2.8pt,
  before upper={\popboxhead{Fiche bilan}{poporange}{L'essentiel à connaître pour l'examen}}}
% Figure encadrée avec légende
\newtcolorbox{popfigure}[1][]{enhanced,breakable=false,colback=white,colframe=popline,boxrule=1.4pt,arc=8pt,
  left=10pt,right=10pt,top=10pt,bottom=8pt,before skip=10pt,after skip=12pt,halign=center,
  lower separated=false,halign lower=center,
  fontlower=\popsemi\fontsize{9}{11.5}\selectfont\color{popmuted},
  #1}
\newcommand{\legende}[1]{\par\vspace{6pt}{\popsemi\fontsize{9}{11.5}\selectfont\color{popmuted}#1}}

% ============================================================
% Signature OnBuch+
% ============================================================
\newcommand{\poplogoinline}[1]{{\poplogo\fontsize{#1}{#1}\selectfont\color{popink}OnBuch\color{poporange}+}}
\newcommand{\signatureonbuch}{%
  \begin{tcolorbox}[enhanced,colback=popink,colframe=popink,boxrule=2.2pt,arc=10pt,
      drop shadow=poporange,left=14pt,right=14pt,top=10pt,bottom=10pt,before skip=0pt,after skip=0pt]
    \tikz[baseline=-0.7ex]\node[circle,fill=popgreen,draw=popcream,line width=1.3pt,minimum size=22pt,inner sep=0]{\popcheck{white}};%
    \hspace{9pt}\begin{minipage}[c]{0.62\linewidth}
      {\popxbold\fontsize{12}{13}\selectfont\color{popcream}Signé\ \textbullet\ {\poplogo OnBuch\color{poporange}+}}\par\vspace{2pt}
      {\raggedright\popsemi\fontsize{8}{10.5}\selectfont\color{popline}\MakeUppercase{Équipe pédagogique OnBuch+}\\\MakeUppercase{Conforme au programme officiel MINESEC}\par}
    \end{minipage}\hfill
    {\poplogo\fontsize{22}{22}\selectfont\color{popcream}OnBuch\color{poporange}+}
  \end{tcolorbox}%
}

% ============================================================
% Page de garde
% #1 titre · #2 matière · #3 niveau · #4 module · #5 n° leçon · #6 durée ·
% #7 description / objectifs (texte ou itemize)
% ============================================================
\newcommand{\pagegarde}[7]{%
  \thispagestyle{empty}
  \newgeometry{a4paper,margin=1.7cm,top=1.6cm,bottom=1.4cm}%
  \begin{tikzpicture}[remember picture,overlay]
    \fill[poporange!18] ([xshift=-3.2cm,yshift=-3.6cm]current page.north east) circle (4.6cm);
    \fill[poppurple!14] ([xshift=2.4cm,yshift=7.6cm]current page.south west) circle (3.8cm);
  \end{tikzpicture}%
  {\poplogo\fontsize{30}{30}\selectfont\color{popink}OnBuch\color{poporange}+}\hfill
  \raisebox{0.6ex}{\poppill{Cours signé \textbullet\ Premium}{popink}{popcream}}\par
  \vspace{1pt}\popsquiggle{poppurple}\par
  \vspace{8pt}
  \begin{tcolorbox}[enhanced,colback=poporange,colframe=popink,boxrule=2.8pt,arc=13pt,
      drop shadow=popink,left=18pt,right=18pt,top=12pt,bottom=13pt,after skip=10pt]
    \poppill{#2 \textbullet\ #3}{popink}{popcream}\hspace{5pt}\poppill{#4}{white}{popink}\par\vspace{8pt}
    {\popdisplay\fontsize{14}{14}\selectfont\color{popink}LEÇON #5}\par\vspace{4pt}
    {\raggedright\hyphenpenalty=10000\exhyphenpenalty=10000\popdisplay\fontsize{25}{27}\selectfont\color{popcream}\MakeUppercase{#1}\par}
    \vspace{8pt}
    {\popxbold\fontsize{9.5}{9.5}\selectfont\color{popink}\MakeUppercase{Durée conseillée : #6}}
  \end{tcolorbox}
  \begin{tcolorbox}[enhanced,colback=white,colframe=popink,boxrule=2.2pt,arc=10pt,
      drop shadow=popink,left=16pt,right=16pt,top=10pt,bottom=10pt,after skip=8pt]
    \poppill{Dans cette leçon}{poppurple}{white}\par\vspace{5pt}
    \popsemi\fontsize{9.3}{12.6}\selectfont\color{popink}#7
  \end{tcolorbox}
  \noindent\begin{minipage}[t]{0.487\linewidth}
    \begin{tcolorbox}[enhanced,colback=popgreen,colframe=popink,boxrule=2.2pt,arc=10pt,
        drop shadow=popink,left=11pt,right=11pt,top=10pt,bottom=10pt,height=3.1cm,valign=center]
      \tcbox[colback=white,colframe=popink,boxrule=1.3pt,arc=5pt,boxsep=0pt,nobeforeafter,
        left=3pt,right=3pt,top=3pt,bottom=3pt,valign=center]{\qrcode[height=1.9cm]{https://play.google.com/store/apps/details?id=com.luvvix.onbuch&hl=fr}}%
      \hspace{8pt}\begin{minipage}[c]{0.5\linewidth}
        {\popdisplay\fontsize{11}{12.5}\selectfont\color{white}\MakeUppercase{Télécharge OnBuch}}\par\vspace{4pt}
        {\raggedright\popsemi\fontsize{8.4}{11}\selectfont\color{white}Gratuit \textbullet\ Google Play\\Tuteur IA, annales, concours, résultats\par}
      \end{minipage}
    \end{tcolorbox}
  \end{minipage}\hfill
  \begin{minipage}[t]{0.487\linewidth}
    \begin{tcolorbox}[enhanced,colback=poppurple,colframe=popink,boxrule=2.2pt,arc=10pt,
        drop shadow=popink,left=11pt,right=11pt,top=10pt,bottom=10pt,height=3.1cm,valign=center]
      \tikz[baseline=-0.7ex]\node[circle,fill=white,draw=popink,line width=1.3pt,minimum size=1.6cm,inner sep=0]{\popdisplay\fontsize{24}{24}\selectfont\color{poppurple}+};%
      \hspace{8pt}\begin{minipage}[c]{0.6\linewidth}
        {\popdisplay\fontsize{11}{12.5}\selectfont\color{white}\MakeUppercase{Passe à OnBuch+}}\par\vspace{4pt}
        {\raggedright\popsemi\fontsize{8.4}{11}\selectfont\color{white}Tous les cours signés, Tuteur IA illimité, fiches PDF — onglet OnBuch+ de l'app\par}
      \end{minipage}
    \end{tcolorbox}
  \end{minipage}
  \vspace{8pt}
  \popcontenu
  \vfill
  \signatureonbuch
  \restoregeometry
  \newpage
}

% Rangée « Ce cours contient » (page de garde)
\newcommand{\poptile}[3]{% #1 couleur #2 gros texte #3 légende
  \begin{tcolorbox}[enhanced,colback=#1,colframe=popink,boxrule=2pt,arc=9pt,shadow={2.5pt}{-2.5pt}{0pt}{popink},
      left=6pt,right=6pt,top=9pt,bottom=9pt,halign=center,valign=center,height=1.75cm,width=\linewidth,nobeforeafter]
    {\popdisplay\fontsize{12.5}{13}\selectfont\color{white}\MakeUppercase{#2}}\par\vspace{3pt}
    {\centering\popsemi\fontsize{7.8}{9.5}\selectfont\color{white}#3\par}
  \end{tcolorbox}}
\newcommand{\popcontenu}{%
  \par\noindent{\popxboldls\fontsize{7.5}{7.5}\selectfont\color{popmuted}CE COURS SIGNÉ CONTIENT}\par\vspace{7pt}
  \noindent\begin{minipage}[t]{0.235\linewidth}\poptile{popblue}{Cours}{complet, illustré, pas à pas}\end{minipage}\hfill
  \begin{minipage}[t]{0.235\linewidth}\poptile{poporange}{Méthodes}{et exercices résolus}\end{minipage}\hfill
  \begin{minipage}[t]{0.235\linewidth}\poptile{poppink}{Exercices}{type examen + corrigés}\end{minipage}\hfill
  \begin{minipage}[t]{0.235\linewidth}\poptile{popgreen}{Fiche bilan}{l'essentiel à retenir}\end{minipage}\par}

% Sommaire
\newtcolorbox{sommaire}{enhanced,colback=white,colframe=popink,boxrule=2.2pt,arc=10pt,
  drop shadow=popink,left=18pt,right=18pt,top=13pt,bottom=13pt,before skip=6pt,after skip=20pt,
  before upper={\poppill{Sommaire}{poppurple}{white}\par\vspace{9pt}\popsemi\fontsize{10.6}{14.5}\selectfont\color{popink}}}

% Signature de fin de cours — poussée tout en bas de la dernière page
\newcommand{\findecours}{%
  \par\vfill\nopagebreak
  \begin{center}\popsquiggle{poporange}\end{center}\vspace{4pt}
  \signatureonbuch
}

%% FIGFIX-BEGIN v5 — correctifs globaux des figures (docs/onbuch-generation/tools/figfix)
\usepackage{adjustbox}\usepackage{etoolbox}\tcbuselibrary{raster}
% toute figure TikZ/pgfplots plus large que la ligne est réduite à la largeur disponible
\BeforeBeginEnvironment{tikzpicture}{\begin{adjustbox}{max width=\linewidth}}
\AfterEndEnvironment{tikzpicture}{\end{adjustbox}}
% légendes pgfplots lisibles par-dessus les courbes
\pgfplotsset{every axis/.append style={legend style={fill=white,fill opacity=0.92,text opacity=1,draw=black!15,inner sep=3pt}}}
% étiquettes d'arêtes (syntaxe edge["..."]) avec fond blanc
\tikzset{every edge quotes/.append style={fill=white,inner sep=1.2pt}}
%% FIGFIX-END
\begin{document}

\pagegarde{Réactions nucléaires provoquées et radioprotection}{Physique}{Tle C}{Module 4 : Onde, matière et transformations nucléaires}{P20}{5 h}{Cette leçon étudie les réactions nucléaires provoquées (transmutation, fission, fusion), les lois de conservation qui régissent leurs équations et l'énergie libérée via E = Δm·c². Elle présente ensuite les effets des rayonnements sur l'organisme et les méthodes de radioprotection, dont la loi d'atténuation exponentielle d'un faisceau de photons.
\vspace{4pt}
\begin{multicols}{2}\small
\begin{itemize}
\item À la fin de cette leçon, je sais distinguer transmutation, fission et fusion nucléaires
\item À la fin de cette leçon, je sais écrire et équilibrer l'équation d'une réaction nucléaire en appliquant les lois de conservation de Soddy
\item À la fin de cette leçon, je sais calculer le défaut de masse Δm et l'énergie libérée E = Δm·c² d'une réaction nucléaire
\item À la fin de cette leçon, je sais expliquer le principe d'une réaction en chaîne, d'une centrale nucléaire et de la fusion dans les étoiles
\item À la fin de cette leçon, je sais distinguer irradiation et contamination et citer les effets des rayonnements sur l'organisme
\item À la fin de cette leçon, je sais définir la dose absorbée et la dose équivalente avec leurs unités
\end{itemize}\end{multicols}}

\begin{sommaire}
\begin{multicols}{2}
\begin{enumerate}
\item Réactions nucléaires provoquées : transmutation
\item Fission nucléaire et réaction en chaîne
\item Fusion nucléaire
\item Aspect énergétique : E = Δm·c²
\item Effets biologiques des rayonnements et dosimétrie
\item Radioprotection et atténuation d'un faisceau de photons
\item Activité d'intégration
\item Méthodes et exercices résolus
\item Exercices
\item Corrigés des exercices
\item Fiche bilan
\end{enumerate}
\end{multicols}
\end{sommaire}

\begin{objectifs}
\begin{itemize}
\item À la fin de cette leçon, je sais distinguer transmutation, fission et fusion nucléaires
\item À la fin de cette leçon, je sais écrire et équilibrer l'équation d'une réaction nucléaire en appliquant les lois de conservation de Soddy
\item À la fin de cette leçon, je sais calculer le défaut de masse Δm et l'énergie libérée E = Δm·c² d'une réaction nucléaire
\item À la fin de cette leçon, je sais expliquer le principe d'une réaction en chaîne, d'une centrale nucléaire et de la fusion dans les étoiles
\item À la fin de cette leçon, je sais distinguer irradiation et contamination et citer les effets des rayonnements sur l'organisme
\item À la fin de cette leçon, je sais définir la dose absorbée et la dose équivalente avec leurs unités
\item À la fin de cette leçon, je sais citer et justifier les trois méthodes de radioprotection (distance, temps, écran)
\item À la fin de cette leçon, je sais exploiter la loi N(x) = N0·e\^{}(−μx) et la couche de demi-atténuation pour dimensionner un écran de protection
\end{itemize}
\end{objectifs}

\begin{prerequis}
\begin{itemize}
\item Noyau atomique : notation ᴬzX, isotopes, nucléons (leçon sur la radioactivité)
\item Réactions nucléaires spontanées : radioactivité α, β⁻, β⁺ et lois de Soddy
\item Relation masse-énergie d'Einstein et unités de masse atomique (u, MeV/c²)
\item Fonction exponentielle et logarithme népérien (mathématiques)
\item Énergie, puissance et conversions d'unités (Première)
\end{itemize}
\end{prerequis}

\newpage

\coursec{Réactions nucléaires provoquées : transmutation}

En 2023, le débat sur l'électrification du Cameroun a relancé la question des centrales nucléaires : comment une poignée d'uranium peut-elle éclairer toute une ville ? À l'hôpital de Yaoundé, les techniciens de radiothérapie manipulent chaque jour des sources radioactives derrière des murs de plomb : pourquoi le plomb, et quelle épaisseur suffit-elle ? Dans le Soleil, des noyaux d'hydrogène fusionnent et libèrent l'énergie qui fait vivre la Terre. Comment une minuscule perte de masse peut-elle produire une énergie gigantesque ? Cette leçon répond à ces questions en étudiant les \cle{réactions nucléaires provoquées} et les moyens de s'en protéger. Nous commençons par la plus fondamentale d'entre elles : la \cle{transmutation}, c'est-à-dire l'art de transformer un noyau en un autre noyau.

\courssub{Radioactivité spontanée et réaction nucléaire provoquée}

\textbf{Observation.} Au chapitre précédent, nous avons étudié la \cle{radioactivité} : certains noyaux instables (uranium 238, radium 226, carbone 14\ldots) se désintègrent \emph{d'eux-mêmes}, sans aucune intervention extérieure, en émettant un rayonnement $\alpha$, $\beta^-$ ou $\gamma$. Rien ne peut ni déclencher, ni retarder, ni accélérer cette désintégration : ni la température, ni la pression, ni un champ électrique. C'est un phénomène \emph{spontané} et \emph{aléatoire}.

Mais l'histoire de la physique ne s'est pas arrêtée là. En 1919, Ernest Rutherford eut une idée audacieuse : au lieu d'\emph{attendre} qu'un noyau se transforme, pourquoi ne pas le \emph{frapper} avec un projectile pour le forcer à changer ? En bombardant de l'azote avec des particules $\alpha$, il observa l'apparition de protons : il venait de réaliser la première \emph{transmutation artificielle} de l'histoire. Le vieux rêve des alchimistes --- transformer un élément en un autre --- devenait réalité, non par magie, mais par la physique.

\begin{definition}[Réaction nucléaire provoquée et transmutation]
Une \cle{réaction nucléaire provoquée} est une transformation du noyau atomique déclenchée par l'impact d'un \cle{projectile} (neutron, proton, particule $\alpha$, noyau léger\ldots) sur un noyau \cle{cible}.

On appelle \cle{transmutation nucléaire} toute transformation au cours de laquelle un noyau est changé en un noyau d'\emph{un autre élément} (ou d'un autre isotope), que cette transformation soit spontanée (radioactivité) ou provoquée. Dans ce chapitre, nous nous intéressons aux transmutations \emph{provoquées}, que l'on symbolise par :
\[
\underbrace{\ce{^{A_1}_{Z_1}X}}_{\text{cible}} + \underbrace{\ce{^{A_2}_{Z_2}x}}_{\text{projectile}} \longrightarrow \underbrace{\ce{^{A_3}_{Z_3}Y}}_{\text{noyau fils}} + \underbrace{\ce{^{A_4}_{Z_4}y}}_{\text{particule éjectée}}
\]
\end{definition}

\begin{attention}[Spontanée ou provoquée ?]
Ne confonds pas les deux situations :
\begin{itemize}
\item \textbf{Réaction spontanée (radioactivité)} : un seul noyau à gauche de la flèche, il se désintègre seul ; elle obéit à la loi de décroissance $N(t) = N_0 e^{-\lambda t}$.
\item \textbf{Réaction provoquée} : \emph{deux} partenaires à gauche (cible $+$ projectile) ; elle ne se produit que si le projectile rencontre la cible avec une énergie suffisante.
\end{itemize}
Dans les deux cas, les mêmes lois de conservation s'appliquent.
\end{attention}

\courssub{Les grandes étapes historiques}

\textbf{1919 : la première transmutation artificielle (Rutherford).} Rutherford bombarda de l'azote gazeux avec des particules $\alpha$ (noyaux d'hélium \ce{^{4}_{2}He}) émises par une source radioactive. Il détecta l'émission de protons et comprit que le noyau d'azote s'était transformé en noyau d'oxygène :
\[
\ce{^{14}_{7}N + ^{4}_{2}He -> ^{17}_{8}O + ^{1}_{1}H}
\]
Pour la première fois, l'être humain transformait délibérément un élément chimique en un autre.

\textbf{1932 : la découverte du neutron (Chadwick).} En bombardant du béryllium avec des particules $\alpha$, James Chadwick mit en évidence un rayonnement très pénétrant, électriquement neutre : le \cle{neutron}. L'équation de la réaction s'écrit :
\[
\ce{^{9}_{4}Be + ^{4}_{2}He -> ^{12}_{6}C + ^{1}_{0}n}
\]
Le neutron, sans charge électrique, n'est pas repoussé par le noyau : c'est le projectile idéal pour provoquer des réactions nucléaires, comme nous le verrons avec la fission.

\textbf{1934 : la radioactivité artificielle (Irène et Frédéric Joliot-Curie).} En bombardant de l'aluminium avec des particules $\alpha$, les Joliot-Curie obtinrent du phosphore 30, un isotope \emph{qui n'existe pas dans la nature} et qui est radioactif :
\[
\ce{^{27}_{13}Al + ^{4}_{2}He -> ^{30}_{15}P + ^{1}_{0}n}
\]
Ils venaient d'inventer la \cle{radioactivité artificielle} : on peut désormais \emph{fabriquer} des isotopes radioactifs sur mesure.

\begin{savaistu}[Un Nobel aux conséquences immenses]
Irène et Frédéric Joliot-Curie reçurent le \textbf{prix Nobel de chimie en 1935} pour la découverte de la radioactivité artificielle. Cette découverte est à l'origine de toute la \cle{médecine nucléaire} : le technétium 99m utilisé en scintigraphie, l'iode 131 pour traiter la thyroïde, le fluor 18 des examens TEP\ldots sont tous des isotopes artificiels produits dans des réacteurs ou des cyclotrons. Quand un hôpital camerounais réalise une radiothérapie ou une imagerie médicale, il utilise directement l'héritage des Joliot-Curie.
\end{savaistu}

\courssub{Les lois de conservation dans une réaction nucléaire}

\textbf{Intuition.} Dans une réaction nucléaire, les nucléons ne disparaissent pas et ne surgissent pas de nulle part : ils se \emph{réarrangent}. Les protons et les neutrons présents avant la réaction se retrouvent après, répartis autrement entre les noyaux. Il en résulte deux lois simples et puissantes.

\begin{propriete}[Lois de conservation de Soddy généralisées]
Dans \emph{toute} réaction nucléaire (spontanée ou provoquée), il y a :
\begin{itemize}
\item \textbf{conservation du nombre de nucléons} : la somme des nombres de masse $A$ est la même avant et après ;
\item \textbf{conservation de la charge électrique} : la somme des numéros atomiques $Z$ est la même avant et après.
\end{itemize}
Pour la réaction $\ce{^{A_1}_{Z_1}X + ^{A_2}_{Z_2}x -> ^{A_3}_{Z_3}Y + ^{A_4}_{Z_4}y}$ :
\[
\boxed{A_1 + A_2 = A_3 + A_4} \qquad \text{et} \qquad \boxed{Z_1 + Z_2 = Z_3 + Z_4}
\]
Ces lois sont des \emph{lois expérimentales fondamentales} : aucune réaction nucléaire observée ne les viole. (On admettra aussi que l'énergie totale et la quantité de mouvement se conservent, ce qui servira à l'étude énergétique de la section 4.)
\end{propriete}

\begin{experience}[Vérification sur l'expérience de Rutherford]
Vérifions les lois de Soddy sur la réaction historique $\ce{^{14}_{7}N + ^{4}_{2}He -> ^{17}_{8}O + ^{1}_{1}H}$.
\begin{itemize}
\item \textbf{Nombre de nucléons} : avant, $14 + 4 = 18$ ; après, $17 + 1 = 18$. \checkmark
\item \textbf{Charge} : avant, $7 + 2 = 9$ ; après, $8 + 1 = 9$. \checkmark
\end{itemize}
Les deux bilans sont équilibrés : l'équation est correctement écrite. Cette vérification systématique est \emph{exigible au Baccalauréat} : toute équation nucléaire que tu écris doit être contrôlée ainsi.
\end{experience}

\begin{popfigure}
\begin{center}
\begin{tikzpicture}[scale=1,>=Stealth]
% Projectile alpha
\shade[ball color=poporange] (-4.2,0) circle (0.28);
\shade[ball color=poporange!80] (-3.6,0.12) circle (0.28);
\shade[ball color=poporange!70] (-3.9,-0.22) circle (0.28);
\shade[ball color=poporange!60] (-3.9,0.3) circle (0.28);
\node[below] at (-3.9,-0.75) {projectile $\alpha$ : $\ce{^{4}_{2}He}$};
\draw[->,very thick,popink] (-3.2,0) -- (-1.9,0);
% Cible
\shade[ball color=popblue] (0,0) circle (0.55);
\node[white] at (0,0) {\small $\ce{^{14}_{7}N}$};
\node[below] at (0,-0.85) {noyau cible (au repos)};
% Flèche réaction
\draw[->,ultra thick,popdark] (1.1,0) -- (2.4,0);
\node[above] at (1.75,0.05) {\small transmutation};
% Noyau fils
\shade[ball color=popgreen] (4.1,0.7) circle (0.6);
\node[white] at (4.1,0.7) {\small $\ce{^{17}_{8}O}$};
\draw[->,very thick,popgreen!60!black] (3.4,0.35) -- (2.9,0.1);
\node[right] at (4.75,0.7) {\small noyau fils};
% Proton éjecté
\shade[ball color=poppurple] (4.1,-0.85) circle (0.22);
\node[right] at (4.4,-0.85) {proton éjecté : $\ce{^{1}_{1}H}$};
\draw[->,very thick,poppurple] (3.5,-0.5) -- (3.85,-0.75);
\end{tikzpicture}
\end{center}
\legende{Schéma de la transmutation de Rutherford (1919) : la particule $\alpha$ frappe le noyau d'azote 14 ; il se forme un noyau d'oxygène 17 et un proton est éjecté. Bilan : $14+4=17+1$ et $7+2=8+1$.}
\end{popfigure}

\courssub{Méthode : équilibrer une équation nucléaire}

Au Baccalauréat, l'exercice classique consiste à \emph{identifier une particule inconnue} dans une équation. Voici la démarche infaillible.

\begin{methode}[Équilibrer une équation nucléaire en 3 étapes]
On cherche la particule inconnue $\ce{^{A}_{Z}X}$ dans une équation nucléaire.
\begin{enumerate}
\item \textbf{Écrire la conservation du nombre de nucléons} : somme des $A$ à gauche $=$ somme des $A$ à droite. On en tire la valeur de $A$.
\item \textbf{Écrire la conservation de la charge} : somme des $Z$ à gauche $=$ somme des $Z$ à droite. On en tire la valeur de $Z$.
\item \textbf{Identifier la particule} à l'aide du couple $(Z, A)$ : $Z = 0, A = 1$ donne le neutron $\ce{^{1}_{0}n}$ ; $Z = 1, A = 1$ donne le proton $\ce{^{1}_{1}H}$ ; $Z = 2, A = 4$ donne la particule $\alpha$ $\ce{^{4}_{2}He}$ ; $Z = -1, A = 0$ donne l'électron $\ce{^{0}_{-1}e}$ ; sinon, on consulte le tableau périodique pour trouver l'élément de numéro atomique $Z$.
\end{enumerate}
\end{methode}

\begin{exemplebox}[Identifier la particule X]
Identifions $X$ dans la réaction des Joliot-Curie : $\ce{^{27}_{13}Al + ^{4}_{2}He -> ^{30}_{15}P + X}$. Notons $X = \ce{^{A}_{Z}X}$.
\begin{align*}
\text{Conservation de } A : \quad & 27 + 4 = 30 + A \;\Rightarrow\; A = 1 \\
\text{Conservation de } Z : \quad & 13 + 2 = 15 + Z \;\Rightarrow\; Z = 0
\end{align*}
La particule de charge nulle et de nombre de masse 1 est le \textbf{neutron} : $X = \ce{^{1}_{0}n}$.
\end{exemplebox}

\begin{attention}[Ne confonds pas A et Z !]
L'erreur la plus fréquente consiste à mélanger les deux nombres :
\begin{itemize}
\item $Z$ (en \textbf{bas} à gauche) est le \textbf{numéro atomique} : c'est le nombre de \emph{protons}, donc la \emph{charge} du noyau. C'est lui qui identifie l'élément chimique.
\item $A$ (en \textbf{haut} à gauche) est le \textbf{nombre de masse} : c'est le nombre total de \emph{nucléons} (protons $+$ neutrons).
\end{itemize}
On équilibre les $A$ \emph{entre eux} (ligne du haut) et les $Z$ \emph{entre eux} (ligne du bas) --- jamais l'un avec l'autre ! Autre piège : le neutron $\ce{^{1}_{0}n}$ a bien $A = 1$ mais $Z = 0$ ; l'oublier dans le bilan de charge fausse tout le calcul.
\end{attention}

\courssub{Réaction chimique ou réaction nucléaire ?}

Pour bien situer la nouveauté de ce chapitre, comparons ce que tu connais depuis la classe de seconde (la réaction chimique) avec la réaction nucléaire.

\begin{center}
\begin{tabularx}{\linewidth}{|l|X|X|}
\hline
\rowcolor{popblueL} \textbf{Critère} & \textbf{Réaction chimique} & \textbf{Réaction nucléaire} \\
\hline
Siège de la transformation & Électrons périphériques & Noyau de l'atome \\
\hline
Éléments chimiques & Conservés (mêmes atomes) & Modifiés (transmutation) \\
\hline
Grandeurs conservées & Nombre d'atomes de chaque élément, charge & Nombre de nucléons $A$, charge $Z$ \\
\hline
Énergie mise en jeu & Quelques eV par atome ($\sim \SI{e2}{kJ/mol}$) & Quelques MeV par noyau ($\sim 10^{11}$ fois plus) \\
\hline
Influence de la température, de la pression & Importante & Quasi nulle (sauf fusion) \\
\hline
Exemple & $\ce{CH4 + 2O2 -> CO2 + 2H2O}$ & $\ce{^{14}_{7}N + ^{4}_{2}He -> ^{17}_{8}O + ^{1}_{1}H}$ \\
\hline
\end{tabularx}
\end{center}

\begin{aretenir}[L'essentiel de la section]
\begin{itemize}
\item Une \cle{réaction nucléaire provoquée} résulte de l'impact d'un projectile sur un noyau cible ; une \cle{transmutation} transforme un noyau en un noyau d'un autre élément.
\item Étapes fondatrices : Rutherford (1919, première transmutation), Chadwick (1932, découverte du neutron), Joliot-Curie (1934, radioactivité artificielle).
\item Toute équation nucléaire vérifie les \cle{lois de Soddy} : conservation du nombre de nucléons $A$ et de la charge $Z$.
\item Pour identifier une particule inconnue : bilan des $A$, bilan des $Z$, identification à l'aide du tableau périodique.
\end{itemize}
\end{aretenir}

\begin{exoresolu}[Le bombardement du bore]
En bombardant du bore 10 avec des neutrons, on obtient du lithium 7 et une particule $X$.
\begin{enumerate}
\item Écrire l'équation de la réaction en identifiant $X$.
\item Vérifier les lois de conservation.
\end{enumerate}
\tcblower
\textbf{1.} L'équation s'écrit a priori : $\ce{^{10}_{5}B + ^{1}_{0}n -> ^{7}_{3}Li + ^{A}_{Z}X}$.

Appliquons la méthode en 3 étapes.
\begin{align*}
\text{Conservation de } A : \quad & 10 + 1 = 7 + A \;\Rightarrow\; A = 4 \\
\text{Conservation de } Z : \quad & 5 + 0 = 3 + Z \;\Rightarrow\; Z = 2
\end{align*}
La particule $(Z = 2,\ A = 4)$ est le noyau d'hélium, c'est-à-dire une \textbf{particule $\alpha$} : $X = \ce{^{4}_{2}He}$. L'équation complète est :
\[
\ce{^{10}_{5}B + ^{1}_{0}n -> ^{7}_{3}Li + ^{4}_{2}He}
\]

\textbf{2.} Vérification : nucléons, $10 + 1 = 11$ et $7 + 4 = 11$ \checkmark ; charge, $5 + 0 = 5$ et $3 + 2 = 5$ \checkmark. L'équation est correctement équilibrée.

\emph{Remarque} : cette réaction est utilisée dans les détecteurs de neutrons et dans certaines radiothérapies (thérapie par capture de neutrons).
\end{exoresolu}

\begin{exercice}[À toi de jouer]
\begin{enumerate}
\item Compléter et équilibrer : $\ce{^{6}_{3}Li + ^{1}_{0}n -> ^{3}_{1}H + X}$. Identifier $X$.
\item Compléter : $\ce{^{2}_{1}H + ^{2}_{1}H -> ^{3}_{2}He + X}$. Identifier $X$.
\item Le polonium 210 émet une particule $\alpha$ et se transforme en plomb (Pb, $Z = 82$). Écrire l'équation de cette désintégration.
\end{enumerate}
\end{exercice}

\begin{corrige}[Exercice : À toi de jouer]
\textbf{1.} Conservation de $A$ : $6 + 1 = 3 + A \Rightarrow A = 4$. Conservation de $Z$ : $3 + 0 = 1 + Z \Rightarrow Z = 2$. Donc $X = \ce{^{4}_{2}He}$ (particule $\alpha$) : $\ce{^{6}_{3}Li + ^{1}_{0}n -> ^{3}_{1}H + ^{4}_{2}He}$.

\textbf{2.} Conservation de $A$ : $2 + 2 = 3 + A \Rightarrow A = 1$. Conservation de $Z$ : $1 + 1 = 2 + Z \Rightarrow Z = 0$. Donc $X = \ce{^{1}_{0}n}$ (neutron) : $\ce{^{2}_{1}H + ^{2}_{1}H -> ^{3}_{2}He + ^{1}_{0}n}$.

\textbf{3.} Le polonium a pour numéro atomique $Z = 84$. L'équation s'écrit $\ce{^{210}_{84}Po -> ^{A}_{82}Pb + ^{4}_{2}He}$ avec $210 = A + 4$, soit $A = 206$ : $\ce{^{210}_{84}Po -> ^{206}_{82}Pb + ^{4}_{2}He}$. Vérification : $84 = 82 + 2$ \checkmark.
\end{corrige}

Dans la suite de la leçon, nous verrons deux transmutations provoquées d'une importance capitale : la \cle{fission}, qui alimente les centrales nucléaires, et la \cle{fusion}, qui fait briller le Soleil --- et nous calculerons l'énergie gigantesque qu'elles libèrent grâce à la célèbre relation $E = \Delta m\, c^2$.

\coursec{Fission nucléaire et réaction en chaîne}

Dans la section précédente, nous avons vu qu'il est possible de \emph{provoquer} la transformation d'un noyau en le bombardant avec une particule : c'est la \cle{transmutation}. La question qui se pose alors est naturelle : peut-on trouver des réactions nucléaires provoquées qui, en plus de transformer la matière, \emph{libèrent une énergie considérable} et \emph{s'entretiennent d'elles-mêmes} ? La réponse est oui, et elle a changé le cours du XX\textsuperscript{e} siècle : c'est la \cle{fission nucléaire}.

\courssub{Qu'est-ce que la fission nucléaire ?}

\textbf{Intuition.} Un noyau lourd comme celui de l'uranium contient plus de 200 nucléons. Les forces nucléaires qui le lient ont une très courte portée : chaque nucléon n'attire que ses voisins immédiats. En revanche, la répulsion électrique entre les 92 protons agit à longue distance, sur tout le volume du noyau. Un noyau lourd est donc comme une grosse goutte d'eau tendue, prête à se scinder : il suffit d'une petite perturbation — l'absorption d'un neutron — pour qu'il se déforme puis se rompe en deux morceaux.

\begin{definition}[Fission nucléaire]
La \cle{fission nucléaire} est une réaction nucléaire \emph{provoquée} au cours de laquelle un noyau lourd ($A > 200$), dit \cle{fissile}, se scinde en deux noyaux plus légers, appelés \cle{fragments de fission}, sous l'impact d'un neutron. Elle s'accompagne de l'émission de 2 ou 3 neutrons et d'un dégagement d'énergie.
\end{definition}

\begin{attention}[Fission et radioactivité : ne pas confondre]
La fission est une réaction \emph{provoquée} : elle nécessite l'absorption d'un neutron. Ce n'est \textbf{pas} une désintégration spontanée comme les radioactivités $\alpha$, $\beta$ ou $\gamma$. Certains noyaux très lourds peuvent aussi fissionner spontanément, mais ce phénomène est rare et hors programme.
\end{attention}

\courssub{Exemple fondateur : la fission de l'uranium-235}

L'isotope \ce{^{235}_{92}U} est le combustible fissile le plus utilisé. Lorsqu'il capture un neutron (de préférence lent, dit \emph{thermique}), il forme un noyau intermédiaire \ce{^{236}_{92}U} très instable, qui se scinde presque immédiatement. Une des fissions possibles s'écrit :
\[
\ce{^{235}_{92}U + ^{1}_{0}n -> ^{94}_{38}Sr + ^{139}_{54}Xe + 3^{1}_{0}n}
\]

Vérifions les lois de conservation de Soddy :
\begin{itemize}
\item \textbf{Conservation du nombre de nucléons :} $235 + 1 = 236$ à gauche ; $94 + 139 + 3 = 236$ à droite : la loi est vérifiée.
\item \textbf{Conservation de la charge :} $92 + 0 = 92$ à gauche ; $38 + 54 + 0 = 92$ à droite : la loi est vérifiée.
\end{itemize}

\begin{methode}[Équilibrer une équation de fission et compter les neutrons]
\begin{enumerate}
\item Écrire l'équation avec le noyau inconnu $\ce{^{A}_{Z}X}$.
\item Appliquer la conservation du nombre de nucléons : somme des $A$ à gauche = somme des $A$ à droite (ne pas oublier de multiplier le nombre de masse du neutron par le nombre de neutrons émis).
\item Appliquer la conservation du nombre de charge : somme des $Z$ à gauche = somme des $Z$ à droite (le neutron a $Z = 0$).
\item Identifier l'élément $X$ grâce à son numéro atomique $Z$ dans la classification périodique.
\item Conclure en écrivant l'équation complète et en vérifiant une dernière fois.
\end{enumerate}
\end{methode}

\begin{exemplebox}[Identifier un fragment de fission]
Équilibrons : $\ce{^{235}_{92}U + ^{1}_{0}n -> ^{141}_{56}Ba + ^{A}_{Z}X + 3^{1}_{0}n}$.

Conservation de $A$ : $235 + 1 = 141 + A + 3 \Rightarrow A = 236 - 144 = 92$.

Conservation de $Z$ : $92 = 56 + Z \Rightarrow Z = 36$ : c'est le krypton.

L'équation complète est donc :
\[
\ce{^{235}_{92}U + ^{1}_{0}n -> ^{141}_{56}Ba + ^{92}_{36}Kr + 3^{1}_{0}n}
\]
\end{exemplebox}

\courssub{Un caractère statistique : les produits de fission ne sont pas uniques}

\begin{attention}[Erreur fréquente]
Beaucoup d'élèves croient que la fission de l'uranium-235 produit \emph{toujours} les mêmes fragments (strontium et xénon, par exemple). C'est faux ! Il existe des \textbf{dizaines de fissions possibles} : seules les \emph{sommes} des nombres de masse et des numéros atomiques sont imposées par les lois de conservation. La nature des fragments est \emph{aléatoire} : on parle de \cle{distribution statistique} des fragments de fission.
\end{attention}

Si l'on mesure expérimentalement la proportion de fragments produits en fonction de leur nombre de masse $A$, on obtient une courbe remarquable à \emph{deux bosses} : les fragments les plus probables ont des masses voisines de $A \approx 95$ et $A \approx 140$, et la fission en deux fragments de masses égales ($A \approx 118$) est au contraire très peu probable.

\begin{popfigure}
\begin{center}
\begin{tikzpicture}
\begin{axis}[
width=0.9\linewidth, height=6cm,
xlabel={Nombre de masse $A$ du fragment},
ylabel={Rendement de fission (\%)},
xmin=60, xmax=165, ymin=0, ymax=8,
xtick={70,80,90,95,100,110,118,130,140,150},
ytick={0,2,4,6},
grid=both, grid style={popline!40},
legend style={at={(0.5,1.02)}, anchor=south, draw=none}
]
\addplot[popcurve, domain=70:160, samples=200]
{6*exp(-((x-95)^2)/60) + 5.5*exp(-((x-140)^2)/70) + 0.05};
\addlegendentry{Distribution expérimentale (allure)}
\draw[dashed, poporange, thick] (axis cs:95,0) -- (axis cs:95,6.2) node[above, font=\small]{$A \approx 95$};
\draw[dashed, popgreen, thick] (axis cs:140,0) -- (axis cs:140,5.8) node[above, font=\small]{$A \approx 140$};
\draw[dashed, popmuted] (axis cs:118,0) -- (axis cs:118,1.5) node[above, font=\small, text=popmuted]{$A \approx 118$};
\end{axis}
\end{tikzpicture}
\end{center}
\legende{Allure de la distribution en masse des fragments de fission de \ce{^{235}_{92}U} : les fragments de masses voisines de 95 et 140 sont les plus probables ; la fission symétrique ($A \approx 118$) est rare.}
\end{popfigure}

\courssub{La réaction en chaîne}

Regarde à nouveau l'équation de fission : \textbf{un} neutron suffit à déclencher la fission, mais \textbf{trois} neutrons sont émis. Chacun de ces neutrons peut à son tour provoquer la fission d'un autre noyau d'uranium-235, qui émettra à nouveau des neutrons, et ainsi de suite. Le nombre de fissions croît alors de façon exponentielle : $1$, puis $3$, puis $9$, puis $27$\ldots

\begin{definition}[Réaction en chaîne]
Une \cle{réaction en chaîne} est une suite de fissions dans laquelle les neutrons émis par chaque fission provoquent de nouvelles fissions. Si, en moyenne, plus d'un neutron par fission déclenche une nouvelle fission, le nombre de réactions croît exponentiellement.
\end{definition}

\begin{popfigure}
\begin{center}
\begin{tikzpicture}[
scale=0.95, transform shape,
noyau/.style={circle, draw=popdark, fill=poporangeL, minimum size=8mm, inner sep=1pt, font=\small},
neutron/.style={circle, draw=popblue, fill=popblueL, minimum size=3.5mm, inner sep=0pt, font=\scriptsize},
fleche/.style={-{Stealth[length=2mm]}, thick, popblue}
]
% Étape 1
\node[neutron] (n0) at (0,0) {n};
\node[noyau] (u1) at (1.8,0) {\ce{^{235}U}};
\draw[fleche] (n0) -- (u1);
% Fission 1 -> 3 neutrons
\node[neutron] (n1) at (4,-2) {n};
\node[neutron] (n2) at (4,0) {n};
\node[neutron] (n3) at (4,2) {n};
\draw[fleche, poporange] (u1) -- (n1);
\draw[fleche, poporange] (u1) -- (n2);
\draw[fleche, poporange] (u1) -- (n3);
% Étape 2 : 3 noyaux
\node[noyau] (u2a) at (6.2,-2) {\ce{^{235}U}};
\node[noyau] (u2b) at (6.2,0) {\ce{^{235}U}};
\node[noyau] (u2c) at (6.2,2) {\ce{^{235}U}};
\draw[fleche] (n1) -- (u2a);
\draw[fleche] (n2) -- (u2b);
\draw[fleche] (n3) -- (u2c);
% Étape 3 : neutrons multiples
\foreach \y/\u in {-2/u2a, 0/u2b, 2/u2c} {
  \draw[fleche, poporange] (\u) -- ++(1.5,0.5);
  \draw[fleche, poporange] (\u) -- ++(1.5,0);
  \draw[fleche, poporange] (\u) -- ++(1.5,-0.5);
}
\node[font=\small, text=popmuted] at (8.6,-2) {$3$ n};
\node[font=\small, text=popmuted] at (8.6,0) {$3$ n};
\node[font=\small, text=popmuted] at (8.6,2) {$3$ n};
% Étiquettes d'étapes
\node[font=\small\bfseries, text=popdark] at (0.9,1.1) {1 neutron};
\node[font=\small\bfseries, text=popdark] at (4,3.1) {1 fission $\to$ 3 neutrons};
\node[font=\small\bfseries, text=popdark] at (6.2,3.1) {3 fissions};
\node[font=\small\bfseries, text=popdark] at (8.6,3.1) {9 neutrons \ldots};
\end{tikzpicture}
\end{center}
\legende{Schéma en arbre d'une réaction en chaîne : chaque fission libère 3 neutrons susceptibles de provoquer 3 nouvelles fissions. Le nombre de réactions croît géométriquement.}
\end{popfigure}

\courssub{Masse critique : réaction contrôlée ou divergente}

Tous les neutrons émis ne provoquent pas une nouvelle fission : certains s'échappent de l'échantillon, d'autres sont absorbés sans fission par des impuretés. Tout dépend donc du \emph{bilan} des neutrons. On définit le facteur de multiplication $k$ : nombre moyen de fissions provoquées par les neutrons issus d'une fission.

\begin{itemize}
\item Si $k < 1$ : la réaction \emph{s'éteint} (régime sous-critique) ;
\item si $k = 1$ : la réaction se \emph{maintient} à puissance constante (régime critique, cas d'un réacteur en fonctionnement stable) ;
\item si $k > 1$ : la réaction \emph{diverge} (régime sur-critique, cas de la bombe A).
\end{itemize}

\begin{definition}[Masse critique]
La \cle{masse critique} d'un combustible fissile est la masse \emph{minimale} de ce combustible, dans des conditions données (forme géométrique, pureté, présence d'un modérateur), nécessaire pour entretenir une réaction en chaîne ($k \geq 1$). En dessous, trop de neutrons s'échappent par la surface par rapport au nombre de fissions créées dans le volume.
\end{definition}

\textbf{Pourquoi la géométrie compte-t-elle ?} Les neutrons perdus s'échappent par la \emph{surface} de l'échantillon (proportionnelle à $R^2$ pour une sphère), tandis que les fissions se produisent dans tout le \emph{volume} (proportionnel à $R^3$). Le rapport pertes/production varie donc comme $1/R$ : plus l'échantillon est gros, moins les pertes relatives sont importantes. Pour l'uranium-235 pur en forme de sphère, la masse critique est de l'ordre de \SI{50}{kg} ; elle tombe à quelques kilogrammes si la sphère est entourée d'un réflecteur de neutrons.

\begin{attention}[Deux applications, deux philosophies]
\begin{itemize}
\item \textbf{Bombe A} : on réunit brutalement deux masses sous-critiques pour former une masse très sur-critique ($k \gg 1$) : la réaction diverge en une fraction de seconde et libère une énergie gigantesque (Hiroshima, 1945).
\item \textbf{Réacteur nucléaire} : on maintient soigneusement $k = 1$ grâce à des barres de contrôle ; la puissance est stable et l'énergie est récupérée progressivement pour produire de l'électricité.
\end{itemize}
\end{attention}

\courssub{Schéma de principe d'une centrale nucléaire}

Une centrale nucléaire est, en définitive, une « machine à vapeur » dont la chaudière est un réacteur à fission. Ses éléments essentiels :

\begin{itemize}
\item le \cle{combustible} : barres d'uranium enrichi en \ce{^{235}U} (environ 3 à 5 \%), placées dans le \emph{cœur} du réacteur, où se déroulent les fissions ;
\item le \cle{modérateur} (eau ordinaire dans les réacteurs à eau pressurisée) : il \emph{ralentit} les neutrons rapides émis par les fissions, car les neutrons lents sont beaucoup plus efficaces pour provoquer de nouvelles fissions de \ce{^{235}U} ;
\item les \cle{barres de contrôle} (en cadmium ou en bore) : elles \emph{absorbent} les neutrons ; en les enfonçant ou les retirant, on ajuste $k$ et donc la puissance du réacteur ;
\item le \cle{circuit de refroidissement} (circuit primaire) : l'eau sous pression évacue la chaleur du cœur vers un générateur de vapeur ;
\item le circuit secondaire : la vapeur produite fait tourner une \emph{turbine} couplée à un \emph{alternateur} qui produit l'électricité ; un condenseur (refroidi par un fleuve ou l'air des tours aéroréfrigérantes) referme le cycle.
\end{itemize}

\begin{popfigure}
\begin{center}
\begin{tikzpicture}[scale=0.9, transform shape,
boite/.style={draw=popdark, thick, rounded corners=2pt, fill=popcream, minimum height=9mm, align=center, font=\small},
fleche/.style={-{Stealth[length=2.2mm]}, very thick}]
% Enceinte du réacteur
\draw[draw=popdark, very thick, fill=poporangeL!60, rounded corners=4pt] (0,-1.6) rectangle (4.4,2.2);
\node[font=\small\bfseries, text=popdark] at (2.2,1.85) {C\oe ur du réacteur};
% Combustible
\foreach \x in {0.7,1.3,1.9} {
  \draw[fill=poporange, draw=popdark] (\x,-1.3) rectangle (\x+0.25,1.2);
}
\node[font=\scriptsize, text=popdark] at (1.3,-1.95) {combustible \ce{^{235}U}};
% Barres de contrôle
\foreach \x in {2.6,3.2} {
  \draw[fill=popblue, draw=popdark] (\x,0.2) rectangle (\x+0.25,2.0);
}
\draw[-{Stealth}, thick, popblue] (3.1,2.9) -- (3.1,2.25);
\node[font=\scriptsize, text=popblue, align=center] at (3.1,3.35) {barres de\\contrôle};
% Modérateur
\node[font=\scriptsize, text=popmuted] at (3.6,-0.8) {eau (modérateur)};
% Circuit primaire
\draw[fleche, poppink] (4.4,0.9) -- (5.8,0.9);
\node[boite, fill=poppinkL, align=center] (ech) at (7.1,0.9) {Générateur\\de vapeur};
\draw[fleche, poppink] (7.1,-0.1) -- (7.1,-0.7) -- (4.4,-0.7);
\node[font=\scriptsize, text=poppink] at (5.6,-1.05) {circuit primaire (eau chaude)};
% Circuit secondaire
\draw[fleche, popgreen] (8.4,0.9) -- (9.6,0.9);
\node[boite, fill=popgreenL] (turb) at (10.8,0.9) {Turbine};
\draw[fleche, popgreen] (10.8,-0.1) -- (10.8,-0.7) -- (8.4,-0.7);
\node[font=\scriptsize, text=popgreen!60!black] at (9.6,-1.05) {vapeur (circuit secondaire)};
% Alternateur
\node[boite, fill=popgoldL] (alt) at (10.8,2.3) {Alternateur};
\draw[-{Stealth}, thick, popdark] (10.8,1.5) -- (10.8,1.85);
\draw[fleche, popgold] (11.8,2.3) -- (13.0,2.3) node[right, font=\small\bfseries, text=popdark] {électricité};
\end{tikzpicture}
\end{center}
\legende{Schéma simplifié d'une centrale nucléaire : la fission chauffe l'eau du circuit primaire ; la chaleur vaporise l'eau du circuit secondaire qui entraîne la turbine et l'alternateur.}
\end{popfigure}

\courssub{Avantages et risques du nucléaire}

\begin{center}
\begin{tabularx}{\linewidth}{|X|X|}
\hline
\rowcolor{popblueL} \textbf{Avantages} & \textbf{Risques et inconvénients} \\
\hline
Production massive et continue d'électricité, sans émission de \ce{CO2} en fonctionnement (atout contre le réchauffement climatique). & Déchets radioactifs à vie longue (certains fragments de fission restent dangereux pendant des milliers d'années) : problème du stockage. \\
\hline
Très grande densité énergétique : peu de combustible pour beaucoup d'énergie. & Risque d'accident majeur en cas de perte de contrôle de la réaction en chaîne ou de refroidissement : Tchernobyl (1986), Fukushima (2011). \\
\hline
Indépendance énergétique des pays producteurs. & Coût élevé de construction et de démantèlement ; risque de prolifération militaire. \\
\hline
\end{tabularx}
\end{center}

\begin{savaistu}[Un gramme d'uranium contre deux tonnes de pétrole]
La fission complète d'\textbf{un gramme} d'uranium-235 libère environ \SI{8e10}{J}, soit autant d'énergie que la combustion d'environ \textbf{2 tonnes de pétrole} ! C'est cette densité énergétique extraordinaire qui explique l'intérêt du nucléaire. Au Cameroun, où la production électrique repose essentiellement sur l'hydroélectricité (barrages d'Edéa, Song Loulou, Lom Pangar, Nachtigal), certains experts évoquent à long terme le nucléaire comme complément possible pour l'industrialisation, sous réserve de maîtriser la sûreté et la gestion des déchets.
\end{savaistu}

\begin{exoresolu}[Compter les neutrons et équilibrer une fission]
\textbf{Énoncé.} Lors de la fission d'un noyau d'uranium-235 provoquée par un neutron, on observe parmi les produits le xénon \ce{^{140}_{54}Xe}, le strontium \ce{^{A}_{38}Sr} et $n = 2$ neutrons. Aucun autre produit n'est formé. (1) Déterminer $A$ et écrire l'équation complète de la réaction. (2) Expliquer pourquoi cette réaction peut s'entretenir d'elle-même.
\tcblower
\textbf{Solution.} (1) L'équation s'écrit :
\[
\ce{^{235}_{92}U + ^{1}_{0}n -> ^{140}_{54}Xe + ^{A}_{38}Sr + 2^{1}_{0}n}
\]
Conservation de la charge : $92 + 0 = 54 + 38 + 0$, soit $92 = 92$ : la loi est vérifiée, ce qui confirme que $Z_{\text{Sr}} = 38$.

Conservation du nombre de nucléons :
\[
235 + 1 = 140 + A + 2 \quad\Rightarrow\quad A = 236 - 142 = 94.
\]
L'équation complète est donc :
\[
\ce{^{235}_{92}U + ^{1}_{0}n -> ^{140}_{54}Xe + ^{94}_{38}Sr + 2^{1}_{0}n}
\]
Vérification finale : $140 + 94 + 2 = 236$ et $54 + 38 = 92$ : les deux lois de Soddy sont satisfaites.

(2) Chaque fission est déclenchée par \emph{un} neutron mais en libère \emph{deux} : si la masse d'uranium dépasse la masse critique, ces neutrons provoquent à leur tour de nouvelles fissions : la réaction s'entretient et s'amplifie, c'est la \cle{réaction en chaîne}.
\end{exoresolu}

\begin{aretenir}[L'essentiel de la section]
\begin{itemize}
\item La \cle{fission} est la rupture provoquée d'un noyau lourd ($A > 200$) en deux fragments plus légers, sous l'impact d'un neutron, avec émission de 2 ou 3 neutrons et dégagement d'énergie.
\item Les produits de fission sont \emph{variés} (distribution statistique à deux bosses, $A \approx 95$ et $A \approx 140$) ; seules les sommes des $A$ et des $Z$ sont conservées.
\item Les neutrons émis entretiennent une \cle{réaction en chaîne}, contrôlée ($k = 1$) dans un réacteur, divergente ($k > 1$) dans une bombe A.
\item La \cle{masse critique} est la masse minimale de combustible fissile permettant d'entretenir la réaction en chaîne.
\item Dans une centrale : le combustible fissionne, le modérateur ralentit les neutrons, les barres de contrôle les absorbent, et la chaleur dégagée produit de la vapeur qui entraîne un turbo-alternateur.
\end{itemize}
\end{aretenir}

\coursec{Fusion nucléaire}

Dans la section précédente, nous avons vu comment un noyau \emph{lourd} comme l'uranium 235 peut se \emph{scinder} en deux fragments plus légers en libérant une énergie considérable : c'est la fission, exploitée dans les centrales nucléaires. Mais la nature offre une voie symétrique et encore plus prodigieuse : au lieu de casser les gros noyaux, on peut \emph{assembler} les tout petits. C'est ce processus, la \cle{fusion nucléaire}, qui fait briller le Soleil au-dessus de nos têtes, chaque jour, depuis 4,6 milliards d'années. Comprendre la fusion, c'est comprendre la source ultime de presque toute l'énergie disponible sur Terre — et entrevoir peut-être l'énergie du futur pour nos villes, de Douala à Maroua.

\courssub{Définition de la fusion nucléaire}

\textbf{Observation et intuition.} Reprenons la courbe d'Aston rencontrée lors de l'étude de l'énergie de liaison : l'énergie de liaison \emph{par nucléon} $E_\ell/A$ est maximale autour du fer ($A \approx 56$), où elle vaut environ $\SI{8,8}{MeV}$ par nucléon. Les noyaux très légers (hydrogène, hélium) sont \emph{mal liés} : leur énergie de liaison par nucléon est faible (par exemple $\approx \SI{1,1}{MeV}$ par nucléon pour le deutérium). Si l'on parvient à réunir deux noyaux légers pour former un noyau plus lourd et mieux lié, la masse du produit est \emph{inférieure} à la somme des masses des réactifs : le défaut de masse $\Delta m$ se transforme en énergie libérée, conformément à la relation d'Einstein $E = \Delta m\, c^2$. C'est exactement le raisonnement inverse de celui de la fission, où l'on partait d'un noyau lourd mal lié pour aboutir à des fragments mieux liés. Dans les deux cas, on « remonte » vers le sommet de la courbe d'Aston, et c'est cette remontée qui libère l'énergie.

\begin{definition}[Fusion nucléaire]
La \cle{fusion nucléaire} est une réaction nucléaire provoquée au cours de laquelle deux noyaux \emph{légers} s'unissent pour former un noyau \emph{plus lourd}, avec libération d'énergie. Comme toute réaction nucléaire, elle obéit aux \textbf{lois de conservation de Soddy} :
\begin{itemize}
\item conservation du nombre de nucléons : $\sum A_{\text{réactifs}} = \sum A_{\text{produits}}$ ;
\item conservation de la charge électrique : $\sum Z_{\text{réactifs}} = \sum Z_{\text{produits}}$.
\end{itemize}
\end{definition}

\begin{exemplebox}[La réaction deutérium -- tritium]
La réaction de fusion la plus étudiée, car la plus « facile » à réaliser, met en jeu deux isotopes de l'hydrogène : le \cle{deutérium} \ce{^{2}_{1}H} (un proton + un neutron) et le \cle{tritium} \ce{^{3}_{1}H} (un proton + deux neutrons) :
\[ \ce{^{2}_{1}H + ^{3}_{1}H -> ^{4}_{2}He + ^{1}_{0}n} \]
Vérifions les lois de conservation :
\begin{itemize}
\item nucléons : $2 + 3 = 5$ à gauche, $4 + 1 = 5$ à droite : la loi est vérifiée ;
\item charges : $1 + 1 = 2$ à gauche, $2 + 0 = 2$ à droite : la loi est vérifiée.
\end{itemize}
Cette réaction libère environ $\SI{17,6}{MeV}$ par fusion, emportés sous forme d'énergie cinétique par le noyau d'hélium (environ $\SI{3,5}{MeV}$) et par le neutron (environ $\SI{14,1}{MeV}$).
\end{exemplebox}

\begin{popfigure}
\begin{center}
\begin{tikzpicture}[scale=1.0]
% Deutérium
\fill[popblue] (0,0) circle (0.28);
\fill[popblue!60] (0.32,0.18) circle (0.28);
\node[below=0.35cm of {(0.15,-0.1)}, popdark] {\ce{^{2}_{1}H}};
% Tritium
\fill[poporange] (2.2,0.1) circle (0.28);
\fill[poporange!70] (2.55,-0.15) circle (0.28);
\fill[poporange!50] (2.5,0.35) circle (0.28);
\node[below=0.35cm of {(2.4,-0.2)}, popdark] {\ce{^{3}_{1}H}};
% Flèche
\draw[-{Stealth[length=4mm]},very thick,popdark] (3.3,0.1) -- (4.6,0.1);
% Hélium
\fill[popgreen] (5.6,0.05) circle (0.30);
\fill[popgreen!70] (5.95,0.28) circle (0.30);
\fill[popgreen!55] (6.0,-0.22) circle (0.30);
\fill[popgreen!85] (6.32,0.05) circle (0.30);
\node[below=0.4cm of {(5.95,-0.25)}, popdark] {\ce{^{4}_{2}He}};
% Neutron
\fill[popmuted] (7.6,0.1) circle (0.22);
\node[below=0.3cm of {(7.6,0.05)}, popdark] {\ce{^{1}_{0}n}};
% Energie
\node[popgold,font=\bfseries] at (6.6,1.0) {$+\ E = \Delta m\, c^2 \approx \SI{17,6}{MeV}$};
\end{tikzpicture}
\end{center}
\legende{Fusion deutérium -- tritium : deux noyaux légers s'unissent pour former un noyau d'hélium 4 plus un neutron rapide, avec libération d'énergie.}
\end{popfigure}

\courssub{Conditions de la fusion : vaincre la répulsion coulombienne}

\textbf{Le problème fondamental.} Les deux noyaux que l'on veut faire fusionner sont tous deux chargés \emph{positivement} (ce sont des assemblages de protons et de neutrons, $Z \geq 1$). Or deux charges de même signe se \emph{repoussent} : c'est la force électrostatique (ou coulombienne), dont la norme est donnée par la loi de Coulomb :
\[ F = k\,\frac{Z_1 Z_2\, e^2}{r^2} \qquad \text{avec} \quad k = \SI{9,0e9}{N.m^2.C^{-2}}. \]
Tant que les noyaux sont distants de plus de quelques femtomètres, cette répulsion domine et les empêche de s'approcher. Mais si l'on parvient à les rapprocher à une distance de l'ordre de $r \approx \SI{e-15}{m}$, l'\cle{interaction forte} — attractive et de très courte portée — prend le dessus et « soude » les deux noyaux en un seul. Il existe donc une \emph{barrière de potentiel} électrostatique à franchir.

\textbf{Ordre de grandeur de la barrière.} Pour deux noyaux de deutérium ($Z_1 = Z_2 = 1$) au contact ($r \approx \SI{2e-15}{m}$), l'énergie potentielle électrostatique vaut :
\[ E_p = k\,\frac{e^2}{r} \approx 9{,}0\times 10^{9} \times \frac{(1{,}6\times 10^{-19})^2}{2\times 10^{-15}} \approx \SI{1,2e-13}{J} \approx \SI{0,7}{MeV}. \]
Vérifions l'homogénéité : $k e^2 / r$ s'exprime en $\si{N.m^2.C^{-2}} \times \si{C^2} / \si{m} = \si{N.m} = \si{J}$ : le résultat est bien une énergie. Pour franchir cette barrière, les noyaux doivent posséder une énergie cinétique d'agitation thermique considérable, ce qui impose des températures gigantesques.

\begin{propriete}[Condition de température pour la fusion — admise]
La fusion nucléaire ne peut s'amorcer qu'à des \textbf{températures de l'ordre de $10^8$ K} (cent millions de degrés). À de telles températures :
\begin{itemize}
\item l'agitation thermique donne aux noyaux une énergie cinétique suffisante pour vaincre (ou franchir par effet tunnel) la répulsion coulombienne ;
\item la matière n'est plus un gaz ordinaire : les atomes sont totalement ionisés. On obtient un \cle{plasma}, quatrième état de la matière, constitué d'un mélange de noyaux et d'électrons libres.
\end{itemize}
\emph{Justification qualitative} : l'énergie cinétique moyenne d'agitation est de l'ordre de $E_c \sim k_B T$ avec $k_B = \SI{1,38e-23}{J.K^{-1}}$. Pour approcher l'énergie de la barrière coulombienne ($\sim \SI{e-13}{J}$), il faudrait $T \sim \dfrac{E_p}{k_B} \sim \dfrac{10^{-13}}{1{,}4\times 10^{-23}} \sim 10^{10}$ K ; grâce à la distribution statistique des vitesses (les noyaux les plus rapides de la distribution) et à l'effet tunnel quantique, la fusion devient effective dès $T \approx 10^8$ K. (Cette justification est qualitative et non exigible au Baccalauréat.)
\end{propriete}

\begin{attention}[Ne pas confondre fission et fusion]
Dans une équation nucléaire, demande-toi toujours : \textbf{le noyau de départ est-il lourd ou léger ?}
\begin{itemize}
\item \textbf{Fission} : un noyau \emph{lourd} ($A \gtrsim 230$, ex. \ce{^{235}_{92}U}) frappé par un neutron se \emph{divise} en deux noyaux moyens + des neutrons.
\item \textbf{Fusion} : deux noyaux \emph{légers} ($A \leq 3$ en général, ex. \ce{^{2}_{1}H}, \ce{^{3}_{1}H}) s'\emph{unissent} en un noyau plus lourd.
\end{itemize}
Erreur fréquente au Bac : écrire $\ce{^{235}_{92}U + ^{1}_{0}n -> ...}$ et parler de « fusion » parce qu'un neutron « fusionne » avec l'uranium. Non ! Le neutron est \emph{capturé}, puis le noyau composé se \emph{scinde} : c'est une fission. De même, la fusion exige des températures de $10^8$ K, alors que la fission se déclenche à température ambiante par simple capture d'un neutron lent.
\end{attention}

\begin{exoresolu}[Équilibrer une équation de fusion]
\textbf{Énoncé.} Deux noyaux de deutérium peuvent fusionner selon la réaction :
\[ \ce{^{2}_{1}H + ^{2}_{1}H -> ^{3}_{2}He + X} \]
Identifier la particule $X$ en précisant son symbole complet.
\tcblower
\textbf{Solution détaillée.} Écrivons $X = \ce{^{A}_{Z}X}$ et appliquons les lois de conservation de Soddy.
\begin{itemize}
\item Conservation du nombre de nucléons : $2 + 2 = 3 + A$, d'où $A = 1$.
\item Conservation de la charge : $1 + 1 = 2 + Z$, d'où $Z = 0$.
\end{itemize}
La particule $X$ a un nucléon et aucune charge : c'est un \textbf{neutron} \ce{^{1}_{0}n}. L'équation complète s'écrit :
\[ \ce{^{2}_{1}H + ^{2}_{1}H -> ^{3}_{2}He + ^{1}_{0}n} \]
Cette réaction libère environ $\SI{3,3}{MeV}$. (Remarque : une autre voie possible est $\ce{^{2}_{1}H + ^{2}_{1}H -> ^{3}_{1}H + ^{1}_{1}H}$, qui libère environ $\SI{4,0}{MeV}$.)
\end{exoresolu}

\courssub{La fusion dans les étoiles : le Soleil, un réacteur naturel}

\textbf{Le cycle proton-proton.} Au cœur du Soleil, la température atteint environ $1,5\times 10^7$ K et la pression est colossale : les conditions (aidées par l'effet tunnel) permettent la fusion des protons. Pour les étoiles de la taille du Soleil, le mécanisme dominant est la \cle{chaîne proton-proton}, qui se déroule en plusieurs étapes :
\begin{enumerate}
\item deux protons fusionnent en deutérium : $\ce{^{1}_{1}H + ^{1}_{1}H -> ^{2}_{1}H + ^{0}_{1}e + \nu}$ (émission d'un positon et d'un neutrino) ;
\item le deutérium capture un proton : $\ce{^{2}_{1}H + ^{1}_{1}H -> ^{3}_{2}He + \gamma}$ ;
\item deux noyaux d'hélium 3 fusionnent : $\ce{^{3}_{2}He + ^{3}_{2}He -> ^{4}_{2}He + 2\,^{1}_{1}H}$.
\end{enumerate}

\begin{aretenir}[Bilan global de la fusion solaire]
En additionnant toutes les étapes, les particules intermédiaires s'éliminent et le bilan global de la chaîne proton-proton s'écrit :
\[ \ce{4\,^{1}_{1}H -> ^{4}_{2}He + 2\,^{0}_{1}e + 2\,\nu} \quad + \text{énergie} \]
Quatre protons se transforment en un noyau d'hélium 4, deux positons et deux neutrinos. L'énergie libérée est d'environ $\SI{25}{MeV}$ par noyau d'hélium formé (en tenant compte de l'annihilation des positons avec les électrons). Chaque seconde, le Soleil convertit ainsi environ \textbf{600 millions de tonnes} d'hydrogène en hélium, et environ \textbf{4 millions de tonnes de matière} sont transformées en énergie rayonnée — dont une infime fraction suffit à éclairer et chauffer toute la Terre, à faire pousser le manioc et le maïs de nos champs et à alimenter nos panneaux solaires.
\end{aretenir}

\begin{exemplebox}[Énergie massique de la fusion solaire]
Calculons l'énergie libérée par kilogramme d'hydrogène fusionné. Former un noyau \ce{^{4}_{2}He} consomme 4 protons de masse totale $4 \times 1{,}6726\times 10^{-27} \approx \SI{6,69e-27}{kg}$ et libère $E \approx \SI{25}{MeV} = 25\times 10^6 \times 1{,}6\times 10^{-19} = \SI{4,0e-12}{J}$. L'énergie massique vaut donc :
\[ \frac{E}{m} = \frac{4{,}0\times 10^{-12}}{6{,}69\times 10^{-27}} \approx \SI{6e14}{J.kg^{-1}}. \]
C'est environ \textbf{dix millions de fois} l'énergie dégagée par la combustion d'un kilogramme d'essence ($\approx \SI{4,5e7}{J.kg^{-1}}$) ! La fusion est, à masse égale, la source d'énergie la plus puissante connue.
\end{exemplebox}

\courssub{La bombe H : la fusion non contrôlée}

L'Humanité a réalisé la fusion nucléaire avant de savoir la contrôler : la \cle{bombe H} (bombe à hydrogène, ou bombe thermonucléaire), testée pour la première fois en 1952, exploite la fusion non contrôlée des isotopes de l'hydrogène.

\begin{propriete}[Amorçage de la bombe H]
Les températures de $10^8$ K nécessaires à la fusion ne peuvent pas être atteintes par des moyens chimiques classiques. Dans une bombe H, l'amorçage est réalisé par une \textbf{bombe A} (bombe à fission) : l'explosion d'une charge de plutonium ou d'uranium porte brutalement la matière à des températures et des pressions suffisantes pour déclencher la fusion du deutérium et du tritium. La bombe H est donc un dispositif \emph{fission -- fusion}, d'une puissance destructrice pouvant dépasser \textbf{mille fois} celle de la bombe d'Hiroshima.
\end{propriete}

\courssub{La fusion contrôlée : l'enjeu énergétique du XXI\textsuperscript{e} siècle}

\textbf{Le défi.} Produire de l'électricité par fusion exige de maintenir un plasma à $10^8$ K... sans qu'aucune paroi matérielle ne puisse le toucher : à cette température, tout solide fondrait instantanément, et le contact avec les parois refroidirait le plasma. La solution retenue est le \cle{confinement magnétique} : le plasma étant constitué de particules chargées, un champ magnétique intense et convenablement dessiné peut courber leurs trajectoires et les maintenir en lévitation au centre d'une chambre à vide. Le dispositif le plus abouti est le \cle{tokamak} (acronyme russe de « chambre torique magnétique ») : une chambre en forme de tore (anneau) entourée de puissantes bobines supraconductrices créant le champ magnétique de confinement.

\begin{popfigure}
\begin{center}
\begin{tikzpicture}[scale=1.0]
% Chambre torique (vue en coupe stylisée)
\draw[very thick,popblue,fill=popblueL] (0,0) ellipse (3.2 and 1.5);
\draw[very thick,popblue,fill=white] (0,0) ellipse (1.5 and 0.7);
\node[popdark] at (0,-0.05) {\small vide};
% Plasma
\draw[very thick,poporange,decorate,decoration={coil,segment length=6pt,amplitude=3pt}] (-2.35,0) arc (180:0:2.35 and 1.1);
\draw[very thick,poporange,decorate,decoration={coil,segment length=6pt,amplitude=3pt}] (2.35,0) arc (0:-180:2.35 and 1.1);
\node[poporange,font=\bfseries] at (0,1.15) {\small plasma à $10^8$ K};
% Bobines
\foreach \x in {-2.6,-1.3,0,1.3,2.6}{
  \draw[thick,popgreen,fill=popgreenL] (\x-0.18,-1.9) rectangle (\x+0.18,1.9);
}
\node[popgreen!60!black] at (0,-2.35) {\small bobines de confinement magnétique};
% Champ B
\draw[-{Stealth[length=3mm]},thick,poppurple] (-3.9,0.4) -- (-3.9,-0.4) node[midway,left,font=\small] {$\vec{B}$};
\end{tikzpicture}
\end{center}
\legende{Schéma simplifié d'un tokamak : le plasma (orange) circule dans une chambre torique à vide (bleu), confiné par le champ magnétique $\vec{B}$ créé par des bobines supraconductrices (vert).}
\end{popfigure}

\begin{savaistu}[ITER : le Soleil en laboratoire]
Le projet international \cle{ITER} (35 nations coopérant, construction à Cadarache, en France) est le plus grand tokamak jamais conçu. Son objectif : démontrer la faisabilité scientifique et industrielle de la fusion en produisant \textbf{500 MW de puissance de fusion pour seulement 50 MW injectés} pour chauffer le plasma — un gain d'un facteur 10, jamais atteint jusqu'ici. Le combustible ? Le deutérium, extractible de l'eau de mer (environ 33 g par mètre cube), et le tritium, régénéré à partir du lithium. Quelques grammes de combustible de fusion libèrent autant d'énergie que plusieurs tonnes de pétrole ! Si la fusion contrôlée devient industrielle dans la seconde moitié du siècle, tous les pays en bénéficieront, y compris le Cameroun : une énergie quasi inépuisable, sans émission de gaz à effet de serre, qui pourrait compléter nos grands barrages hydroélectriques (Lom Pangar, Nachtigal, Memve'ele) et accélérer l'électrification rurale.
\end{savaistu}

\courssub{Comparaison fission -- fusion}

\begin{center}
\begin{tabularx}{\linewidth}{|l|X|X|}
\hline
\rowcolor{popblueL} \textbf{Critère} & \textbf{Fission} & \textbf{Fusion} \\
\hline
Noyaux mis en jeu & lourds (\ce{^{235}U}, \ce{^{239}Pu}) qui se scindent & légers (\ce{^{2}H}, \ce{^{3}H}) qui s'unissent \\
\hline
Condition de déclenchement & capture d'un neutron lent, à température ambiante & $T \approx 10^8$ K (plasma) \\
\hline
Combustible & uranium, minerai rare et concentré & deutérium extrait de l'eau de mer, quasi inépuisable \\
\hline
Déchets & produits de fission radioactifs à vie longue & hélium (inerte) ; structures activées à vie plus courte \\
\hline
Risque d'emballement & réaction en chaîne à contrôler & aucun : la réaction s'arrête dès que le confinement cesse \\
\hline
Maîtrise industrielle & acquise (centrales en service) & en cours (ITER), pas encore de production nette durable \\
\hline
\end{tabularx}
\end{center}

\begin{aretenir}[L'essentiel de la section]
\begin{itemize}
\item La \textbf{fusion} est l'union de deux noyaux légers en un noyau plus lourd, avec libération d'énergie $E = \Delta m\, c^2$ ; elle respecte les lois de conservation de Soddy.
\item Elle exige des températures de l'ordre de $10^8$ K pour vaincre la répulsion coulombienne : la matière est alors à l'état de \textbf{plasma}.
\item Réaction de référence : $\ce{^{2}_{1}H + ^{3}_{1}H -> ^{4}_{2}He + ^{1}_{0}n}$ ($\approx \SI{17,6}{MeV}$).
\item Dans les étoiles, le bilan global est $\ce{4\,^{1}_{1}H -> ^{4}_{2}He + 2\,^{0}_{1}e + 2\nu}$ : le Soleil est un réacteur à fusion naturel.
\item La bombe H (fusion non contrôlée) est amorcée par une bombe A ; la fusion contrôlée (tokamak, projet ITER) est un enjeu énergétique majeur : combustible abondant, déchets réduits, sûreté intrinsèque.
\end{itemize}
\end{aretenir}

\begin{exercice}[Pour s'entraîner]
On considère la réaction de fusion : $\ce{^{2}_{1}H + ^{3}_{1}H -> ^{4}_{2}He + ^{1}_{0}n}$.
Données : $m(\ce{^{2}_{1}H}) = \SI{2,01355}{u}$ ; $m(\ce{^{3}_{1}H}) = \SI{3,01550}{u}$ ; $m(\ce{^{4}_{2}He}) = \SI{4,00150}{u}$ ; $m_n = \SI{1,00866}{u}$ ; $\SI{1}{u} \cdot c^2 = \SI{931,5}{MeV}$.
\begin{enumerate}
\item Calculer le défaut de masse $\Delta m$ de la réaction.
\item En déduire l'énergie libérée en MeV, puis en joules.
\item Combien de réactions de fusion par seconde faudrait-il pour alimenter un quartier consommant une puissance de $\SI{100}{MW}$ ?
\end{enumerate}
\end{exercice}

\begin{corrige}[Exercice — Fusion deutérium-tritium]
\begin{enumerate}
\item Masse des réactifs : $2{,}01355 + 3{,}01550 = \SI{5,02905}{u}$. Masse des produits : $4{,}00150 + 1{,}00866 = \SI{5,01016}{u}$. D'où :
\[ \Delta m = 5{,}02905 - 5{,}01016 = \SI{0,01889}{u}. \]
\item Énergie libérée : $E = \Delta m\, c^2 = 0{,}01889 \times 931{,}5 \approx \SI{17,6}{MeV}$, soit en joules :
\[ E = 17{,}6\times 10^6 \times 1{,}60\times 10^{-19} \approx \SI{2,8e-12}{J}. \]
\item Nombre de réactions par seconde pour fournir $P = \SI{100}{MW} = \SI{e8}{W}$ :
\[ N = \frac{P}{E} = \frac{10^8}{2{,}8\times 10^{-12}} \approx \num{3,6e19} \text{ réactions par seconde}. \]
Cela correspond à une consommation de deutérium-tritium d'environ $5{,}03$ u par réaction, soit seulement quelques grammes par jour : la fusion est d'une sobriété en « combustible » extraordinaire.
\end{enumerate}
\end{corrige}

Nous disposons désormais du cadre énergétique complet des réactions nucléaires provoquées. La section suivante reprendra de manière systématique le bilan énergétique $E = \Delta m\, c^2$, avant que nous abordions les effets des rayonnements sur l'organisme et les méthodes de radioprotection.

\coursec{Aspect énergétique : $E = \Delta m \cdot c^2$}

Dans les sections précédentes, nous avons écrit les équations bilans de la fission de l'uranium-235 et de la fusion deutérium-tritium, en vérifiant à chaque fois les lois de conservation de Soddy (conservation du nombre de nucléons $A$ et de la charge $Z$). Mais une question fondamentale reste en suspens : \emph{d'où vient l'énergie considérable dégagée par ces réactions ?} Pourquoi une centrale nucléaire produit-elle de l'électricité avec quelques kilogrammes de combustible là où une centrale thermique classique brûle des wagons entiers de charbon ? La réponse, l'une des plus célèbres de toute la physique, tient en une équation : $E = \Delta m \cdot c^2$. C'est l'objet de cette section.

\courssub{Le constat expérimental : la masse a « disparu »}

Pesons, en pensée, les réactifs et les produits d'une réaction nucléaire. Les spectromètres de masse modernes mesurent les masses des noyaux avec une précision extraordinaire (mieux que $10^{-9}$ en valeur relative). Or, pour toutes les réactions nucléaires qui \emph{dégagent} de l'énergie (fission de l'uranium, fusion de l'hydrogène, désintégrations radioactives), on observe systématiquement le même fait :

\begin{center}
\emph{La masse totale des produits est inférieure à la masse totale des réactifs.}
\end{center}

Prenons l'exemple de la fusion deutérium-tritium, étudiée dans la section précédente :
\[
\ce{^{2}_{1}H} + \ce{^{3}_{1}H} \rightarrow \ce{^{4}_{2}He} + \ce{^{1}_{0}n}
\]
Les masses mesurées sont :
\begin{itemize}
\item $m(\ce{^{2}_{1}H}) = \num{2,0136}$~u et $m(\ce{^{3}_{1}H}) = \num{3,0155}$~u, soit une masse totale des réactifs $m_{\text{réactifs}} = \num{5,0291}$~u ;
\item $m(\ce{^{4}_{2}He}) = \num{4,0015}$~u et $m(n) = \num{1,0087}$~u, soit une masse totale des produits $m_{\text{produits}} = \num{5,0102}$~u.
\end{itemize}
Il manque environ $\num{0,0189}$~u ! Cette « masse perdue » n'a pas été détruite : elle a été \emph{convertie en énergie}, emportée par les produits sous forme d'énergie cinétique (et parfois de rayonnement $\gamma$). C'est précisément cette énergie qui chauffe l'eau des centrales ou fait briller le Soleil.

\begin{definition}[Défaut de masse d'une réaction nucléaire]
Le \cle{défaut de masse} $\Delta m$ d'une réaction nucléaire est la différence entre la masse totale des réactifs et la masse totale des produits, toutes mesurées au repos :
\[
\Delta m = m(\text{réactifs}) - m(\text{produits})
\]
Pour une réaction \cle{exoénergétique} (qui libère de l'énergie), on a $\Delta m > 0$. Le défaut de masse s'exprime en unité de masse atomique (u) ou en kilogramme (kg).
\end{definition}

\begin{attention}[Ne pas confondre avec le défaut de masse du noyau]
Il existe deux « défauts de masse » qu'il ne faut pas mélanger :
\begin{itemize}
\item le défaut de masse \emph{d'un noyau} : $\Delta m_{\text{noyau}} = Z\,m_p + (A-Z)\,m_n - m_{\text{noyau}}$, lié à l'énergie de liaison du noyau (courbe d'Aston) ;
\item le défaut de masse \emph{d'une réaction} : $\Delta m = m(\text{réactifs}) - m(\text{produits})$, lié à l'énergie libérée par la réaction.
\end{itemize}
Au Baccalauréat, c'est presque toujours le second qui intervient dans les calculs de $E = \Delta m c^2$. Lis bien l'énoncé pour savoir de quelles masses tu disposes.
\end{attention}

\begin{popfigure}
\begin{center}
\begin{tikzpicture}[scale=1, every node/.style={font=\small}]
% Plateau gauche : réactifs
\draw[thick, popdark] (-4.6,0) -- (-1.4,0);
\draw[thick, popdark] (-3,0) -- (-3,-0.5);
\draw[thick, popdark] (-3.6,-0.5) -- (-2.4,-0.5);
\fill[popblue] (-4.1,0.28) circle (0.28);
\fill[popblue!60] (-3.4,0.22) circle (0.22);
\fill[popblue!40] (-2.6,0.18) circle (0.18);
\fill[popblue!75] (-1.95,0.14) circle (0.14);
\node[popblue, below] at (-3,-0.65) {Réactifs : masse $m_1$};
% Flèche réaction
\draw[-{Stealth[length=3mm]}, very thick, poporange] (-1.1,0.3) -- (1.1,0.3);
\node[poporange, above] at (0,0.35) {réaction};
% Plateau droit : produits plus légers
\draw[thick, popdark] (1.4,0) -- (4.6,0);
\draw[thick, popdark] (3,0) -- (3,-0.5);
\draw[thick, popdark] (2.4,-0.5) -- (3.6,-0.5);
\fill[popgreen] (2.2,0.22) circle (0.22);
\fill[popgreen!60] (3.0,0.16) circle (0.16);
\fill[popgreen!40] (3.7,0.12) circle (0.12);
\node[popgreen!60!black, below] at (3,-0.65) {Produits : masse $m_2 < m_1$};
% Énergie libérée
\draw[-{Stealth[length=3mm]}, very thick, poppink] (3,0.55) -- (3,1.6);
\node[poppink, right, align=left] at (3.15,1.1) {$E = \Delta m\, c^2$\\ $\Delta m = m_1 - m_2 > 0$};
\end{tikzpicture}
\end{center}
\legende{Bilan de masse d'une réaction nucléaire exoénergétique : les produits sont plus légers que les réactifs ; la différence de masse $\Delta m$ est convertie en énergie.}
\end{popfigure}

\courssub{La relation d'Einstein : $E = \Delta m \cdot c^2$}

\begin{propriete}[Équivalence masse-énergie — relation d'Einstein]
Lors d'une réaction nucléaire exoénergétique, l'énergie libérée est proportionnelle au défaut de masse :
\[
\boxed{E_{\text{libérée}} = \Delta m \cdot c^2}
\]
où $\Delta m$ est le défaut de masse (en kg) et $c = \num{3,00e8}$~m.s$^{-1}$ la célérité de la lumière dans le vide. $E_{\text{libérée}}$ s'exprime en joules (J).
\par\smallskip
\emph{Cette relation est admise} : elle découle de la relativité restreinte d'Einstein (1905), dont la démonstration est hors programme. Elle traduit l'\cle{équivalence entre la masse et l'énergie} : la masse est une forme d'énergie « concentrée ».
\end{propriete}

Vérifions l'homogénéité : $[\Delta m \cdot c^2] = \text{kg} \times (\text{m.s}^{-1})^2 = \text{kg.m}^2\text{.s}^{-2} = \text{J}$. La relation est bien dimensionnellement cohérente.

Ce qui rend cette relation spectaculaire, c'est le facteur $c^2 \approx \num{9e16}$~m$^2$.s$^{-2}$, un nombre gigantesque. Une masse minuscule correspond à une énergie énorme : la conversion complète d'\emph{un seul gramme} de matière libérerait $E = 10^{-3} \times (\num{3e8})^2 = \num{9e13}$~J, soit l'équivalent de l'énergie électrique consommée par une ville comme Douala pendant plusieurs jours !

\courssub{Unités d'énergie en physique nucléaire}

Le joule est l'unité SI, mais il est mal adapté à l'échelle d'un noyau : les énergies mises en jeu par noyau sont de l'ordre de $10^{-13}$ à $10^{-11}$~J. On utilise donc l'électronvolt et son multiple le mégaélectronvolt.

\begin{aretenir}[Unités et conversions à connaître]
\begin{itemize}
\item \cle{Électronvolt} : $1~\text{eV} = \num{1,60e-19}$~J ; \quad $1~\text{MeV} = 10^6~\text{eV} = \num{1,60e-13}$~J.
\item \cle{Unité de masse atomique} : $1~\text{u} = \num{1,66054e-27}$~kg.
\item Équivalence masse-énergie de l'unité u : $1~\text{u} = 931{,}5~\text{MeV}/c^2$, c'est-à-dire qu'un défaut de masse de 1~u correspond à une énergie de $931{,}5$~MeV :
\[
E~(\text{MeV}) = \Delta m~(\text{u}) \times 931{,}5
\]
\end{itemize}
\end{aretenir}

D'où vient ce facteur $931{,}5$ ? Appliquons la relation d'Einstein à une masse de 1~u :
\[
E = \num{1,66054e-27} \times (\num{2,998e8})^2 = \num{1,4924e-10}~\text{J}
\]
puis convertissons en MeV : $E = \dfrac{\num{1,4924e-10}}{\num{1,602e-13}} = 931{,}5$~MeV. Retiens ce facteur de conversion, il est utilisé dans presque tous les exercices du Bac.

\courssub{Méthode de calcul d'une énergie libérée}

\begin{methode}[Calculer l'énergie libérée par une réaction nucléaire en 4 étapes]
\begin{enumerate}
\item \textbf{Écrire l'équation bilan} de la réaction en vérifiant les lois de conservation de Soddy (conservation de $A$ et de $Z$).
\item \textbf{Calculer le défaut de masse} $\Delta m = m(\text{réactifs}) - m(\text{produits})$, en unités de masse atomique (u). Vérifier que $\Delta m > 0$ pour une réaction exoénergétique.
\item \textbf{Convertir en MeV} : $E~(\text{MeV}) = \Delta m~(\text{u}) \times 931{,}5$.
\item \textbf{Convertir en joules} si demandé : $E~(\text{J}) = E~(\text{MeV}) \times \num{1,60e-13}$.
\end{enumerate}
\emph{Variante} : si les masses sont données en kg, appliquer directement $E = \Delta m \cdot c^2$ avec $c = \num{3,00e8}$~m.s$^{-1}$.
\end{methode}

\begin{attention}[Les deux erreurs classiques du Bac]
\begin{itemize}
\item \textbf{Oublier la conversion MeV $\rightarrow$ J} : si la question demande une énergie \emph{en joules} (par exemple pour calculer ensuite une puissance en watts), un résultat laissé en MeV est faux. $1~\text{MeV} = \num{1,60e-13}$~J, pas $10^6$~J !
\item \textbf{Mélanger les systèmes d'unités} : n'utilise \emph{jamais} $E = \Delta m \cdot c^2$ avec $\Delta m$ en u et $c = \num{3e8}$~m.s$^{-1}$ sans conversion préalable des u en kg. Soit tu convertis $\Delta m$ en kg ($1~\text{u} = \num{1,66e-27}$~kg), soit tu utilises le facteur $931{,}5$~MeV/u. Choisis une route et tiens-t'y.
\end{itemize}
\end{attention}

\courssub{Exemples chiffrés fondamentaux}

\begin{exemplebox}[Énergie libérée par la fission d'un noyau d'uranium-235]
Considérons la réaction de fission :
\[
\ce{^{235}_{92}U} + \ce{^{1}_{0}n} \rightarrow \ce{^{94}_{38}Sr} + \ce{^{139}_{54}Xe} + 3\,\ce{^{1}_{0}n}
\]
Vérification des conservations : nucléons $235 + 1 = 236$ et $94 + 139 + 3 = 236$ ; charges $92 = 38 + 54$. Les lois de Soddy sont bien respectées.
\par\smallskip
Le calcul du défaut de masse à partir des masses des noyaux donne $\Delta m \approx 0{,}193$~u. L'énergie libérée est donc :
\begin{align*}
E &= 0{,}193 \times 931{,}5 \approx 180~\text{MeV}\\
E &\approx 180 \times \num{1,60e-13} \approx \num{2,9e-11}~\text{J par noyau}
\end{align*}
(en comptant l'énergie emportée ensuite par les désintégrations des produits, on aboutit au chiffre usuel d'environ $200$~MeV par fission). Une seule fission libère environ $\num{3e-11}$~J : c'est microscopique, mais un gramme d'uranium-235 contient environ $\num{2,6e21}$ noyaux, soit une énergie totale de l'ordre de $\num{8e10}$~J par gramme — l'équivalent de plus d'une tonne de pétrole !
\end{exemplebox}

\begin{exemplebox}[Énergie libérée par la fusion deutérium-tritium]
Pour la réaction $\ce{^{2}_{1}H} + \ce{^{3}_{1}H} \rightarrow \ce{^{4}_{2}He} + \ce{^{1}_{0}n}$, le défaut de masse vaut :
\[
\Delta m = (2{,}0136 + 3{,}0155) - (4{,}0015 + 1{,}0087) = 0{,}0189~\text{u}
\]
d'où :
\[
E = 0{,}0189 \times 931{,}5 \approx 17{,}6~\text{MeV} \approx \num{2,8e-12}~\text{J}
\]
C'est moins que la fission en valeur absolue, mais rapportons à la masse : la fusion met en jeu 5 nucléons contre 236 pour la fission, soit environ $3{,}5$~MeV par nucléon contre $0{,}85$~MeV par nucléon pour la fission. \textbf{Par unité de masse, la fusion est environ 4 fois plus énergétique que la fission}, et des millions de fois plus que les réactions chimiques (quelques eV par réaction). C'est pourquoi le projet ITER et les étoiles misent sur la fusion.
\end{exemplebox}

\begin{aretenir}[Ordre de grandeur à retenir]
L'énergie libérée par une réaction nucléaire est environ \cle{$10^6$ fois} (un million de fois) plus grande que celle d'une réaction chimique rapportée au même nombre de particules : quelques MeV par noyau contre quelques eV par molécule. C'est ce facteur un million qui explique la compacité du combustible nucléaire... et la puissance destructrice des armes atomiques.
\end{aretenir}

\courssub{Lien avec la courbe d'Aston (complément Bac)}

Pourquoi la fission des noyaux lourds et la fusion des noyaux légers libèrent-elles toutes deux de l'énergie ? La réponse se lit sur la \cle{courbe d'Aston}, qui représente l'énergie de liaison par nucléon $\dfrac{E_\ell}{A}$ en fonction du nombre de masse $A$. Plus cette valeur est élevée, plus le noyau est stable. Le maximum (environ $8{,}8$~MeV/nucléon) se situe autour du fer ($A \approx 56$).

\begin{itemize}
\item \textbf{Fission} : un noyau lourd ($A > 200$, faible $E_\ell/A$) se scinde en deux noyaux moyens, plus proches du fer, donc plus liés. L'augmentation de $E_\ell/A$ se traduit par une libération d'énergie.
\item \textbf{Fusion} : deux noyaux très légers ($A < 20$, faible $E_\ell/A$) s'unissent en un noyau plus lourd, plus lié. Là encore, $E_\ell/A$ augmente et l'énergie est libérée.
\end{itemize}

\begin{popfigure}
\begin{center}
\begin{tikzpicture}
\begin{axis}[
width=0.92\linewidth, height=6.2cm,
xlabel={$A$ (nombre de masse)}, ylabel={$E_\ell/A$ (MeV/nucléon)},
xmin=0, xmax=250, ymin=0, ymax=10,
xtick={0,50,100,150,200,250}, ytick={0,2,4,6,8,10},
grid=major, grid style={popline!40},
legend style={draw=none, fill=none, font=\small}
]
\addplot[popcurve, domain=2:245, samples=150] {8.8 - 8.5*exp(-x/12) - 0.0095*(x-56>0 ? x-56 : 0)};
\addplot[only marks, mark=*, mark size=1.6pt, poppink] coordinates {(2,1.1) (4,7.1) (12,7.7) (56,8.8) (235,7.6)};
\node[font=\small, poppink, anchor=south] at (axis cs:56,8.8) {Fe};
\node[font=\small, poppink, anchor=south west] at (axis cs:228,7.7) {U};
\node[font=\small, poppink, anchor=south] at (axis cs:2,1.4) {H};
% Zones
\draw[-{Stealth}, very thick, popgreen] (axis cs:8,2.2) -- (axis cs:35,6.3);
\node[font=\small, popgreen!60!black, align=center] at (axis cs:22,3.4) {\textbf{fusion}\\ zone $A$ petit};
\draw[-{Stealth}, very thick, poporange] (axis cs:238,6.6) -- (axis cs:150,7.9);
\node[font=\small, poporange!80!black, align=center] at (axis cs:196,6.0) {\textbf{fission}\\ zone $A$ grand};
\end{axis}
\end{tikzpicture}
\end{center}
\legende{Courbe d'Aston simplifiée (complément) : énergie de liaison par nucléon en fonction de $A$. Les réactions qui font « monter » vers le maximum (fer) libèrent de l'énergie : fusion des noyaux légers, fission des noyaux lourds.}
\end{popfigure}

\courssub{Application : puissance d'une centrale et consommation de combustible}

\begin{exoresolu}[Combustible d'une centrale nucléaire]
\textbf{Énoncé.} Une centrale nucléaire fournit une puissance électrique $P = 900$~MW avec un rendement global $\eta = 33\,\%$. Chaque fission d'uranium-235 libère $E_1 = 200$~MeV. On donne $M(\ce{^{235}U}) = 235$~g.mol$^{-1}$, $\mathcal{N}_A = \num{6,02e23}$~mol$^{-1}$, $1~\text{MeV} = \num{1,60e-13}$~J. Calculer la masse d'uranium-235 consommée par jour de fonctionnement.
\tcblower
\textbf{Solution détaillée.}
\par\smallskip
\textbf{Étape 1 : puissance thermique nécessaire.} Le rendement étant de $33\,\%$ :
\[
P_{\text{th}} = \frac{P}{\eta} = \frac{\num{900e6}}{0{,}33} \approx \num{2,7e9}~\text{W}
\]
\textbf{Étape 2 : énergie thermique pour un jour.} Avec $\Delta t = 24 \times 3600 = \num{8,64e4}$~s :
\[
E_{\text{th}} = P_{\text{th}} \times \Delta t = \num{2,7e9} \times \num{8,64e4} \approx \num{2,35e14}~\text{J}
\]
\textbf{Étape 3 : énergie par fission en joules.}
\[
E_1 = 200~\text{MeV} = 200 \times \num{1,60e-13} = \num{3,2e-11}~\text{J}
\]
\textbf{Étape 4 : nombre de fissions puis masse.}
\[
N = \frac{E_{\text{th}}}{E_1} = \frac{\num{2,35e14}}{\num{3,2e-11}} \approx \num{7,3e24}~\text{fissions}
\]
La masse correspondante :
\[
m = \frac{N}{\mathcal{N}_A} \times M = \frac{\num{7,3e24}}{\num{6,02e23}} \times 235 \approx \num{2,9e3}~\text{g} \approx 2{,}9~\text{kg}
\]
\textbf{Conclusion.} Environ \textbf{3 kg d'uranium-235 par jour} suffisent à alimenter une centrale de 900~MW, là où une centrale à charbon de même puissance brûlerait plusieurs milliers de tonnes de charbon. Voilà l'illustration concrète du facteur $10^6$ entre énergies nucléaire et chimique.
\end{exoresolu}

\begin{savaistu}[Le Soleil maigrit de 4 millions de tonnes chaque seconde]
Dans le cœur du Soleil, la fusion de l'hydrogène en hélium libère une puissance d'environ $\num{3,8e26}$~W. D'après la relation d'Einstein, la masse convertie en énergie chaque seconde est :
\[
\Delta m = \frac{E}{c^2} = \frac{\num{3,8e26}}{(\num{3e8})^2} \approx \num{4,2e9}~\text{kg}
\]
soit environ \textbf{4 millions de tonnes de masse perdues par seconde}, rayonnées dans l'espace sous forme de lumière et de chaleur — dont une infime partie fait vivre la Terre, fait pousser le manioc et le maïs de nos champs et alimente les panneaux solaires de l'électrification rurale camerounaise. Et pourtant, à ce rythme, le Soleil peut briller encore environ 5 milliards d'années !
\end{savaistu}

\begin{exercice}[Fusion dans ITER]
On donne : $m(\ce{^{2}_{1}H}) = \num{2,0136}$~u ; $m(\ce{^{3}_{1}H}) = \num{3,0155}$~u ; $m(\ce{^{4}_{2}He}) = \num{4,0015}$~u ; $m_n = \num{1,0087}$~u ; $1~\text{u} = 931{,}5$~MeV/$c^2$ ; $1~\text{MeV} = \num{1,60e-13}$~J.
\begin{enumerate}
\item Calculer le défaut de masse de la réaction $\ce{^{2}_{1}H} + \ce{^{3}_{1}H} \rightarrow \ce{^{4}_{2}He} + \ce{^{1}_{0}n}$.
\item En déduire l'énergie libérée en MeV puis en joules.
\item Calculer l'énergie libérée par la fusion de 1~g de mélange deutérium-tritium (on prendra $N = \num{1,5e23}$ réactions pour 1~g de mélange).
\end{enumerate}
\end{exercice}

\begin{corrige}[Exercice : Fusion dans ITER]
\begin{enumerate}
\item $\Delta m = (2{,}0136 + 3{,}0155) - (4{,}0015 + 1{,}0087) = 5{,}0291 - 5{,}0102 = 0{,}0189$~u.
\item $E = 0{,}0189 \times 931{,}5 \approx 17{,}6$~MeV, soit $E = 17{,}6 \times \num{1,60e-13} \approx \num{2,8e-12}$~J.
\item $E_{\text{tot}} = N \times E = \num{1,5e23} \times \num{2,8e-12} \approx \num{4,2e11}$~J par gramme, soit l'équivalent énergétique d'environ 10 tonnes de pétrole : la fusion est bien la source d'énergie la plus concentrée accessible à l'humanité.
\end{enumerate}
\end{corrige}

En résumé, le défaut de masse $\Delta m = m(\text{réactifs}) - m(\text{produits}) > 0$ est la signature d'une réaction nucléaire exoénergétique, et la relation d'Einstein $E = \Delta m \cdot c^2$ en donne l'énergie libérée. Mais cette énergie se manifeste sous forme de rayonnements et de particules rapides : quels sont leurs effets sur la matière vivante, et comment s'en protéger ? C'est l'objet des deux dernières sections de cette leçon.

\coursec{Effets biologiques des rayonnements et dosimétrie}

Dans la section précédente, nous avons vu qu'une réaction nucléaire provoquée — fission de l'uranium 235 ou fusion de noyaux légers — libère une énergie considérable, $E = \Delta m\,c^{2}$, sous forme d'énergie cinétique des fragments et de rayonnements. Mais que deviennent ces rayonnements lorsqu'ils rencontrent la matière vivante, par exemple le corps d'un technicien d'une centrale ou celui d'un patient en radiothérapie à l'hôpital général de Douala ? C'est toute la question de cette section : comprendre comment les rayonnements interagissent avec nos cellules, mesurer l'exposition à l'aide de grandeurs adaptées (le gray, le sievert), et distinguer deux situations très différentes que le langage courant confond trop souvent : l'\cle{irradiation} et la \cle{contamination}.

\courssub{Les rayonnements ionisants : nature, pouvoir ionisant et pouvoir pénétrant}

\textbf{Quatre types de rayonnements.} Les désintégrations et réactions nucléaires émettent des projectiles de natures très différentes :

\begin{itemize}
\item le rayonnement \cle{alpha} ($\alpha$) : noyaux d'hélium \ce{^{4}_{2}He}, chargés positivement ($q = +2e$), massifs ($m \approx 4\ \mathrm{u}$) ;
\item le rayonnement \cle{bêta} ($\beta$) : électrons \ce{^{0}_{-1}e} (ou positrons), légers, chargés ;
\item le rayonnement \cle{gamma} ($\gamma$) : photons de très haute énergie, sans masse ni charge ;
\item les \cle{neutrons} \ce{^{1}_{0}n} : particules neutres, émis en abondance lors des fissions.
\end{itemize}

\textbf{Pouvoir ionisant.} En traversant la matière, ces particules arrachent des électrons aux atomes rencontrés : on dit qu'elles \emph{ionisent} le milieu. C'est précisément cette ionisation qui les rend dangereuses pour le vivant, d'où le nom de \cle{rayonnements ionisants}. Le pouvoir ionisant dépend de la charge et de la masse : une particule $\alpha$, lourde et doublement chargée, interagit violemment avec les électrons des atomes et ionise énormément sur un court trajet ; un photon $\gamma$, sans charge, ionise peu mais sur un long parcours.

\textbf{Pouvoir pénétrant.} C'est l'inverse : plus une particule ionise fortement, plus elle perd vite son énergie et moins elle pénètre. On retiendra la hiérarchie :

\begin{center}
\begin{tabularx}{\linewidth}{|l|X|X|X|}
\hline
\rowcolor{popblueL} \textbf{Rayonnement} & \textbf{Arrêté par} & \textbf{Portée dans l'air} & \textbf{Dans les tissus} \\
\hline
$\alpha$ & une feuille de papier, la peau morte & quelques cm & quelques dizaines de $\mu$m \\
\hline
$\beta$ & une feuille d'aluminium de quelques mm & quelques m & quelques mm \\
\hline
$\gamma$ & atténué par du plomb épais ou du béton & plusieurs centaines de m & traverse le corps \\
\hline
neutrons & eau, paraffine, béton & grande & traverse le corps \\
\hline
\end{tabularx}
\end{center}

\begin{popfigure}
\begin{tikzpicture}[scale=0.95, >=Stealth]
% écran papier
\fill[popgoldL] (2.6,-0.5) rectangle (2.9,2.6);
\node[rotate=90, font=\footnotesize] at (2.75,1.05) {papier};
% écran aluminium
\fill[popmuted!60] (5.4,-0.5) rectangle (5.8,2.6);
\node[rotate=90, font=\footnotesize] at (5.6,1.05) {aluminium};
% écran plomb
\fill[popdark!70] (8.2,-0.5) rectangle (9.0,2.6);
\node[rotate=90, font=\footnotesize, text=white] at (8.6,1.05) {plomb};
% source
\node[circle, fill=popink, inner sep=2.5pt, label=left:\footnotesize source] (S) at (0,1.05) {};
% alpha
\draw[->, very thick, poppink] (S) -- (2.55,2.2) node[midway, above, font=\footnotesize] {$\alpha$};
\node[font=\footnotesize, poppink] at (3.6,2.2) {arrêté};
% beta
\draw[->, very thick, poporange] (S) -- (5.35,1.05) node[midway, above, font=\footnotesize] {$\beta$};
\node[font=\footnotesize, poporange] at (6.9,1.05) {arrêté};
% gamma
\draw[->, very thick, popblue] (S) -- (10.6,-0.1) node[midway, above, font=\footnotesize, pos=0.35] {$\gamma$};
\draw[->, very thick, popblue!45] (9.0,-0.1) -- (10.6,-0.1);
\node[font=\footnotesize, popblue] at (10.9,-0.1) {atténué};
\end{tikzpicture}
\legende{Pouvoir pénétrant comparé : les $\alpha$ sont stoppés par le papier, les $\beta$ par quelques millimètres d'aluminium, les $\gamma$ ne sont qu'atténués par une épaisseur importante de plomb ou de béton.}
\end{popfigure}

\begin{attention}[Ne pas confondre les deux pouvoirs]
Le rayonnement $\alpha$ a le \emph{plus grand pouvoir ionisant} mais le \emph{plus faible pouvoir pénétrant}. Le $\gamma$ est le plus pénétrant mais le moins ionisant. Conséquence pratique : une source $\alpha$ extérieure au corps est presque inoffensive (arrêtée par la couche de peau morte), mais une source $\alpha$ \emph{avalée ou inhalée} dépose toute son énergie dans quelques cellules : elle devient extrêmement dangereuse.
\end{attention}

\courssub{Mécanisme d'action sur la matière vivante}

Une cellule vivante est constituée d'environ 70 \% d'eau et de molécules organiques, dont la plus précieuse est l'\cle{ADN}, support de l'information génétique. Lorsqu'un rayonnement ionisant traverse un tissu, deux mécanismes se combinent :

\begin{itemize}
\item \textbf{Action indirecte} : le rayonnement ionise les molécules d'eau, produisant des \emph{radicaux libres} (espèces chimiques très réactives comme OH$^{\bullet}$). Ces radicaux attaquent ensuite les molécules voisines, notamment l'ADN.
\item \textbf{Action directe} : le rayonnement ionise ou brise directement une molécule d'ADN (cassure d'un ou des deux brins).
\end{itemize}

La cellule possède des enzymes de réparation, mais si les lésions sont trop nombreuses ou mal réparées, trois issues sont possibles : la mort de la cellule, sa réparation correcte, ou sa \emph{transformation} en cellule mutée capable de se diviser de façon incontrôlée — c'est le point de départ d'un cancer.

\textbf{Effets déterministes et effets stochastiques.} On distingue deux grandes catégories d'effets :

\begin{itemize}
\item Les \cle{effets déterministes} apparaissent au-dessus d'un \textbf{seuil de dose} et leur gravité augmente avec la dose : rougeurs et brûlures de la peau, nausées, chute des globules blancs, et au-delà d'environ 1 Sv reçu en une seule fois, le \emph{syndrome d'irradiation aiguë} pouvant conduire à la mort. C'est ce qui a frappé les liquidateurs de Tchernobyl.
\item Les \cle{effets stochastiques} (aléatoires) sont des effets \textbf{sans seuil} : même une faible dose a une probabilité non nulle de provoquer, des années plus tard, un cancer ou une mutation génétique transmissible. On applique donc le principe de précaution : toute exposition inutile doit être évitée.
\end{itemize}

\courssub{Irradiation et contamination : une distinction fondamentale}

\begin{definition}[Irradiation et contamination]
\begin{itemize}
\item L'\cle{irradiation} est l'exposition du corps à un rayonnement émis par une source \textbf{extérieure}. Elle cesse dès que l'on s'éloigne de la source ou que celle-ci est retirée. Une personne irradiée \textbf{n'est pas radioactive} : elle ne présente aucun danger pour son entourage.
\item La \cle{contamination} est la présence de \textbf{matière radioactive} sur le corps (contamination \emph{externe} : peau, cheveux, vêtements) ou à l'intérieur du corps (contamination \emph{interne} : par \emph{inhalation} de poussières ou gaz radioactifs, ou par \emph{ingestion} d'aliments ou d'eau contaminés). Elle persiste tant que la matière radioactive n'a pas été éliminée ou désintégrée. Une personne contaminée \textbf{est elle-même radioactive} et peut contaminer son entourage.
\end{itemize}
\end{definition}

\begin{popfigure}
\begin{tikzpicture}[scale=0.9, >=Stealth]
% ---- Irradiation ----
\node[font=\bfseries\small] at (2.2,3.1) {Irradiation (source externe)};
\node[circle, fill=popink, inner sep=3pt, label=above:\footnotesize source] (S) at (0.4,1.2) {};
% personne
\draw[thick, popdark] (3.6,0.4) circle (0.28);
\draw[thick, popdark] (3.6,0.12) -- (3.6,-0.9);
\draw[thick, popdark] (3.0,-0.35) -- (4.2,-0.35);
\draw[thick, popdark] (3.6,-0.9) -- (3.15,-1.7);
\draw[thick, popdark] (3.6,-0.9) -- (4.05,-1.7);
\foreach \y in {0.6,1.2,1.8} {\draw[->, thick, poppink] (S) ++(0.15,0) -- ++(2.4,\y-1.2);}
\node[font=\footnotesize, poppink, align=center] at (2.2,-2.1) {le rayonnement cesse\\ si l'on s'éloigne};
% ---- Contamination ----
\node[font=\bfseries\small] at (8.4,3.1) {Contamination (matière radioactive)};
\draw[thick, popdark] (8.4,0.4) circle (0.28);
\draw[thick, popdark] (8.4,0.12) -- (8.4,-0.9);
\draw[thick, popdark] (7.8,-0.35) -- (9.0,-0.35);
\draw[thick, popdark] (8.4,-0.9) -- (7.95,-1.7);
\draw[thick, popdark] (8.4,-0.9) -- (8.85,-1.7);
\foreach \p in {(8.15,0.55),(8.6,-0.2),(8.3,-0.6),(8.75,-1.2),(8.05,-1.4)} {\node[circle, fill=popgreen, inner sep=1.6pt] at \p {};}
\node[font=\footnotesize, popgreen!60!black, align=center] at (8.4,-2.1) {particules déposées sur/dans le corps :\\ l'exposition persiste};
\end{tikzpicture}
\legende{À gauche, l'irradiation : la source est à distance, l'exposition s'arrête quand on s'éloigne. À droite, la contamination : des particules radioactives (points verts) sont fixées sur la peau ou incorporées, l'exposition continue.}
\end{popfigure}

\begin{attention}[L'erreur classique du Baccalauréat et des médias]
Confondre irradiation et contamination est l'erreur la plus fréquente. Un patient qui a reçu une radiographie pulmonaire a été \emph{irradié} : en sortant du cabinet, il n'émet aucun rayonnement. En revanche, un travailleur dont les vêtements sont souillés de poussières radioactives est \emph{contaminé} : il faut le doucher, changer ses vêtements et traiter ses déchets. Retiens : \textbf{irradié $\neq$ radioactif ; contaminé $=$ radioactif}.
\end{attention}

\courssub{Dosimétrie : la dose absorbée (gray) et la dose équivalente (sievert)}

Pour protéger les personnes, il faut d'abord \emph{mesurer} l'exposition. L'idée naturelle est de compter l'énergie que le rayonnement dépose dans les tissus.

\begin{definition}[Dose absorbée]
La \cle{dose absorbée} $D$ par un tissu est l'énergie $E$ cédée par le rayonnement à ce tissu, rapportée à la masse $m$ du tissu :
\[ D = \dfrac{E}{m} \]
Son unité SI est le \cle{gray} (Gy) : $\SI{1}{Gy} = \SI{1}{J.kg^{-1}}$.
\end{definition}

Vérifions l'homogénéité : $[D] = \dfrac{[E]}{[m]} = \mathrm{J \cdot kg^{-1}}$, ce qui correspond bien au gray. Une dose de 1 Gy signifie que chaque kilogramme de tissu a absorbé 1 joule d'énergie de rayonnement. Cela peut sembler dérisoire — 1 J correspond au travail du poids d'un objet de 100 g soulevé d'environ 1 m ($mgh \approx 0{,}1 \times 10 \times 1 = \SI{1}{J}$) — mais cette énergie, concentrée à l'échelle des molécules d'ADN, suffit à provoquer des millions de cassures par cellule.

\textbf{Pourquoi le gray ne suffit pas.} L'expérience montre qu'à \emph{même dose absorbée}, tous les rayonnements ne sont pas également nocifs : 1 Gy de rayonnement $\alpha$ endommage beaucoup plus les tissus que 1 Gy de photons $\gamma$, car les $\alpha$ déposent leur énergie de façon très concentrée sur un trajet microscopique. On introduit donc une dose « corrigée » de l'efficacité biologique.

\begin{definition}[Dose équivalente]
La \cle{dose équivalente} $H$ mesure la dangerosité biologique réelle de l'exposition. Son unité est le \cle{sievert} (Sv), avec $\SI{1}{Sv} = \SI{1}{J.kg^{-1}}$ (même dimension que le gray, mais signification différente).
\end{definition}

\begin{propriete}[Relation entre dose équivalente et dose absorbée — admise]
\[ H = D \times Q \]
où $Q$ est le \cle{facteur de qualité} (ou facteur de pondération) du rayonnement, sans dimension :
\begin{center}
\begin{tabular}{|c|c|}
\hline
\rowcolor{popblueL} \textbf{Rayonnement} & \textbf{Facteur $Q$} \\
\hline
$\gamma$, $\beta$, rayons X & 1 \\
\hline
neutrons & 5 à 20 selon l'énergie \\
\hline
$\alpha$ & 20 \\
\hline
\end{tabular}
\end{center}
Cette relation est admise au niveau Terminale ; elle doit être \emph{exploitée numériquement}. À dose absorbée égale, un rayonnement $\alpha$ est 20 fois plus nocif qu'un rayonnement $\gamma$.
\end{propriete}

\begin{methode}[Calculer une dose équivalente en trois étapes]
\begin{enumerate}
\item \textbf{Convertir} toutes les données en unités SI : énergie en joules, masse en kilogrammes.
\item \textbf{Calculer la dose absorbée} $D = \dfrac{E}{m}$ en gray.
\item \textbf{Pondérer} par le facteur de qualité du rayonnement : $H = D \times Q$, en sievert, puis comparer aux repères (dose naturelle, limites réglementaires).
\end{enumerate}
\end{methode}

\begin{exemplebox}[Calcul d'une dose équivalente]
Un tissu pulmonaire absorbe une dose $D = \SI{0,05}{Gy}$ de rayonnement $\alpha$ ($Q = 20$) provenant de poussières inhalées. La dose équivalente vaut :
\[ H = D \times Q = 0{,}05 \times 20 = \SI{1}{Sv} \]
C'est une dose considérable, de l'ordre du seuil d'apparition d'effets graves. Remarque que la même énergie déposée par des photons $\gamma$ n'aurait donné que $H = \SI{0,05}{Sv}$ : le facteur de qualité change tout.
\end{exemplebox}

\begin{exoresolu}[Radiographie ou scanner : quelle dose ?]
\textbf{Énoncé.} Lors d'un examen, un organe de masse $m = \SI{2,0}{kg}$ absorbe une énergie $E = \SI{4,0e-4}{J}$ de rayons X ($Q = 1$). (1) Calculer la dose absorbée $D$. (2) En déduire la dose équivalente $H$ et la comparer à la dose naturelle annuelle moyenne de $\SI{2,4}{mSv}$.
\tcblower
\textbf{Solution.} (1) Par définition de la dose absorbée :
\[ D = \dfrac{E}{m} = \dfrac{4{,}0 \times 10^{-4}}{2{,}0} = \SI{2,0e-4}{Gy} = \SI{0,20}{mGy} \]
(2) Pour des rayons X, $Q = 1$, donc :
\[ H = D \times Q = 0{,}20 \times 1 = \SI{0,20}{mSv} \]
Soit environ $\dfrac{0{,}20}{2{,}4} \approx 8\ \%$ de la dose naturelle annuelle, soit l'équivalent d'un mois d'exposition naturelle : un examen médical reste une exposition faible, mais qui s'ajoute et doit être justifiée.
\end{exoresolu}

\courssub{Ordres de grandeur et échelle des doses}

Le sievert étant une unité « grande » à l'échelle des expositions courantes, on utilise le plus souvent le millisievert (mSv) et le microsievert ($\mu$Sv). Quelques repères à connaître :

\begin{itemize}
\item radioactivité naturelle moyenne (sol, radon, cosmos, alimentation) : $\approx 2$ à $\SI{3}{mSv/an}$ ;
\item radiographie pulmonaire : $\approx \SI{0,1}{mSv}$ ; scanner abdominal : $\approx \SI{10}{mSv}$ ;
\item limite réglementaire pour le public (hors exposition naturelle et médicale) : $\SI{1}{mSv/an}$ ;
\item seuil d'apparition d'effets déterministes graves : $\approx \SI{1}{Sv}$ en exposition brève ;
\item dose létale pour 50 \% des personnes exposées sans soins : $\approx 4$ à $\SI{5}{Sv}$.
\end{itemize}

\begin{popfigure}
\begin{tikzpicture}[scale=1.0, >=Stealth]
% axe logarithmique de 1e-5 à 1e1 Sv
\draw[->, very thick, popdark] (0,0) -- (12.6,0) node[right, font=\footnotesize] {dose équivalente (Sv)};
\foreach \x/\l in {0/{$10^{-5}$},2/{$10^{-4}$},4/{$10^{-3}$},6/{$10^{-2}$},8/{$10^{-1}$},10/{$10^{0}$},12/{$10^{1}$}} {
  \draw[thick, popdark] (\x,0.12) -- (\x,-0.12);
  \node[below, font=\footnotesize] at (\x,-0.15) {\l};
}
% repères
\draw[->, thick, popgreen] (3.4,0) -- (3.4,0.9); \node[above, font=\footnotesize, popgreen!60!black, align=center] at (3.4,0.95) {radio naturelle :\\ $\approx 2{,}4$ mSv/an};
\draw[->, thick, popblue] (4.0,0) -- (4.0,1.7); \node[above, font=\footnotesize, popblue, align=center] at (4.0,1.75) {radiographie :\\ $\approx 0{,}1$ mSv};
\draw[->, thick, poporange] (6.0,0) -- (6.0,0.9); \node[above, font=\footnotesize, poporange, align=center] at (6.0,0.95) {scanner :\\ $\approx 10$ mSv};
\draw[->, thick, poppink] (10.0,0) -- (10.0,1.7); \node[above, font=\footnotesize, poppink, align=center] at (10.0,1.75) {seuil d'effets\\ graves : $\approx 1$ Sv};
\draw[->, thick, poppurple] (11.3,0) -- (11.3,0.9); \node[above, font=\footnotesize, poppurple, align=center] at (11.3,0.95) {dose létale 50\,\% :\\ $\approx 4$--$5$ Sv};
\end{tikzpicture}
\legende{Échelle logarithmique des doses équivalentes. Entre la radiographie et le seuil d'effets graves, il y a un facteur $10^{4}$ : l'échelle logarithmique est indispensable pour tout représenter.}
\end{popfigure}

\begin{attention}[Bien choisir l'unité et le préfixe]
Aux examens, les doses sont souvent données en mSv ou $\mu$Sv : convertis toujours en sievert avant tout calcul ($\SI{1}{mSv} = \SI{e-3}{Sv}$, $\SI{1}{\mu Sv} = \SI{e-6}{Sv}$). Autre piège : ne jamais additionner des grays et des sieverts — ce ne sont pas les mêmes grandeurs, même si leurs unités ont la même dimension.
\end{attention}

\begin{savaistu}[Le radon, première source d'exposition naturelle]
Le \cle{radon} 222 est un gaz radioactif issu de la désintégration de l'uranium naturellement présent dans les roches. Il s'accumule dans les habitations mal ventilées, surtout celles bâties sur des sols granitiques. Inhalé, il dépose dans les poumons des descendants solides émetteurs $\alpha$ : c'est la \textbf{première source d'exposition naturelle} des populations (environ la moitié de la dose naturelle moyenne) et, dans plusieurs pays, la deuxième cause de cancer du poumon après le tabac. Au Cameroun, les régions volcaniques et granitiques de l'Ouest et de l'Adamaoua appellent une vigilance particulière : une simple aération quotidienne des pièces réduit fortement la concentration en radon.
\end{savaistu}

\begin{exercice}[Dose reçue par un technicien]
Un technicien manipule une source émettant des neutrons. Une masse $m = \SI{0,50}{kg}$ de tissu de sa main absorbe une énergie $E = \SI{1,0e-3}{J}$. On prend $Q = 10$ pour ces neutrons. (1) Calculer la dose absorbée $D$. (2) Calculer la dose équivalente $H$ et commenter au regard du seuil d'effets déterministes.
\end{exercice}

\begin{corrige}[Exercice : dose reçue par un technicien]
(1) Dose absorbée :
\[ D = \dfrac{E}{m} = \dfrac{1{,}0 \times 10^{-3}}{0{,}50} = \SI{2,0e-3}{Gy} = \SI{2,0}{mGy} \]
(2) Dose équivalente :
\[ H = D \times Q = 2{,}0 \times 10^{-3} \times 10 = \SI{2,0e-2}{Sv} = \SI{20}{mSv} \]
Cette dose reste très inférieure au seuil d'effets déterministes ($\approx \SI{1}{Sv}$), mais elle représente environ 20 fois la limite annuelle réglementaire pour le public : une telle exposition doit déclencher une enquête et un renforcement des protections (écrans, distance, temps de manipulation réduit).
\end{corrige}

\begin{aretenir}[L'essentiel de la section]
\begin{itemize}
\item Les rayonnements $\alpha$, $\beta$, $\gamma$ et les neutrons ionisent la matière vivante, directement ou via les molécules d'eau, et peuvent léser l'ADN.
\item Effets \textbf{déterministes} (avec seuil, gravité croissante) et effets \textbf{stochastiques} (sans seuil : cancers, mutations).
\item \textbf{Irradiation} = exposition à une source externe, qui cesse quand on s'éloigne ; \textbf{contamination} = matière radioactive sur ou dans le corps, qui persiste.
\item Dose absorbée $D = \dfrac{E}{m}$, en gray ($\SI{1}{Gy} = \SI{1}{J.kg^{-1}}$) ; dose équivalente $H = D \times Q$, en sievert, avec $Q = 1$ pour $\beta,\gamma$ et $Q = 20$ pour $\alpha$.
\item Repères : radioactivité naturelle $\approx 2$ à $3$ mSv/an, radiographie $\approx 0{,}1$ mSv, effets graves à partir de $\approx 1$ Sv.
\end{itemize}
\end{aretenir}

Maintenant que nous savons quantifier l'exposition et ses dangers, une question pratique s'impose : comment s'en protéger concrètement ? C'est l'objet de la dernière section, où nous verrons les trois piliers de la radioprotection — la distance, le temps d'exposition et les écrans — et la loi d'atténuation exponentielle $N(x) = N_{0}\,e^{-\mu x}$ d'un faisceau de photons.

\coursec{Radioprotection et atténuation d'un faisceau de photons}

Dans la section précédente, nous avons vu que les rayonnements ionisants déposent de l'énergie dans les tissus vivants et que la \cle{dose absorbée} (en gray) ainsi que la \cle{dose équivalente} (en sievert) permettent de quantifier le risque biologique. Une question pratique se pose alors immédiatement : comment se protéger concrètement de ces rayonnements ? Comment le radiologue de l'hôpital général de Douala, le technicien du centre de médecine nucléaire de Yaoundé, ou le soudeur qui contrôle des pipelines par gammagraphie peuvent-ils travailler sans danger ? Toute la réponse tient en trois règles simples — le \cle{temps}, la \cle{distance}, l'\cle{écran} — et en une loi expérimentale remarquable : l'atténuation exponentielle d'un faisceau de photons traversant la matière.

\courssub{Les trois principes de la radioprotection}

\textbf{Premier principe : diminuer le temps d'exposition.}\par
La dose absorbée $D$ est l'énergie déposée par unité de masse de tissu. À débit de dose $\dot{D}$ constant, l'énergie reçue — donc la dose — est \emph{proportionnelle à la durée} $\Delta t$ d'exposition :
\[ D = \dot{D}\times \Delta t. \]
Diviser le temps d'exposition par deux divise la dose par deux. C'est pourquoi les manipulateurs en médecine nucléaire préparent leurs gestes à blanc, travaillent vite et ne s'attardent jamais près d'une source inutilement.

\textbf{Deuxième principe : augmenter la distance à la source.}\par
Pour une source ponctuelle émettant de façon isotrope (dans toutes les directions), les photons se répartissent uniformément sur la surface d'une sphère de rayon $d$, dont l'aire vaut $4\pi d^{2}$. Le nombre de photons reçus par unité de surface — et donc le débit de dose — varie donc comme l'inverse du carré de la distance :
\[ \dot{D}(d) \propto \frac{1}{d^{2}} \qquad \text{(loi en } 1/d^{2}\text{, complément Bac).} \]
Doubler la distance divise le débit de dose par quatre ; le multiplier par dix le divise par cent. C'est pourquoi on manipule les sources radioactives avec des \cle{pinces} et des tiges longues, et pourquoi le simple fait de s'éloigner de quelques mètres d'une source est déjà une protection spectaculaire.

\begin{exemplebox}[La loi en $1/d^{2}$ en action]
Un technicien reçoit un débit de dose de $\dot{D}_{1} = \SI{0,80}{mSv/h}$ à $d_{1} = \SI{1,0}{m}$ d'une source de gammagraphie. À quelle distance le débit de dose tombe-t-il à $\dot{D}_{2} = \SI{0,020}{mSv/h}$ ?
Comme $\dot{D}\propto 1/d^{2}$, on a $\dot{D}_{1}d_{1}^{2} = \dot{D}_{2}d_{2}^{2}$, d'où
\[ d_{2} = d_{1}\sqrt{\frac{\dot{D}_{1}}{\dot{D}_{2}}} = 1{,}0\times\sqrt{\frac{0{,}80}{0{,}020}} = \sqrt{40} \approx \SI{6,3}{m}. \]
Sans aucun matériel, simplement en reculant, on a divisé l'exposition par 40.
\end{exemplebox}

\textbf{Troisième principe : interposer un écran.}\par
Lorsqu'on ne peut ni raccourcir le temps ni augmenter la distance, on place un matériau absorbant entre la source et la personne. Le choix du matériau dépend du \emph{pouvoir de pénétration} du rayonnement :

\begin{center}
\begin{tabularx}{\linewidth}{|l|X|X|}
\hline
\rowcolor{popblueL} \textbf{Rayonnement} & \textbf{Pouvoir pénétrant} & \textbf{Écran adapté} \\
\hline
$\alpha$ (\ce{^{4}_{2}He}) & Très faible : quelques centimètres d'air & Une feuille de papier, la peau suffisent \\
\hline
$\beta$ (électrons) & Moyen : quelques mètres d'air & Quelques mm d'aluminium ou de plexiglas \\
\hline
$\gamma$ et X (photons) & Très élevé : très pénétrant & Plomb épais ou béton \\
\hline
\end{tabularx}
\end{center}

\begin{attention}[Pourquoi du plexiglas plutôt que du plomb pour les $\beta$ ?]
Arrêter brutalement des électrons rapides dans un matériau de numéro atomique $Z$ élevé (comme le plomb) produit un rayonnement X secondaire (\emph{rayonnement de freinage} ou \emph{bremsstrahlung}), plus pénétrant que les $\beta$ eux-mêmes ! Pour les sources $\beta$, on utilise donc des matériaux légers (plexiglas, aluminium), et l'on ne place du plomb qu'en deuxième couche si nécessaire.
\end{attention}

\begin{popfigure}
\begin{center}
\begin{tikzpicture}[scale=0.95, font=\small]
% Règle 1 : temps
\draw[fill=poporangeL, draw=poporange, rounded corners] (0,0) rectangle (3.6,2.2);
\node[popdark, font=\bfseries] at (1.8,1.85) {1. TEMPS};
\draw[fill=white, draw=poporange] (1.35,0.55) circle (0.45);
\draw[popdark, thick, -{Stealth}] (1.35,0.55) -- (1.35,0.9);
\draw[popdark, thick, -{Stealth}] (1.35,0.55) -- (1.6,0.42);
\node[popdark, align=center] at (1.8,0.15) {exposer le moins\\ longtemps possible};
% Règle 2 : distance
\draw[fill=popgreenL, draw=popgreen, rounded corners] (4.2,0) rectangle (7.8,2.2);
\node[popdark, font=\bfseries] at (6.0,1.85) {2. DISTANCE};
\draw[fill=popgold, draw=popdark] (4.7,0.7) circle (0.18);
\node[popdark] at (4.7,0.35) {source};
\draw[popdark, thick, -{Stealth}] (4.9,0.7) -- (7.1,0.7) node[midway, above]{$d$};
\draw[fill=popblue, draw=popdark] (7.2,0.7) circle (0.14);
\node[popdark] at (6.0,0.15) {$\dot{D}\propto 1/d^{2}$};
% Règle 3 : écran
\draw[fill=poppurpleL, draw=poppurple, rounded corners] (8.4,0) rectangle (12,2.2);
\node[popdark, font=\bfseries] at (10.2,1.85) {3. ÉCRAN};
\draw[fill=popgold, draw=popdark] (8.9,0.7) circle (0.18);
\draw[popmuted, thick, decorate, decoration={coil, amplitude=1.5pt, segment length=4pt}] (9.1,0.7) -- (10.1,0.7);
\draw[fill=popmuted!60, draw=popdark] (10.1,0.25) rectangle (10.5,1.15);
\node[popdark] at (10.3,0.05) {Pb};
\draw[popmuted, thick, decorate, decoration={coil, amplitude=1.5pt, segment length=4pt}] (10.5,0.7) -- (11.0,0.7);
\draw[fill=popblue, draw=popdark] (11.3,0.7) circle (0.14);
\end{tikzpicture}
\end{center}
\legende{Les trois règles fondamentales de la radioprotection : réduire le temps d'exposition, augmenter la distance (loi en $1/d^{2}$), interposer un écran adapté au rayonnement.}
\end{popfigure}

\begin{aretenir}[Les trois principes de la radioprotection]
\begin{enumerate}
\item \cle{Temps} : la dose reçue est proportionnelle à la durée d'exposition.
\item \cle{Distance} : pour une source ponctuelle, le débit de dose décroît en $1/d^{2}$.
\item \cle{Écran} : papier ou peau pour les $\alpha$, aluminium ou plexiglas pour les $\beta$, plomb ou béton pour les photons $\gamma$ et X.
\end{enumerate}
\end{aretenir}

\courssub{Atténuation d'un faisceau de photons : loi exponentielle}

\textbf{Observation.}\par
Plaçons un compteur Geiger--Müller derrière des plaques de plomb d'épaisseur croissante, éclairées par une source $\gamma$ fixe. On constate que chaque fois qu'on ajoute une \emph{même} épaisseur de plomb, le nombre de photons détectés est divisé par un \emph{même} facteur. C'est la signature d'une décroissance exponentielle — exactement comme la décroissance radioactive en fonction du temps.

\textbf{Interprétation microscopique.}\par
Un photon $\gamma$ qui traverse la matière peut, à chaque instant de son trajet, interagir avec un atome (effet photoélectrique, diffusion Compton, production de paires) et disparaître du faisceau. Pour une tranche de matériau d'épaisseur infinitésimale $dx$, le nombre de photons absorbés est proportionnel au nombre de photons présents et à l'épaisseur traversée :
\[ dN = -\mu\, N\, dx, \]
où $\mu$ est une constante positive caractéristique du matériau et de l'énergie des photons. En séparant les variables et en intégrant entre $x=0$ (où $N=N_{0}$) et l'épaisseur $x$ :
\[ \int_{N_{0}}^{N(x)} \frac{dN}{N} = -\mu\int_{0}^{x} dx \quad\Longrightarrow\quad \ln\frac{N(x)}{N_{0}} = -\mu x. \]

\begin{propriete}[Loi d'atténuation d'un faisceau de photons — admise]
Un faisceau homogène de $N_{0}$ photons incidents, traversant un écran d'épaisseur $x$, voit le nombre de photons transmis sans interaction donné par :
\[ N(x) = N_{0}\,e^{-\mu x}, \]
où $\mu$ est le \cle{coefficient d'atténuation linéique} du matériau pour l'énergie considérée. Cette loi expérimentale est admise ; son exploitation numérique est exigible.
\end{propriete}

\begin{definition}[Coefficient d'atténuation linéique]
Le \cle{coefficient d'atténuation linéique} $\mu$ est l'inverse de la distance caractéristique d'atténuation du matériau. Il s'exprime en $\SI{}{m^{-1}}$ ou en $\SI{}{cm^{-1}}$ et dépend :
\begin{itemize}
\item de la \textbf{nature du matériau} (plus le numéro atomique $Z$ et la densité sont élevés, plus $\mu$ est grand : le plomb est un excellent absorbeur) ;
\item de l'\textbf{énergie des photons} (en général, $\mu$ diminue quand l'énergie augmente : les photons très énergétiques sont plus pénétrants).
\end{itemize}
\emph{Vérification dimensionnelle} : l'exposant $-\mu x$ doit être sans dimension, donc $[\mu] = \mathrm{L}^{-1}$ : $\mu$ s'exprime bien en $\SI{}{m^{-1}}$.
\end{definition}

\begin{attention}[Analogie à connaître, confusion à éviter]
La loi $N(x) = N_{0}e^{-\mu x}$ est \emph{formellement identique} à la loi de décroissance radioactive $N(t) = N_{0}e^{-\lambda t}$ : $\mu$ joue le rôle de la constante radioactive $\lambda$, et l'épaisseur $x$ celui du temps $t$. Mais il ne faut \textbf{pas confondre} :
\begin{itemize}
\item la \cle{couche de demi-atténuation} $x_{1/2}$ (en cm ou m) : grandeur \emph{spatiale}, épaisseur qui divise le faisceau par deux ;
\item la \cle{demi-vie} $T$ (en s) : grandeur \emph{temporelle}, durée au bout de laquelle la moitié des noyaux s'est désintégrée.
\end{itemize}
Les formules sont analogues ($x_{1/2} = \ln 2/\mu$ et $T = \ln 2/\lambda$), mais les grandeurs physiques sont de natures différentes. Au Bac, mélanger les deux dans une copie coûte des points !
\end{attention}

\begin{popfigure}
\begin{center}
\begin{tikzpicture}[scale=1.0, font=\small]
% source
\draw[fill=popgold, draw=popdark] (0,0) circle (0.25);
\node[popdark] at (0,-0.55) {source $\gamma$};
% faisceau incident
\foreach \y in {-0.35,-0.12,0.12,0.35}{
  \draw[poporange, thick, -{Stealth}, decorate, decoration={coil, amplitude=1.2pt, segment length=4pt}] (0.3,\y) -- (2.3,\y);}
\node[popdark] at (1.3,0.65) {$N_{0}$ photons incidents};
% écran
\draw[fill=popmuted!55, draw=popdark] (2.3,-0.9) rectangle (4.1,0.9);
\node[popdark] at (3.2,0) {écran};
\draw[popdark, -{Stealth}] (2.3,-1.15) -- (4.1,-1.15) node[midway, below]{$x$};
% photons absorbés
\draw[popmuted, decorate, decoration={coil, amplitude=1.2pt, segment length=4pt}] (2.3,0.35) -- (3.0,0.2);
\draw[popmuted, decorate, decoration={coil, amplitude=1.2pt, segment length=4pt}] (2.3,-0.12) -- (3.4,-0.3);
\draw[popmuted, decorate, decoration={coil, amplitude=1.2pt, segment length=4pt}] (2.3,0.12) -- (2.9,0.45);
% faisceau transmis
\draw[poppurple, thick, -{Stealth}, decorate, decoration={coil, amplitude=1.2pt, segment length=4pt}] (4.1,0.1) -- (6.2,0.1);
\node[popdark] at (5.2,0.6) {$N(x) = N_{0}e^{-\mu x}$};
\node[popdark] at (5.2,-0.35) {photons transmis};
\end{tikzpicture}
\end{center}
\legende{Un faisceau de $N_{0}$ photons traverse un écran d'épaisseur $x$ : une partie des photons interagit dans le matériau, seuls $N(x)$ photons sont transmis sans interaction.}
\end{popfigure}

\courssub{Couche de demi-atténuation (CDA)}

\begin{definition}[Couche de demi-atténuation]
La \cle{couche de demi-atténuation} (CDA), notée $x_{1/2}$, est l'épaisseur de matériau qui divise par deux le nombre de photons transmis :
\[ N(x_{1/2}) = \frac{N_{0}}{2}. \]
\end{definition}

\begin{propriete}[Expression de la CDA — démonstration guidée en classe]
La couche de demi-atténuation est reliée au coefficient d'atténuation linéique par :
\[ x_{1/2} = \frac{\ln 2}{\mu} \approx \frac{0{,}693}{\mu}. \]
\end{propriete}

\begin{experience}[Démonstration guidée : de $N(x_{1/2}) = N_{0}/2$ à $x_{1/2} = \ln 2/\mu$]
Partons de la définition de la CDA et appliquons la loi d'atténuation :
\begin{align*}
N(x_{1/2}) &= \frac{N_{0}}{2} \\
N_{0}\,e^{-\mu x_{1/2}} &= \frac{N_{0}}{2} \\
e^{-\mu x_{1/2}} &= \frac{1}{2} \\
-\mu x_{1/2} &= \ln\frac{1}{2} = -\ln 2 \\
x_{1/2} &= \frac{\ln 2}{\mu}.
\end{align*}
Chaque étape est un calcul de niveau Terminale : simplification par $N_{0}$, passage au logarithme népérien, propriété $\ln(1/2) = -\ln 2$. Refais cette démonstration seul, sans regarder : elle tombe régulièrement au Bac.
\end{experience}

\textbf{Conséquence immédiate : atténuation après $n$ CDA.}\par
Après une épaisseur d'une CDA, il reste $N_{0}/2$ photons ; après deux CDA, la moitié de cette moitié, soit $N_{0}/4$ ; et ainsi de suite. Après $n$ couches de demi-atténuation (épaisseur totale $x = n\,x_{1/2}$) :
\[ N = \frac{N_{0}}{2^{n}}. \]
Cette relation est d'une redoutable efficacité pratique : dix CDA divisent le faisceau par $2^{10} = 1024 \approx 1000$.

\begin{popfigure}
\begin{center}
\begin{tikzpicture}
\begin{axis}[
  width=0.85\linewidth, height=6.5cm,
  xlabel={$x/x_{1/2}$}, ylabel={$N(x)/N_{0}$},
  xmin=0, xmax=3.4, ymin=0, ymax=1.05,
  xtick={0,1,2,3}, xticklabels={$0$,$x_{1/2}$,$2x_{1/2}$,$3x_{1/2}$},
  ytick={0.125,0.25,0.5,1}, yticklabels={$N_{0}/8$,$N_{0}/4$,$N_{0}/2$,$N_{0}$},
  grid=major, grid style={popline!40},
  axis lines=left, tick label style={font=\small}, label style={font=\small}
]
\addplot[popcurve, domain=0:3.4, samples=200]{exp(-0.693*x)};
\addplot[popmuted, dashed] coordinates {(1,0) (1,0.5) (0,0.5)};
\addplot[popmuted, dashed] coordinates {(2,0) (2,0.25) (0,0.25)};
\addplot[popmuted, dashed] coordinates {(3,0) (3,0.125) (0,0.125)};
\addplot[only marks, mark=*, mark size=2pt, popink] coordinates {(1,0.5) (2,0.25) (3,0.125)};
\end{axis}
\end{tikzpicture}
\end{center}
\legende{Atténuation exponentielle d'un faisceau de photons : à chaque couche de demi-atténuation supplémentaire, le nombre de photons transmis est divisé par deux.}
\end{popfigure}

\begin{exemplebox}[Le plomb, champion de la radioprotection]
Pour des photons $\gamma$ d'énergie courante en médecine nucléaire, le plomb a un coefficient d'atténuation $\mu = \SI{1,2}{cm^{-1}}$. Sa couche de demi-atténuation vaut :
\[ x_{1/2} = \frac{\ln 2}{\mu} = \frac{0{,}693}{1{,}2} \approx \SI{0,58}{cm}. \]
Moins de 6 mm de plomb suffisent à diviser le faisceau par deux ! Pour le diviser par 1000, il faut $n = 10$ CDA (car $2^{10} = 1024$), soit une épaisseur :
\[ x = 10\times 0{,}58 \approx \SI{5,8}{cm}. \]
C'est l'ordre de grandeur des protections des enceintes de stockage des sources en médecine nucléaire.
\end{exemplebox}

\courssub{Méthode de dimensionnement d'un écran et applications}

\begin{methode}[Calculer l'épaisseur d'écran nécessaire]
Pour réduire un faisceau d'un facteur $F$ (c'est-à-dire passer de $N_{0}$ à $N = N_{0}/F$) :
\begin{enumerate}
\item \textbf{Identifier les données} : coefficient $\mu$ (ou CDA $x_{1/2}$) du matériau, facteur de réduction $F$ souhaité.
\item \textbf{Choisir la voie de calcul} :
  \begin{itemize}
  \item si $F$ est une puissance de 2 ($F = 2^{n}$), compter les CDA : $x = n\,x_{1/2}$ ;
  \item sinon, utiliser le logarithme népérien.
  \end{itemize}
\item \textbf{Voie logarithme} : écrire $N_{0}/F = N_{0}e^{-\mu x}$, simplifier par $N_{0}$, prendre le logarithme :
  \[ e^{-\mu x} = \frac{1}{F} \quad\Longrightarrow\quad x = \frac{\ln F}{\mu}. \]
\item \textbf{Vérifier} : l'épaisseur obtenue est-elle positive, en unité cohérente avec celle de $\mu$ ($\mu$ en $\SI{}{cm^{-1}}$ donne $x$ en cm) ? Le résultat est-il physiquement raisonnable ?
\end{enumerate}
\end{methode}

\begin{exoresolu}[Dimensionner le mur d'une salle de radiographie]
Dans un hôpital de Garoua, on veut que le mur d'une salle de radiographie réduise le faisceau de rayons X d'un facteur $F = 100$. Le béton utilisé a un coefficient d'atténuation linéique $\mu = \SI{0,15}{cm^{-1}}$ pour ces photons. Quelle épaisseur minimale de béton faut-il prévoir ?
\tcblower
On applique la loi d'atténuation avec $N(x) = N_{0}/100$ :
\begin{align*}
\frac{N_{0}}{100} &= N_{0}\,e^{-\mu x} \\
e^{-\mu x} &= \frac{1}{100} \\
x &= \frac{\ln 100}{\mu} = \frac{4{,}61}{0{,}15} \approx \SI{31}{cm}.
\end{align*}
Il faut donc prévoir environ $\SI{31}{cm}$ de béton. Vérification par les CDA : $x_{1/2} = \ln 2/\mu \approx \SI{4,6}{cm}$, et $100 \approx 2^{6{,}64}$, donc $x \approx 6{,}64\times 4{,}6 \approx \SI{31}{cm}$ : les deux méthodes concordent.
\end{exoresolu}

\textbf{Applications concrètes.}\par
\begin{itemize}
\item \textbf{Tabliers de plomb} : le tablier que le radiologue te fait porter lors d'une radiographie dentaire contient l'équivalent de $\SI{0,25}{mm}$ à $\SI{0,5}{mm}$ de plomb ; il protège les organes sensibles (thyroïde, gonades) des photons diffusés.
\item \textbf{Murs des salles de radiographie et de radiothérapie} : dimensionnés en béton (parfois baryté, enrichi en baryum) à l'aide de la loi $N(x) = N_{0}e^{-\mu x}$, comme dans l'exercice résolu ci-dessus.
\item \textbf{Vitrines et châteaux de plomb des laboratoires de médecine nucléaire} : les solutions radioactives sont préparées derrière des écrans de plomb de plusieurs centimètres, avec des vitres au plomb (verre plombé) pour voir sans être exposé.
\item \textbf{Gammagraphie industrielle} : les sources (\ce{^{192}Ir}, \ce{^{60}Co}) sont stockées dans des conteneurs blindés dont l'épaisseur est calculée pour ramener le débit de dose en surface sous les limites réglementaires.
\end{itemize}

\begin{savaistu}[Le radon, première source d'exposition naturelle]
Au Cameroun comme partout dans le monde, l'essentiel de l'exposition naturelle aux rayonnements ne vient pas des centrales ni des hôpitaux, mais du \cle{radon} (\ce{^{222}Rn}), gaz radioactif issu de la désintégration de l'uranium présent dans les sols granitiques. Il s'accumule dans les maisons mal ventilées. La radioprotection, ici, tient en un mot : \emph{aérer} ! C'est une application directe du premier principe : réduire le temps d'exposition à une concentration élevée.
\end{savaistu}

\begin{exercice}[Tablier de plomb et CDA]
Un tablier de protection a une épaisseur équivalente à $e = \SI{0,50}{mm}$ de plomb. Pour les photons X diffusés en radiographie, le plomb a un coefficient $\mu = \SI{4,0}{mm^{-1}}$.
\begin{enumerate}
\item Calculer la couche de demi-atténuation du plomb pour ces photons.
\item De quel facteur le tablier réduit-il le faisceau ?
\item Quelle épaisseur de plomb faudrait-il pour diviser le faisceau par 16 ?
\end{enumerate}
\end{exercice}

\begin{corrige}[Exercice : tablier de plomb et CDA]
\begin{enumerate}
\item $x_{1/2} = \dfrac{\ln 2}{\mu} = \dfrac{0{,}693}{4{,}0} \approx \SI{0,17}{mm}$.
\item $N/N_{0} = e^{-\mu e} = e^{-4{,}0\times 0{,}50} = e^{-2} \approx 0{,}135$. Le faisceau est divisé par $e^{2} \approx 7{,}4$.
\item Diviser par $16 = 2^{4}$ exige $n = 4$ CDA : $e' = 4\times 0{,}17 \approx \SI{0,69}{mm}$.
\end{enumerate}
\end{corrige}

\begin{aretenir}[L'essentiel de la section]
\begin{itemize}
\item Trois principes de radioprotection : \cle{temps} minimal (dose proportionnelle à $\Delta t$), \cle{distance} maximale (débit de dose en $1/d^{2}$), \cle{écran} adapté (papier pour $\alpha$, aluminium/plexiglas pour $\beta$, plomb/béton pour $\gamma$ et X).
\item Loi d'atténuation des photons : $N(x) = N_{0}e^{-\mu x}$, avec $\mu$ en $\SI{}{m^{-1}}$ ou $\SI{}{cm^{-1}}$, dépendant du matériau et de l'énergie des photons.
\item Couche de demi-atténuation : $x_{1/2} = \dfrac{\ln 2}{\mu} \approx \dfrac{0{,}693}{\mu}$ ; après $n$ CDA, $N = N_{0}/2^{n}$.
\item Pour un facteur de réduction $F$ quelconque : $x = \dfrac{\ln F}{\mu}$.
\item Ne pas confondre $x_{1/2}$ (espace) et demi-vie $T$ (temps) : formules analogues, grandeurs différentes.
\end{itemize}
\end{aretenir}

\newpage

\coursec{Activité d'intégration}

\coursec{Activité d'intégration : dimensionner la protection d'une salle de radiothérapie}

\courssub{Objectifs de l'activité}

À l'issue de cette activité d'intégration, tu seras capable de :

\begin{itemize}
\item écrire et équilibrer l'équation d'une réaction nucléaire provoquée en appliquant les \cle{lois de conservation} du nombre de nucléons $A$ et du nombre de charges $Z$ (lois de Soddy) ;
\item exploiter la relation d'Einstein $E = \Delta m\,c^{2}$ pour relier l'énergie libérée par une transformation nucléaire au défaut de masse, et passer de l'échelle du noyau à l'échelle macroscopique (un gramme de matière) ;
\item utiliser la loi d'atténuation d'un faisceau de photons $N(x) = N_{0}\,e^{-\mu x}$ et la notion de \cle{couche de demi-atténuation} (CDA) pour dimensionner un écran de protection ;
\item rédiger une note de synthèse argumentée mobilisant les trois principes de la \cle{radioprotection} (distance, temps, écran) et la distinction entre \cle{irradiation} et \cle{contamination}.
\end{itemize}

Cette activité mobilise donc l'ensemble des savoir-faire de la leçon dans une situation authentique : le dimensionnement de la protection d'une salle de radiothérapie.

\courssub{Situation}

L'hôpital général de Douala, l'un des plus grands centres hospitaliers du Cameroun, modernise son service de cancérologie. La direction a décidé d'installer une nouvelle salle de \cle{radiothérapie} destinée au traitement des tumeurs par rayonnements. L'appareil choisi utilise une source de \cle{cobalt-60} $\ce{^{60}_{27}Co}$, noyau radioactif émetteur de rayonnements $\gamma$ très pénétrants, produit industriellement par irradiation neutronique du cobalt-59 naturel dans un réacteur nucléaire.

Tu es le(la) physicien(ne) recruté(e) comme conseiller technique par la direction de l'hôpital. Ta mission comporte quatre volets : comprendre la production du cobalt-60, estimer l'énergie mise en jeu, dimensionner le mur de béton qui protégera le personnel et les visiteurs, et rédiger une note d'information sur les règles de radioprotection.

\textbf{Données numériques :}

\begin{center}
\begin{tabularx}{\linewidth}{|l|X|}
\hline
\rowcolor{popblueL} Grandeur & Valeur \\
\hline
Masse molaire du cobalt-60 & $M = \SI{60}{g.mol^{-1}}$ \\
\hline
Nombre d'Avogadro & $N_{A} = \num{6,02e23} \; \si{mol^{-1}}$ \\
\hline
Énergie libérée par désintégration d'un noyau de $\ce{^{60}_{27}Co}$ & $E_{1} = \SI{2,8}{MeV}$ \\
\hline
Conversion & $\SI{1}{MeV} = \num{1,60e-13} \; \si{J}$ \\
\hline
Unité de masse atomique & $\SI{1}{u} = \SI{931,5}{MeV/c^{2}} = \num{1,66e-27} \; \si{kg}$ \\
\hline
Célérité de la lumière dans le vide & $c = \num{3,00e8} \; \si{m.s^{-1}}$ \\
\hline
Coefficient d'atténuation linéique du béton pour les photons $\gamma$ du cobalt-60 & $\mu = \SI{0,20}{cm^{-1}}$ \\
\hline
Valeurs utiles & $\ln 2 \approx 0,693$ ; $\ln 1000 \approx 6,91$ \\
\hline
\end{tabularx}
\end{center}

Le cahier des charges de la direction impose que le mur de béton séparant la salle de traitement du couloir du personnel réduise l'intensité du faisceau de photons à \textbf{un millième} de sa valeur initiale, c'est-à-dire $\dfrac{N}{N_{0}} = \dfrac{1}{1000}$.

\courssub{Consignes}

\textbf{Partie 1 — Production du cobalt-60}

\begin{enumerate}
\item Le cobalt-60 est produit en bombardant des cibles de cobalt-59 $\ce{^{59}_{27}Co}$ par des neutrons $\ce{^{1}_{0}n}$. Écrire l'équation complète de la réaction nucléaire provoquée, sachant qu'un photon $\gamma$ est également émis.
\item Vérifier explicitement que les deux lois de conservation (nombre de nucléons $A$ et nombre de charges $Z$) sont satisfaites. De quel type de réaction nucléaire provoquée s'agit-il ?
\end{enumerate}

\textbf{Partie 2 — Bilan énergétique}

\begin{enumerate}
\setcounter{enumi}{2}
\item La désintégration d'un noyau de cobalt-60 libère $E_{1} = \SI{2,8}{MeV}$. Calculer le défaut de masse $\Delta m$ correspondant, en unités de masse atomique (u) puis en kilogrammes.
\item Calculer le nombre de noyaux contenus dans $m = \SI{1}{g}$ de cobalt-60, puis l'énergie totale $E_{\text{tot}}$ libérée par la désintégration de tous ces noyaux, en joules. Commenter l'ordre de grandeur obtenu.
\end{enumerate}

\textbf{Partie 3 — Dimensionnement du mur de protection}

\begin{enumerate}
\setcounter{enumi}{4}
\item Calculer la couche de demi-atténuation (CDA) du béton pour ces photons, c'est-à-dire l'épaisseur $x_{1/2}$ telle que $N(x_{1/2}) = \dfrac{N_{0}}{2}$.
\item En utilisant directement la loi $N(x) = N_{0}\,e^{-\mu x}$, calculer l'épaisseur $x$ du mur nécessaire pour obtenir $\dfrac{N}{N_{0}} = \dfrac{1}{1000}$.
\item Retrouver ce résultat par la méthode des CDA : combien de couches de demi-atténuation faut-il pour diviser l'intensité par 1000 ? Conclure sur l'épaisseur du mur à construire (arrondir au centimètre supérieur par sécurité).
\end{enumerate}

\textbf{Partie 4 — Note de synthèse}

\begin{enumerate}
\setcounter{enumi}{7}
\item Rédiger une note de synthèse d'une quinzaine de lignes à l'attention du directeur de l'hôpital, rappelant les trois principes de la radioprotection (distance, temps, écran) et expliquant clairement la différence entre irradiation et contamination, avec une consigne concrète pour le personnel soignant.
\end{enumerate}

\courssub{Ressources à mobiliser}

\begin{aretenir}[Boîte à outils de l'activité]
\begin{itemize}
\item \textbf{Lois de conservation (Soddy)} : dans toute réaction nucléaire, la somme des nombres de masse $A$ et la somme des nombres de charge $Z$ se conservent.
\item \textbf{Relation d'Einstein} : $E = \Delta m\,c^{2}$, où $\Delta m$ est le défaut de masse. Avec $\SI{1}{u} = \SI{931,5}{MeV/c^{2}}$, un défaut de masse de 1 u correspond à une énergie de $\SI{931,5}{MeV}$.
\item \textbf{Nombre de noyaux} dans une masse $m$ d'échantillon : $N = \dfrac{m}{M}\,N_{A}$.
\item \textbf{Loi d'atténuation} d'un faisceau de photons à travers un écran d'épaisseur $x$ : $N(x) = N_{0}\,e^{-\mu x}$, où $\mu$ est le coefficient d'atténuation linéique (en $\si{cm^{-1}}$ si $x$ est en cm).
\item \textbf{Couche de demi-atténuation} : $x_{1/2} = \dfrac{\ln 2}{\mu}$. Après $n$ CDA, il reste $\dfrac{N_{0}}{2^{n}}$ photons.
\item \textbf{Radioprotection} : trois principes — augmenter la \cle{distance}, réduire le \cle{temps} d'exposition, interposer un \cle{écran}. L'\cle{irradiation} est une exposition à distance à un rayonnement (elle cesse quand la source est éloignée) ; la \cle{contamination} est la présence de substance radioactive sur ou dans l'organisme (elle persiste tant que la substance n'est pas éliminée).
\end{itemize}
\end{aretenir}

\courssub{Démarche et corrigé commenté}

\begin{exoresolu}[Dimensionnement de la salle de radiothérapie de l'hôpital général de Douala]
Énoncé : reprendre l'ensemble des questions des parties 1 à 4 à partir des données de la situation.
\tcblower
\textbf{Partie 1 — Production du cobalt-60}

\textbf{Question 1.} Le cobalt-59 capture un neutron ; le noyau formé est le cobalt-60, et l'excès d'énergie est évacué par un photon $\gamma$ (qui ne porte ni masse ni charge) :
\[ \ce{^{59}_{27}Co + ^{1}_{0}n -> ^{60}_{27}Co + \gamma} \]

\textbf{Question 2.} Vérification des lois de conservation de Soddy :
\begin{itemize}
\item conservation du nombre de nucléons : $59 + 1 = 60 + 0$ \checkmark
\item conservation du nombre de charges : $27 + 0 = 27 + 0$ \checkmark
\end{itemize}
Il s'agit d'une \cle{transmutation} provoquée (ici une capture neutronique) : un noyau est transformé en un autre noyau sous l'action d'un projectile. Le photon $\gamma$, sans masse ni charge, n'intervient pas dans les bilans de $A$ et $Z$.

\textbf{Partie 2 — Bilan énergétique}

\textbf{Question 3.} D'après la relation d'Einstein, l'énergie libérée correspond au défaut de masse $\Delta m = \dfrac{E_{1}}{c^{2}}$. En utilisant l'équivalence $\SI{1}{u} \leftrightarrow \SI{931,5}{MeV}$ :
\[ \Delta m = \frac{\SI{2,8}{MeV}}{\SI{931,5}{MeV/u}} \approx \num{3,0e-3} \; \si{u} \]
Conversion en kilogrammes :
\[ \Delta m = \num{3,0e-3} \times \num{1,66e-27} \approx \num{5,0e-30} \; \si{kg} \]
Vérification par le calcul direct en unités SI : $E_{1} = 2,8 \times \num{1,60e-13} = \num{4,48e-13} \; \si{J}$, d'où
\[ \Delta m = \frac{E_{1}}{c^{2}} = \frac{\num{4,48e-13}}{(\num{3,00e8})^{2}} = \frac{\num{4,48e-13}}{\num{9,0e16}} \approx \num{5,0e-30} \; \si{kg} \]
Les deux méthodes concordent. L'analyse dimensionnelle confirme le résultat : $\dfrac{[\text{J}]}{[\text{m}^{2}\text{.s}^{-2}]} = \dfrac{\si{kg.m^{2}.s^{-2}}}{\si{m^{2}.s^{-2}}} = \si{kg}$.

\textbf{Question 4.} Nombre de noyaux dans $\SI{1}{g}$ de cobalt-60 :
\[ N = \frac{m}{M}\,N_{A} = \frac{1}{60} \times \num{6,02e23} \approx \num{1,0e22} \; \text{noyaux} \]
Énergie totale libérée :
\[ E_{\text{tot}} = N \times E_{1} = \num{1,0e22} \times \num{4,48e-13} \approx \num{4,5e9} \; \si{J} \]
Soit environ \textbf{4,5 milliards de joules} pour un seul gramme : c'est l'équivalent de l'énergie électrique consommée par un quartier de Douala pendant plusieurs jours. Cet ordre de grandeur illustre la densité énergétique extraordinaire des transformations nucléaires, environ un million de fois supérieure à celle des réactions chimiques.

\textbf{Partie 3 — Dimensionnement du mur}

\textbf{Question 5.} Couche de demi-atténuation :
\[ x_{1/2} = \frac{\ln 2}{\mu} = \frac{0,693}{0,20} \approx \SI{3,5}{cm} \]
Tous les $\SI{3,5}{cm}$, le béton divise le nombre de photons transmis par deux.

\textbf{Question 6.} On impose $\dfrac{N(x)}{N_{0}} = e^{-\mu x} = \dfrac{1}{1000}$. En prenant le logarithme népérien des deux membres :
\[ -\mu x = \ln\left(\frac{1}{1000}\right) = -\ln 1000 \quad\Longrightarrow\quad x = \frac{\ln 1000}{\mu} = \frac{6,91}{0,20} \approx \SI{34,6}{cm} \]

\textbf{Question 7.} Méthode des CDA : après $n$ couches de demi-atténuation, $\dfrac{N}{N_{0}} = \dfrac{1}{2^{n}}$. On cherche le plus petit entier $n$ tel que $2^{n} \geq 1000$. Or $2^{9} = 512$ (insuffisant) et $2^{10} = 1024$ (suffisant) : il faut donc $n = 10$ CDA.
\[ x = 10 \times x_{1/2} = 10 \times 3,5 = \SI{35}{cm} \]
Les deux méthodes concordent (34,6 cm $\approx$ 35 cm, l'écart vient des arrondis sur $\ln 2$ et sur $x_{1/2}$). \textbf{Conclusion} : par sécurité, on retiendra un mur de béton d'au moins \textbf{35 cm} d'épaisseur, arrondi au centimètre supérieur, conformément au cahier des charges.

\textbf{Partie 4 — Note de synthèse (éléments attendus)}

\emph{À l'attention du Directeur de l'hôpital général de Douala — Objet : règles de radioprotection autour de la salle de cobaltthérapie.}

La radioprotection du personnel repose sur trois principes complémentaires. \textbf{La distance} : l'exposition décroît fortement quand on s'éloigne de la source ; le personnel ne pénètre dans la salle que lorsque la source est en position fermée. \textbf{Le temps} : toute durée de présence inutile à proximité de la source doit être supprimée ; les interventions sont préparées et chronométrées. \textbf{L'écran} : le mur de béton de 35 cm calculé ci-dessus réduit le faisceau à un millième de son intensité ; les portes plombées complètent ce dispositif.

Il faut enfin distinguer \textbf{irradiation} et \textbf{contamination}. L'irradiation est une exposition au rayonnement émis par une source extérieure : elle cesse dès que l'on s'éloigne de la source ou que celle-ci est obturée ; le personnel irradié n'est pas lui-même radioactif. La contamination est le dépôt de substance radioactive sur la peau, les vêtements, ou à l'intérieur du corps (inhalation, ingestion) : elle persiste tant que la substance n'est pas éliminée. Consigne pratique : port obligatoire du dosimètre, des gants et de la blouse, lavage des mains systématique, et interdiction de manger ou boire dans la zone contrôlée.
\end{exoresolu}

\begin{attention}[Pièges fréquents dans ce type d'exercice]
\begin{itemize}
\item Ne pas oublier que le photon $\gamma$ n'a ni $A$ ni $Z$ : il ne modifie pas les bilans de conservation.
\item Homogénéité des unités : avec $\mu$ en $\si{cm^{-1}}$, l'épaisseur $x$ s'obtient en \textbf{cm}. Convertir $\mu$ en $\si{m^{-1}}$ si l'on veut $x$ en mètres.
\item Ne pas confondre $x_{1/2} = \dfrac{\ln 2}{\mu}$ (CDA, en cm) avec la période radioactive $t_{1/2} = \dfrac{\ln 2}{\lambda}$ (en s) : les formules se ressemblent mais concernent des phénomènes différents — l'une décrit l'atténuation \emph{dans l'espace}, l'autre la décroissance \emph{dans le temps}.
\item Dans la méthode des CDA, arrondir toujours le nombre de couches à l'entier \textbf{supérieur} : la protection doit être au moins égale à la valeur exigée.
\item Dans $E = \Delta m\,c^{2}$, ne pas oublier d'élever $c$ au carré et de convertir les MeV en joules avant tout calcul en unités SI.
\end{itemize}
\end{attention}

\courssub{Bilan de l'activité}

Cette activité t'a permis de mettre en évidence, dans un contexte hospitalier camerounais réaliste, les points essentiels de la leçon :

\begin{itemize}
\item les \textbf{lois de conservation} de Soddy suffisent à écrire et vérifier l'équation de n'importe quelle réaction nucléaire provoquée, comme la production du cobalt-60 par capture neutronique ;
\item la relation $E = \Delta m\,c^{2}$ établit le pont entre l'infiniment petit (un défaut de masse de $\num{5e-30}$ kg par noyau) et le monde macroscopique (plus de 4 milliards de joules pour un gramme), ce qui explique à la fois l'intérêt médical et la dangerosité des sources radioactives ;
\item la loi $N(x) = N_{0}\,e^{-\mu x}$ et la notion de CDA fournissent un outil de calcul concret pour \textbf{dimensionner un écran} : ici, 35 cm de béton divisent le rayonnement par 1000 ;
\item la radioprotection n'est pas seulement un calcul d'épaisseur : elle repose sur la combinaison des trois principes \textbf{distance -- temps -- écran} et sur la compréhension de la différence entre irradiation et contamination, indispensable à la sécurité du personnel soignant.
\end{itemize}

Tu es désormais capable de traiter au Baccalauréat tout exercice combinant équation de réaction nucléaire, bilan énergétique et atténuation d'un faisceau de photons, et de justifier scientifiquement des mesures de protection dans la vie professionnelle.

\newpage

\coursec{Méthodes et exercices résolus}

\courssub{Méthode 1 : Écrire l'équation d'une réaction nucléaire provoquée}

Contrairement à la radioactivité, qui est \emph{spontanée}, une réaction nucléaire \cle{provoquée} résulte du choc entre un projectile (neutron, proton, particule $\alpha$\ldots) et un noyau cible. Dans tous les cas, l'outil de travail est le même : les deux lois de conservation de Soddy.

\begin{methode}[Équilibrer une réaction nucléaire avec les lois de Soddy]
\begin{enumerate}
\item \cle{Identifier} les noyaux et particules en présence : projectile, cible, produits. Écrire chacun sous la forme $\ce{^{A}_{Z}X}$ (A = nombre de nucléons, Z = nombre de charges).
\item \cle{Écrire les deux lois de conservation} (lois de Soddy) :
\begin{itemize}
\item conservation du nombre de nucléons : $\sum A_{\text{réactifs}} = \sum A_{\text{produits}}$ ;
\item conservation de la charge électrique : $\sum Z_{\text{réactifs}} = \sum Z_{\text{produits}}$.
\end{itemize}
\item \cle{Résoudre} le petit système à deux inconnues (A et Z du noyau ou de la particule inconnue).
\item \cle{Identifier l'élément} à partir de Z grâce au tableau périodique (Z définit l'élément, jamais A).
\item \cle{Vérifier} : réécrire l'équation complète et contrôler mentalement les deux sommes. Rappel des particules usuelles : neutron $\ce{^{1}_{0}n}$, proton $\ce{^{1}_{1}p}$, particule $\alpha$ $\ce{^{4}_{2}He}$, électron $\ce{^{0}_{-1}e}$, positon $\ce{^{0}_{1}e}$, photon $\gamma$ (A = 0, Z = 0).
\end{enumerate}
\end{methode}

\begin{exoresolu}[La première transmutation artificielle (Rutherford, 1919)]
En bombardant de l'azote $\ce{^{14}_{7}N}$ avec des particules $\alpha$, Rutherford obtient un noyau d'oxygène $\ce{^{17}_{8}O}$ et une particule inconnue $\ce{^{A}_{Z}X}$.
\begin{enumerate}
\item Écrire l'équation de la réaction et déterminer A et Z.
\item Identifier la particule émise.
\end{enumerate}
\tcblower
\textbf{1.} L'équation s'écrit :
\[ \ce{^{14}_{7}N} + \ce{^{4}_{2}He} \longrightarrow \ce{^{17}_{8}O} + \ce{^{A}_{Z}X} \]
Appliquons les lois de Soddy :
\begin{itemize}
\item conservation des nucléons : $14 + 4 = 17 + A \Rightarrow A = 1$ ;
\item conservation de la charge : $7 + 2 = 8 + Z \Rightarrow Z = 1$.
\end{itemize}
\textbf{2.} La particule $\ce{^{1}_{1}X}$ est un \cle{proton} $\ce{^{1}_{1}p}$ (c'est d'ailleurs ainsi que Rutherford découvrit le proton). L'équation complète :
\[ \ce{^{14}_{7}N} + \ce{^{4}_{2}He} \longrightarrow \ce{^{17}_{8}O} + \ce{^{1}_{1}p} \]
Vérification : nucléons $18 = 18$ ; charges $9 = 9$. C'est bien une \cle{transmutation} : un élément (azote) est transformé en un autre (oxygène).
\end{exoresolu}

\courssub{Méthode 2 : Calculer l'énergie libérée par une réaction nucléaire}

Pourquoi ces réactions libèrent-elles de l'énergie ? Parce que la masse totale des produits est \emph{légèrement inférieure} à celle des réactifs : le « défaut de masse » $\Delta m$ s'est converti en énergie selon la célèbre relation d'Einstein.

\begin{methode}[Exploiter $E_{\text{libérée}} = \Delta m \cdot c^{2}$]
\begin{enumerate}
\item \cle{Calculer le défaut de masse} : $\Delta m = \sum m_{\text{réactifs}} - \sum m_{\text{produits}}$. Attention au sens : pour une réaction qui libère de l'énergie, $\Delta m > 0$ (la masse diminue).
\item \cle{Appliquer la relation d'Einstein} : $E = \Delta m \cdot c^{2}$ avec $\Delta m$ en kilogrammes et $c = \num{3,00}\times 10^{8}\ \mathrm{m\,s^{-1}}$ ; E s'obtient alors en joules.
\item \cle{Convertir si besoin} : si les masses sont données en unités de masse atomique (u), utiliser $1\ \mathrm{u} \cdot c^{2} = \num{931,5}\ \mathrm{MeV}$, ce qui donne directement E en MeV. Rappels : $1\ \mathrm{u} = \num{1,66}\times 10^{-27}\ \mathrm{kg}$ et $1\ \mathrm{MeV} = \num{1,60}\times 10^{-13}\ \mathrm{J}$.
\item \cle{Conclure par une phrase} : préciser le signe (énergie libérée, donc cédée au milieu extérieur) et donner le résultat avec unité et chiffres significatifs cohérents.
\end{enumerate}
\textbf{Réflexe analyse dimensionnelle} : $[\Delta m \cdot c^{2}] = \mathrm{kg}\cdot\mathrm{m^{2}\,s^{-2}} = \mathrm{J}$. La relation est homogène.
\end{methode}

\begin{exoresolu}[Énergie libérée par la fission de l'uranium 235]
On considère la fission d'un noyau d'uranium 235 :
\[ \ce{^{235}_{92}U} + \ce{^{1}_{0}n} \longrightarrow \ce{^{94}_{38}Sr} + \ce{^{140}_{54}Xe} + 2\,\ce{^{1}_{0}n} \]
Masses : $m(\ce{^{235}U}) = \num{234,9934}\ \mathrm{u}$ ; $m(\ce{^{94}Sr}) = \num{93,8945}\ \mathrm{u}$ ; $m(\ce{^{140}Xe}) = \num{139,8920}\ \mathrm{u}$ ; $m_n = \num{1,0087}\ \mathrm{u}$.
\begin{enumerate}
\item Vérifier que l'équation est équilibrée.
\item Calculer l'énergie libérée par la fission d'un noyau, en MeV puis en joules.
\item Calculer l'énergie libérée par la fission de 1,0 g d'uranium 235. On donne $N_A = \num{6,02}\times 10^{23}\ \mathrm{mol^{-1}}$ et $M(\ce{^{235}U}) = 235\ \mathrm{g\,mol^{-1}}$.
\end{enumerate}
\tcblower
\textbf{1.} Nucléons : $235 + 1 = 236$ et $94 + 140 + 2\times 1 = 236$. Charges : $92 + 0 = 92$ et $38 + 54 + 0 = 92$. L'équation est équilibrée.

\textbf{2.} Défaut de masse :
\begin{align*}
\Delta m &= \left[m(\ce{^{235}U}) + m_n\right] - \left[m(\ce{^{94}Sr}) + m(\ce{^{140}Xe}) + 2\,m_n\right] \\
&= (234{,}9934 + 1{,}0087) - (93{,}8945 + 139{,}8920 + 2\times 1{,}0087) \\
&= 236{,}0021 - 235{,}8039 = 0{,}1982\ \mathrm{u}
\end{align*}
Énergie libérée :
\[ E = 0{,}1982 \times 931{,}5 \approx 185\ \mathrm{MeV} \]
Conversion : $E = 185 \times 1{,}60\times 10^{-13} \approx \num{2,95}\times 10^{-11}\ \mathrm{J}$ par noyau.

\textbf{3.} Nombre de noyaux dans 1,0 g :
\[ N = \frac{m}{M}\,N_A = \frac{1{,}0}{235}\times 6{,}02\times 10^{23} \approx \num{2,56}\times 10^{21}\ \text{noyaux} \]
Énergie totale :
\[ E_{\text{tot}} = N \times E = 2{,}56\times 10^{21} \times 2{,}95\times 10^{-11} \approx \num{7,6}\times 10^{10}\ \mathrm{J} \]
Soit environ 76 GJ pour un seul gramme : c'est l'équivalent de l'énergie produite par la combustion d'environ 2 tonnes de pétrole. C'est ce qui explique l'intérêt énergétique des centrales à fission et la densité énergétique exceptionnelle du combustible nucléaire.
\end{exoresolu}

\courssub{Méthode 3 : Distinguer fission et fusion, analyser une réaction en chaîne}

\begin{methode}[Reconnaître et exploiter fission, fusion et réaction en chaîne]
\begin{enumerate}
\item \cle{Fission} : un noyau lourd ($A > 200$, ex. $\ce{^{235}U}$, $\ce{^{239}Pu}$) frappé par un neutron se scinde en deux noyaux plus légers + 2 ou 3 neutrons + énergie. Réaction \emph{provoquée}, contrôlable (centrale) ou non (bombe A).
\item \cle{Fusion} : deux noyaux légers (isotopes de l'hydrogène : deutérium $\ce{^{2}_{1}H}$, tritium $\ce{^{3}_{1}H}$) s'unissent pour former un noyau plus lourd + énergie. Nécessite des températures de l'ordre de $10^{8}$ K pour vaincre la répulsion électrique entre noyaux (étoiles, bombe H, projet ITER).
\item \cle{Réaction en chaîne} : compter les neutrons. Si chaque fission émet en moyenne plus d'un neutron susceptible de provoquer une nouvelle fission, la réaction s'emballe (régime \emph{critique} si exactement 1 neutron utile par fission, \emph{supercritique} si plus, \emph{sous-critique} si moins). Dans un réacteur, des barres de contrôle absorbent l'excès de neutrons pour maintenir le régime critique.
\item \cle{Comparer les énergies} : toujours raisonner \emph{par unité de masse} de « combustible » ; la fusion libère davantage d'énergie par kilogramme que la fission.
\end{enumerate}
\end{methode}

\begin{exoresolu}[Fusion dans le Soleil et projet ITER]
Dans les étoiles comme le Soleil, le bilan global du cycle proton-proton peut s'écrire :
\[ 4\,\ce{^{1}_{1}H} \longrightarrow \ce{^{4}_{2}He} + 2\,\ce{^{0}_{1}e} + 2\,\gamma \]
Masses : $m(\ce{^{1}_{1}H}) = \num{1,0073}\ \mathrm{u}$ ; $m(\ce{^{4}_{2}He}) = \num{4,0015}\ \mathrm{u}$ ; $m(\ce{^{0}_{1}e}) = \num{0,00055}\ \mathrm{u}$.
\begin{enumerate}
\item Vérifier les lois de conservation.
\item Calculer l'énergie libérée par cette réaction en MeV.
\item Le Soleil rayonne une puissance $P \approx \num{3,8}\times 10^{26}\ \mathrm{W}$. Estimer la masse d'hydrogène convertie en hélium chaque seconde.
\end{enumerate}
\tcblower
\textbf{1.} Nucléons : $4\times 1 = 4$ et $4 + 0 + 0 = 4$. Charges : $4\times 1 = 4$ et $2 + 2\times 1 + 0 = 4$. Les lois de Soddy sont vérifiées.

\textbf{2.} Défaut de masse :
\begin{align*}
\Delta m &= 4\,m(\ce{^{1}_{1}H}) - \left[m(\ce{^{4}_{2}He}) + 2\,m(\ce{^{0}_{1}e})\right] \\
&= 4\times 1{,}0073 - (4{,}0015 + 2\times 0{,}00055) \\
&= 4{,}0292 - 4{,}0026 = 0{,}0266\ \mathrm{u}
\end{align*}
\[ E = 0{,}0266 \times 931{,}5 \approx 24{,}8\ \mathrm{MeV} \approx \num{4,0}\times 10^{-12}\ \mathrm{J} \]

\textbf{3.} Masse de 4 noyaux d'hydrogène consommés par réaction :
\[ m_{\text{réac}} = 4 \times 1{,}0073 \times 1{,}66\times 10^{-27} \approx \num{6,69}\times 10^{-27}\ \mathrm{kg} \]
Nombre de réactions par seconde :
\[ n = \frac{P}{E} = \frac{3{,}8\times 10^{26}}{4{,}0\times 10^{-12}} \approx \num{9,5}\times 10^{37}\ \text{réactions/s} \]
Masse d'hydrogène convertie par seconde :
\[ \dot{m} = n \times m_{\text{réac}} \approx 9{,}5\times 10^{37} \times 6{,}69\times 10^{-27} \approx \num{6,4}\times 10^{11}\ \mathrm{kg\,s^{-1}} \]
Le Soleil « brûle » donc environ 640 millions de tonnes d'hydrogène par seconde. Reproduire la fusion sur Terre est l'objectif du projet international ITER : une énergie abondante, sans réaction en chaîne incontrôlable ni déchets de haute activité à vie longue.
\end{exoresolu}

\begin{savaistu}[Et le Cameroun dans tout cela ?]
Le Cameroun ne possède pas de centrale nucléaire : son électricité provient surtout de l'hydroélectricité (barrages d'Edéa, de Song Loulou, de Lom Pangar et bientôt Nachtigal). Mais la physique nucléaire y est bien présente : radiothérapie et imagerie médicale dans les hôpitaux de Yaoundé et Douala (sources de $\ce{^{60}Co}$), gammagraphie industrielle pour contrôler les soudures des oléoducs, et irradiation des denrées pour la conservation. D'où l'importance des règles de radioprotection étudiées ci-dessous.
\end{savaistu}

\courssub{Méthode 4 : Distinguer irradiation et contamination, citer les effets biologiques}

\begin{methode}[Irradiation, contamination, effets sur l'organisme]
\begin{enumerate}
\item \cle{Irradiation} : exposition à des rayonnements émis par une source \emph{extérieure} au corps. Elle cesse dès qu'on s'éloigne de la source ou qu'on interpose un écran. Le corps ne devient pas radioactif.
\item \cle{Contamination} : présence de substances radioactives \emph{sur} la peau (externe) ou \emph{dans} l'organisme (interne : inhalation, ingestion, blessure). Elle persiste tant que la substance n'est pas éliminée.
\item \cle{Effets biologiques} : ionisation des molécules d'eau et de l'ADN $\Rightarrow$ lésions cellulaires. Effets \emph{déterministes} (brûlures, syndrome aigu d'irradiation, avec seuil de dose) et effets \emph{stochastiques} (cancers, mutations, sans seuil, probabilité croissant avec la dose).
\item \cle{Dosimétrie} : dose absorbée $D = E/m$ en gray (Gy), avec E l'énergie reçue et m la masse de tissu exposé ; dose équivalente $H = D \times w_R$ en sievert (Sv), où $w_R$ est le facteur de pondération du rayonnement ($w_R = 1$ pour $\gamma$ et $\beta$ ; $w_R = 20$ pour $\alpha$, très ionisant).
\end{enumerate}
\end{methode}

\begin{exoresolu}[Accident dans un service de radiothérapie]
Lors d'une maintenance dans un hôpital, un technicien reste 30 min près d'une source de $\ce{^{60}Co}$ mal protégée. Il reçoit une énergie de rayonnement $\gamma$ de $E = \num{0,42}\ \mathrm{J}$ répartie sur une masse corporelle $m = 70\ \mathrm{kg}$.
\begin{enumerate}
\item S'agit-il d'une irradiation ou d'une contamination ? Justifier.
\item Calculer la dose absorbée puis la dose équivalente reçue.
\item Commenter en sachant que la dose annuelle due à la radioactivité naturelle est d'environ 2 mSv et qu'une dose aiguë de 1 Sv provoque des troubles sanguins.
\end{enumerate}
\tcblower
\textbf{1.} Il s'agit d'une \cle{irradiation} : la source est extérieure au corps, aucun radionucléide n'a été déposé sur ou dans l'organisme du technicien. L'exposition cesse dès qu'il s'éloigne de la source.

\textbf{2.} Dose absorbée :
\[ D = \frac{E}{m} = \frac{0{,}42}{70} = \num{6,0}\times 10^{-3}\ \mathrm{Gy} = 6{,}0\ \mathrm{mGy} \]
Pour des photons $\gamma$, $w_R = 1$, donc :
\[ H = D \times w_R = 6{,}0\ \mathrm{mSv} \]

\textbf{3.} La dose reçue (6 mSv) représente environ 3 ans d'irradiation naturelle : elle est notable mais reste très inférieure au seuil de 1 Sv des effets déterministes. Seul un léger surcroît de risque stochastique (cancer) est à considérer. Cet exemple illustre l'importance des trois règles de radioprotection : réduire le \cle{temps} d'exposition, augmenter la \cle{distance}, interposer un \cle{écran}.
\end{exoresolu}

\courssub{Méthode 5 : Exploiter la loi d'atténuation d'un faisceau de photons}

Lorsqu'un faisceau de photons traverse un écran de matière, chaque photon a une certaine probabilité d'être absorbé ou diffusé par unité d'épaisseur traversée. Il en résulte une décroissance \emph{exponentielle} du nombre de photons transmis, mathématiquement analogue à la décroissance radioactive.

\begin{methode}[Utiliser $N(x) = N_0\,e^{-\mu x}$ et la couche de demi-atténuation]
\begin{enumerate}
\item \cle{Identifier les données} : $N_0$ (nombre de photons incidents), $\mu$ (coefficient d'atténuation linéique, en $\mathrm{m^{-1}}$ ou $\mathrm{cm^{-1}}$), x (épaisseur traversée, dans l'unité cohérente avec $\mu$).
\item \cle{Appliquer la loi} : $N(x) = N_0\,e^{-\mu x}$. Pour trouver x connaissant $N/N_0$, prendre le logarithme népérien : $x = \dfrac{1}{\mu}\ln\!\left(\dfrac{N_0}{N}\right)$.
\item \cle{Couche de demi-atténuation (CDA)} : épaisseur $x_{1/2}$ qui divise le nombre de photons par 2 :
\[ x_{1/2} = \frac{\ln 2}{\mu} \approx \frac{0{,}693}{\mu} \]
\item \cle{Raisonner en couches} : après n couches de demi-atténuation, $N = N_0 / 2^{n}$. Très utile pour les questions rapides du type « quelle épaisseur pour réduire à 1/8 ? » (réponse : 3 CDA).
\end{enumerate}
\end{methode}

\begin{exoresolu}[Blindage d'une source de cobalt 60]
Les photons $\gamma$ émis par le $\ce{^{60}Co}$ traversent un mur de plomb. Le coefficient d'atténuation linéique du plomb pour ces photons est $\mu = \num{66}\ \mathrm{m^{-1}}$.
\begin{enumerate}
\item Calculer la couche de demi-atténuation du plomb pour ces photons.
\item Quelle fraction du faisceau subsiste après 5,0 cm de plomb ?
\item Quelle épaisseur de plomb faut-il pour ne laisser passer qu'un millième du faisceau incident ?
\end{enumerate}
\tcblower
\textbf{1.} Couche de demi-atténuation :
\[ x_{1/2} = \frac{\ln 2}{\mu} = \frac{0{,}693}{66} \approx \num{1,05}\times 10^{-2}\ \mathrm{m} \approx 1{,}1\ \mathrm{cm} \]

\textbf{2.} Avec $x = 5{,}0\ \mathrm{cm} = \num{5,0}\times 10^{-2}\ \mathrm{m}$ :
\[ \frac{N}{N_0} = e^{-\mu x} = e^{-66 \times 5{,}0\times 10^{-2}} = e^{-3{,}30} \approx 0{,}037 \]
Il subsiste environ 3,7 \% du faisceau incident (environ 1/27).

\textbf{3.} On veut $N/N_0 = 10^{-3}$ :
\[ x = \frac{1}{\mu}\ln\!\left(\frac{N_0}{N}\right) = \frac{\ln(1000)}{66} = \frac{6{,}91}{66} \approx \num{1,05}\times 10^{-1}\ \mathrm{m} \approx 10{,}5\ \mathrm{cm} \]
Vérification par les couches : $10^{-3} \approx 1/2^{10}$, il faut donc environ 10 CDA, soit $10 \times 1{,}05 \approx 10{,}5\ \mathrm{cm}$. Les deux méthodes concordent : une dizaine de centimètres de plomb suffisent à ramener le rayonnement à un millième de sa valeur initiale, ce qui dimensionne les enceintes blindées des sources médicales.
\end{exoresolu}

\courssub{Méthode 6 : Bilan complet d'une réaction de fission en centrale}

\begin{methode}[Enchaîner équation, énergie et exploitation industrielle]
\begin{enumerate}
\item \cle{Équilibrer} l'équation de fission (souvent un neutron incident, deux fragments, k neutrons émis : trouver k par conservation de A).
\item \cle{Calculer} l'énergie par noyau avec $\Delta m\,c^{2}$ (en MeV).
\item \cle{Passer à l'échelle macroscopique} : nombre de noyaux $N = \dfrac{m}{M}N_A$, puis énergie totale, puis puissance ou durée de fonctionnement $t = E/P$.
\item \cle{Conclure} en comparant avec une source d'énergie classique (hydroélectricité des barrages d'Edéa ou de Nachtigal, thermique à flamme) pour donner du sens au résultat.
\end{enumerate}
\end{methode}

\begin{exoresolu}[Alimenter une ville avec la fission]
Dans un réacteur, la fission de l'uranium 235 libère en moyenne $E_1 = 200\ \mathrm{MeV}$ par noyau. Une ville consomme une puissance électrique moyenne $P = 100\ \mathrm{MW}$. Le rendement de conversion de la centrale est $\eta = 33\ \%$.
\begin{enumerate}
\item Exprimer puis calculer l'énergie thermique nécessaire par jour de fonctionnement.
\item En déduire la masse d'uranium 235 consommée par jour. On donne $M = 235\ \mathrm{g\,mol^{-1}}$, $N_A = \num{6,02}\times 10^{23}\ \mathrm{mol^{-1}}$, $1\ \mathrm{MeV} = \num{1,60}\times 10^{-13}\ \mathrm{J}$.
\end{enumerate}
\tcblower
\textbf{1.} Énergie électrique quotidienne :
\[ E_{\text{él}} = P \times t = 100\times 10^{6} \times 24\times 3600 = \num{8,64}\times 10^{12}\ \mathrm{J} \]
Avec un rendement $\eta = E_{\text{él}}/E_{\text{th}}$ :
\[ E_{\text{th}} = \frac{E_{\text{él}}}{\eta} = \frac{8{,}64\times 10^{12}}{0{,}33} \approx \num{2,62}\times 10^{13}\ \mathrm{J} \]

\textbf{2.} Énergie par noyau en joules :
\[ E_1 = 200 \times 1{,}60\times 10^{-13} = \num{3,20}\times 10^{-11}\ \mathrm{J} \]
Nombre de noyaux fissionnés :
\[ N = \frac{E_{\text{th}}}{E_1} = \frac{2{,}62\times 10^{13}}{3{,}20\times 10^{-11}} \approx \num{8,2}\times 10^{23}\ \text{noyaux} \]
Masse correspondante :
\[ m = \frac{N}{N_A}\times M = \frac{8{,}2\times 10^{23}}{6{,}02\times 10^{23}}\times 235 \approx 320\ \mathrm{g} \]
Il suffit donc d'environ 320 g d'uranium 235 par jour pour alimenter une ville de 100 MW, là où une centrale thermique classique brûlerait des centaines de tonnes de combustible fossile. C'est toute la force — et toute la responsabilité — de l'énergie nucléaire.
\end{exoresolu}

\begin{attention}[Les 5 erreurs qui coûtent des points au Bac]
\begin{enumerate}
\item \cle{Confondre A et Z} dans les lois de conservation : A (nucléons) se conserve en haut, Z (charges) en bas. Un neutron $\ce{^{1}_{0}n}$ compte pour 1 dans A mais pour 0 dans Z : l'oublier fausse tout le bilan.
\item \cle{Se tromper de sens dans $\Delta m$} : pour une énergie \emph{libérée}, $\Delta m = m_{\text{réactifs}} - m_{\text{produits}} > 0$. Un défaut de masse négatif avec une énergie « libérée » positive est une incohérence qui saute aux yeux du correcteur.
\item \cle{Oublier les conversions} : masses en u $\rightarrow$ utiliser $1\ \mathrm{u}\cdot c^{2} = 931{,}5\ \mathrm{MeV}$, ou convertir en kg avant d'appliquer $E = \Delta m c^{2}$ (sinon le joule n'est pas obtenu). Ne jamais mélanger u et kg dans le même calcul.
\item \cle{Confondre irradiation et contamination} : l'irradiation vient d'une source externe et cesse avec l'éloignement ; la contamination implique des radionucléides sur ou dans le corps et persiste. Citer aussi les trois moyens de protection (distance, temps, écran) sans en oublier.
\item \cle{Mal manipuler $N = N_0 e^{-\mu x}$} : unités incohérentes entre $\mu$ (en $\mathrm{cm^{-1}}$) et x (en m), oubli du logarithme népérien pour isoler x, ou confusion entre CDA ($x_{1/2} = \ln 2 / \mu$) et $1/\mu$. Vérifier systématiquement que $\mu x$ est sans dimension.
\end{enumerate}
\end{attention}

\newpage

\coursec{Exercices}

\courssub{Je vérifie mes connaissances}

\begin{exercice}[QCM : réactions nucléaires]
Pour chaque question, choisir la (ou les) bonne(s) réponse(s) en justifiant brièvement.
\begin{enumerate}
\item Dans la réaction $\ce{^{235}_{92}U + ^{1}_{0}n -> ^{90}_{38}Sr + ^{143}_{54}Xe + 3^{1}_{0}n}$, le noyau $\ce{^{143}_{54}Xe}$ est identifié grâce à :
\begin{enumerate}
\item la conservation du nombre de masse seulement ;
\item la conservation du nombre de charge seulement ;
\item les deux lois de conservation de Soddy.
\end{enumerate}
\item La fusion nucléaire :
\begin{enumerate}
\item concerne des noyaux lourds ;
\item concerne des noyaux légers ;
\item est la source d'énergie des étoiles ;
\item nécessite des températures de plusieurs millions de degrés.
\end{enumerate}
\item L'unité SI de la dose absorbée est :
\begin{enumerate}
\item le becquerel (Bq) ; \item le gray (Gy) ; \item le sievert (Sv).
\end{enumerate}
\item Si la couche de demi-atténuation d'un matériau vaut $x_{1/2} = \SI{2}{cm}$, l'épaisseur nécessaire pour ne transmettre que $12,5\,\%$ des photons incidents est :
\begin{enumerate}
\item $\SI{4}{cm}$ ; \item $\SI{6}{cm}$ ; \item $\SI{8}{cm}$.
\end{enumerate}
\end{enumerate}
\end{exercice}

\begin{exercice}[Vrai ou faux ?]
Répondre par vrai ou faux en justifiant chaque réponse.
\begin{enumerate}
\item Dans toute réaction nucléaire, le nombre total de nucléons et la charge électrique totale se conservent.
\item La fission est une réaction nucléaire spontanée qui ne peut pas être provoquée.
\item La masse des produits d'une fission est supérieure à la masse des noyaux initiaux.
\item Une irradiation cesse dès que la source radioactive est éloignée, contrairement à une contamination.
\item Le plomb atténue un faisceau de photons plus efficacement que le béton de même épaisseur si son coefficient d'atténuation linéique $\mu$ est plus grand.
\end{enumerate}
\end{exercice}

\begin{exercice}[Définitions et vocabulaire]
\begin{enumerate}
\item Définir les termes suivants : \cle{transmutation}, \cle{fission}, \cle{fusion}, \cle{réaction en chaîne}.
\item Qu'appelle-t-on \cle{dose absorbée} et \cle{dose équivalente} ? Préciser les unités de chacune.
\item Donner la signification physique du coefficient d'atténuation linéique $\mu$ et celle de la couche de demi-atténuation $x_{1/2}$.
\end{enumerate}
\end{exercice}

\begin{exercice}[Calculs directs]
Données : $c = \SI{3,00e8}{m.s^{-1}}$ ; $\SI{1}{MeV} = \SI{1,60e-13}{J}$ ; $Q(\alpha) = 20$.
\begin{enumerate}
\item Un défaut de masse de $\Delta m = \SI{0,20}{u}$ accompagne une fission. Sachant que $\SI{1}{u} \cdot c^2 = \SI{931,5}{MeV}$, calculer l'énergie libérée en MeV puis en joules.
\item Un organe absorbe une dose de $\SI{0,02}{Gy}$ de particules $\alpha$. Calculer la dose équivalente reçue en sieverts.
\item Un écran a pour coefficient d'atténuation $\mu = \SI{0,50}{cm^{-1}}$. Quelle fraction du nombre initial de photons traverse une épaisseur $x = \SI{4,0}{cm}$ ?
\end{enumerate}
\end{exercice}

\courssub{Je m'entraîne}

\begin{exercice}[Équilibrer et identifier]
Équilibrer les équations suivantes en déterminant la particule $X$ (nature, nombre de masse et charge), puis préciser dans chaque cas s'il s'agit d'une fission ou d'une fusion.
\begin{enumerate}
\item $\ce{^{235}_{92}U + ^{1}_{0}n -> ^{90}_{38}Sr + X + 3^{1}_{0}n}$
\item $\ce{^{2}_{1}H + ^{3}_{1}H -> ^{4}_{2}He + X}$
\item $\ce{^{235}_{92}U + ^{1}_{0}n -> ^{139}_{54}Xe + ^{95}_{38}Sr + X}$ où $X$ est un nombre entier de neutrons.
\end{enumerate}
\end{exercice}

\begin{exercice}[Énergie d'une fission et comparaison avec le pétrole]
Données : $\SI{1}{u} \cdot c^2 = \SI{931,5}{MeV}$ ; $\SI{1}{MeV} = \SI{1,60e-13}{J}$ ; $N_A = \SI{6,02e23}{mol^{-1}}$ ; masse molaire de l'uranium-235 : $M = \SI{235}{g.mol^{-1}}$ ; pouvoir calorifique du pétrole : $\SI{4e7}{J}$ pour $\SI{1}{kg}$.
\begin{enumerate}
\item Le défaut de masse lors de la fission d'un noyau d'uranium-235 est $\Delta m = \SI{0,193}{u}$. Calculer l'énergie libérée par la fission d'un noyau, en MeV puis en joules.
\item Calculer le nombre de noyaux contenus dans $\SI{1}{kg}$ d'uranium-235, puis l'énergie totale libérée par la fission complète de ce kilogramme.
\item Quelle masse de pétrole faudrait-il brûler pour obtenir la même énergie ? Conclure.
\end{enumerate}
\end{exercice}

\begin{exercice}[La masse perdue par le Soleil]
Données : puissance rayonnée par le Soleil $P = \SI{3,9e26}{W}$ ; $c = \SI{3,00e8}{m.s^{-1}}$ ; 1 an $\approx \SI{3,16e7}{s}$.
\begin{enumerate}
\item Rappeler la relation d'équivalence masse-énergie d'Einstein.
\item Calculer l'énergie rayonnée par le Soleil pendant une seconde.
\item En déduire la masse perdue par le Soleil chaque seconde, puis chaque année. Commenter l'ordre de grandeur obtenu.
\end{enumerate}
\end{exercice}

\begin{exercice}[Écran de plomb et couches de demi-atténuation]
Un faisceau de photons traverse un écran de plomb de coefficient d'atténuation linéique $\mu = \SI{1,2}{cm^{-1}}$.
\begin{enumerate}
\item Démontrer, à partir de la loi $N(x) = N_0\,e^{-\mu x}$, la relation $x_{1/2} = \dfrac{\ln 2}{\mu}$.
\item Calculer $x_{1/2}$ pour le plomb.
\item Déterminer l'épaisseur de plomb nécessaire pour que seulement $1\,\%$ des photons incidents soient transmis.
\item Retrouver ce résultat en utilisant les couches de demi-atténuation (on remarquera que $1\,\% \approx \left(\tfrac{1}{2}\right)^{6,64}$).
\item Calculer $x_{1/2}$ pour le béton ($\mu = \SI{0,20}{cm^{-1}}$) et comparer l'efficacité des deux matériaux.
\end{enumerate}
\end{exercice}

\courssub{Je me prépare au Bac}

\begin{exercice}[Irradiation, contamination et dosimétrie en milieu hospitalier]
Dans un service de radiothérapie d'un hôpital de Yaoundé, on utilise une source radioactive pour traiter des tumeurs.
\begin{enumerate}
\item Définir l'\cle{irradiation} et la \cle{contamination} radioactive, en donnant un exemple concret de chaque situation.
\item Citer les trois moyens fondamentaux de radioprotection et expliquer brièvement le principe physique de chacun.
\item Un technicien reçoit au niveau de la main une dose absorbée $D = \SI{0,02}{Gy}$ due à des particules $\alpha$ de facteur de qualité $Q = 20$. Calculer la dose équivalente $H$ reçue, en sieverts. Commenter sachant que la dose limite annuelle réglementaire pour le public est de $\SI{1}{mSv}$.
\item Le faisceau de photons utilisé traverse un écran de protection. Sachant que le coefficient d'atténuation du matériau est $\mu = \SI{0,80}{cm^{-1}}$, calculer l'épaisseur minimale de l'écran pour que le nombre de photons transmis soit inférieur à $0,1\,\%$ du nombre incident.
\end{enumerate}
\end{exercice}

\begin{exercice}[Étude d'une centrale nucléaire]
Une centrale nucléaire fournit une puissance électrique $P_e = \SI{900}{MW}$ avec un rendement global $\eta = 33\,\%$. Chaque fission d'un noyau d'uranium-235 libère en moyenne une énergie $E_1 = \SI{200}{MeV}$.

Données : $\SI{1}{MeV} = \SI{1,60e-13}{J}$ ; $N_A = \SI{6,02e23}{mol^{-1}}$ ; $M(\ce{^{235}U}) = \SI{235}{g.mol^{-1}}$ ; 1 an $\approx \SI{3,16e7}{s}$.
\begin{enumerate}
\item Écrire l'équation d'une fission possible de l'uranium-235 en précisant les lois de conservation utilisées.
\item Calculer la puissance thermique $P_{th}$ fournie par les fissions dans le réacteur.
\item En déduire l'énergie thermique produite en une année de fonctionnement continu.
\item Calculer le nombre de fissions par seconde dans le réacteur, puis le nombre de noyaux fissionnés en un an.
\item En déduire la masse d'uranium-235 consommée par an. Comparer avec la masse de pétrole ($\SI{4e7}{J.kg^{-1}}$) qu'il faudrait brûler pour produire la même énergie thermique, et conclure sur l'intérêt énergétique du nucléaire.
\end{enumerate}
\end{exercice}

\begin{exercice}[Fusion dans les étoiles et projet ITER]
Dans le Soleil, l'énergie provient de la fusion de noyaux d'hydrogène en hélium. Sur Terre, le projet international ITER cherche à reproduire cette réaction avec les isotopes deutérium et tritium :
\[ \ce{^{2}_{1}H + ^{3}_{1}H -> ^{4}_{2}He + ^{1}_{0}n} \]
Données : masses en unités de masse atomique : $m(\ce{^{2}_{1}H}) = \SI{2,0136}{u}$ ; $m(\ce{^{3}_{1}H}) = \SI{3,0155}{u}$ ; $m(\ce{^{4}_{2}He}) = \SI{4,0015}{u}$ ; $m_n = \SI{1,0087}{u}$ ; $\SI{1}{u} \cdot c^2 = \SI{931,5}{MeV}$ ; $\SI{1}{MeV} = \SI{1,60e-13}{J}$.
\begin{enumerate}
\item Vérifier que l'équation de la réaction respecte les lois de conservation de Soddy.
\item Calculer le défaut de masse $\Delta m$ de cette réaction, en unité de masse atomique.
\item En déduire l'énergie libérée par une fusion, en MeV puis en joules.
\item Pourquoi cette réaction nécessite-t-elle des températures de l'ordre de $10^8$ K ? Quel est le rôle du confinement magnétique dans le tokamak d'ITER ?
\item Comparer qualitativement, du point de vue de la radioprotection et des déchets, la fusion et la fission utilisées dans les centrales actuelles.
\end{enumerate}
\end{exercice}



\newpage

\coursec{Corrigés des exercices}

\coursec{Réactions nucléaires provoquées et radioprotection}

\courssub{Réactions nucléaires provoquées : transmutation}

\begin{definition}[Transmutation nucléaire]
Une \cle{transmutation} est la transformation d'un noyau en un noyau d'un \emph{autre élément chimique} (son numéro atomique $Z$ change). Elle peut être :
\begin{itemize}
\item \textbf{spontanée} : radioactivité $\alpha$, $\beta^-$, $\beta^+$ ou capture électronique ;
\item \textbf{provoquée} : le noyau cible est bombardé par une particule incidente (neutron, proton, particule $\alpha$, deutéron...).
\end{itemize}
\end{definition}

\begin{propriete}[Lois de conservation (lois de Soddy) -- Exigibles au Bac]
Dans toute réaction nucléaire (spontanée ou provoquée), se conservent :
\begin{enumerate}
\item Le \textbf{nombre total de nucléons} $A$ (nombre de masse) ;
\item La \textbf{charge électrique totale} $Z$ (nombre atomique).
\end{enumerate}
Ces deux lois permettent d'écrire et d'équilibrer toute équation nucléaire.
\end{propriete}

\begin{methode}[Écrire l'équation d'une réaction nucléaire provoquée]
\begin{enumerate}
\item Identifier le noyau cible $\ce{^{A_1}_{Z_1}X}$ et la particule incidente $\ce{^{a}_{z}x}$ (réactifs, à gauche).
\item Identifier les produits détectés ou attendus (à droite), dont le noyau inconnu $\ce{^{A}_{Z}Y}$.
\item Appliquer la conservation de $A$ : $A_1 + a = \sum A_{\text{produits}}$.
\item Appliquer la conservation de $Z$ : $Z_1 + z = \sum Z_{\text{produits}}$.
\item En déduire $A$ et $Z$ de l'inconnu, puis son symbole chimique via la classification périodique.
\item Vérifier l'équation complète.
\end{enumerate}
\end{methode}

\begin{exemplebox}[Transmutation historique : Rutherford (1919)]
Bombardement d'azote par des particules $\alpha$ :
\[ \ce{^{14}_{7}N + ^{4}_{2}He -> ^{17}_{8}O + ^{1}_{1}p} \]
Vérification : $A : 14+4 = 17+1 = 18$ ; $Z : 7+2 = 8+1 = 9$. C'est la première transmutation artificielle observée.
\end{exemplebox}

\begin{attention}[Piège fréquent : le neutron]
Le neutron $\ce{^{1}_{0}n}$ compte pour $1$ dans le bilan de $A$ et pour $0$ dans celui de $Z$. Ne jamais l'oublier dans les équations de fission ou de capture neutronique.
\end{attention}

\courssub{Fission nucléaire et réaction en chaîne}

\begin{definition}[Fission nucléaire]
La \cle{fission} est la rupture d'un noyau \emph{lourd} (uranium-235, plutonium-239...) en deux noyaux plus légers (produits de fission), provoquée par l'absorption d'un neutron lent (thermique). Elle s'accompagne de l'émission de $2$ à $3$ neutrons rapides et d'une libération d'énergie considérable ($\approx \SI{200}{MeV}$ par fission).
\end{definition}

\begin{exemplebox}[Fission type de l'uranium-235]
\[ \ce{^{235}_{92}U + ^{1}_{0}n -> ^{90}_{38}Sr + ^{143}_{54}Xe + 3\,^{1}_{0}n + E} \]
Vérification Soddy : $A : 235+1 = 90+143+3 = 236$ ; $Z : 92 = 38+54$.
Les produits de fission (Sr, Xe, etc.) sont radioactifs ($\beta^-$ le plus souvent) : ce sont les \emph{déchets à vie longue} (ex. $\ce{^{90}_{38}Sr}$, $T_{1/2}=\SI{28,8}{an}$ ; $\ce{^{137}_{55}Cs}$, $T_{1/2}=\SI{30}{an}$).
\end{exemplebox}

\begin{propriete}[Réaction en chaîne]
Les neutrons émis lors d'une fission peuvent provoquer de nouvelles fissions sur d'autres noyaux $\ce{^{235}_{92}U}$, qui émettent à leur tour des neutrons, etc.
\begin{itemize}
\item \textbf{Non contrôlée} (masse supercritique) : emballement exponentiel $\rightarrow$ \textbf{bombe A}.
\item \textbf{Contrôlée} (réacteur) : on absorbe l'excès de neutrons (barres de contrôle en bore, hafnium, cadmium) pour maintenir exactement \emph{un} neutron par fission qui déclenche la suivante ($k=1$) $\rightarrow$ \textbf{centrale nucléaire}.
\end{itemize}
\end{propriete}

\begin{savaistu}[Le contexte camerounais : énergie nucléaire ?]
Le Cameroun dispose d'un potentiel hydroélectrique majeur (barrages d'\textbf{Edéa}, \textbf{Lom Pangar}, \textbf{Nachtigal}, Memve'ele géré par \textbf{SONATREL}/\textbf{ENEO}). La fission nucléaire n'est pas d'actualité immédiate, mais la maîtrise des concepts (réaction en chaîne, criticité, déchets, radioprotection) est indispensable pour les futurs ingénieurs formés à l'ENSPY ou à la FSM, notamment dans la perspective d'une diversification du mix énergétique ou pour la radiothérapie (hôpital général de Yaoundé, centre de Douala).
\end{savaistu}

\begin{experience}[Modélisation de la réaction en chaîne (analogie)]
\textbf{Matériel} : Plateau percé, billes (neutrons), pièges (barres de contrôle).
\textbf{Protocole} : Lancer une bille ; si elle tombe dans un trou "noyau", elle en éjecte 2 ou 3 nouvelles billes. Compter le nombre de billes actives à chaque "génération".
\textbf{Observation} : Sans pièges, le nombre explose (bombe). Avec des pièges bien placés, le nombre se stabilise (réacteur).
\textbf{Interprétation} : Illustrer le facteur de multiplication $k$ et le rôle des barres de contrôle.
\end{experience}

\courssub{Fusion nucléaire}

\begin{definition}[Fusion nucléaire]
La \cle{fusion} est l'union de deux noyaux \emph{légers} (isotopes de l'hydrogène : deutérium $\ce{^{2}_{1}H}$, tritium $\ce{^{3}_{1}H}$) pour former un noyau plus lourd (hélium $\ce{^{4}_{2}He}$), avec libération d'énergie.
\end{definition}

\begin{propriete}[Conditions de la fusion]
Les noyaux sont chargés positivement : ils se repoussent par force de Coulomb. Pour que l'interaction nucléaire forte (attractive, courte portée $\approx \SI{1}{fm}$) l'emporte, il faut leur donner une énergie cinétique énorme $\rightarrow$ \textbf{températures de l'ordre de $\SI{e8}{K}$}. La matière est alors un \textbf{plasma} (gaz totalement ionisé).
\end{propriete}

\begin{exemplebox}[Réaction D-T (projet ITER)]
\[ \ce{^{2}_{1}H + ^{3}_{1}H -> ^{4}_{2}He + ^{1}_{0}n + \SI{17,6}{MeV}} \]
Vérification : $A: 2+3=4+1=5$ ; $Z: 1+1=2+0=2$.
Le neutron rapide ($\SI{14,1}{MeV}$) dépose son énergie dans la couverture (blanket) pour produire de la chaleur $\rightarrow$ électricité. Le tritium est régénéré in situ par réaction du neutron sur le lithium ($\ce{^{6}_{3}Li + n -> ^{4}_{2}He + ^{3}_{1}H}$).
\end{exemplebox}

\begin{savaistu}[ITER : la voie des étoiles sur Terre]
Le projet international \textbf{ITER} (Cadarahe, France) construit un \textbf{tokamak} : un tore à confinement magnétique. Des champs magnétiques intenses (supraconducteurs Nb$_3$Sn, NbTi) maintiennent le plasma loin des parois. Objectif : produire $\SI{500}{MW}$ thermiques pour $\SI{50}{MW}$ injectés ($Q=10$). Première plasma prévue vers 2025-2030. La fusion ne produit \emph{pas de réaction en chaîne}, son "cendre" (hélium) n'est pas radioactive, et les déchets à vie longue sont minimes (activation des structures par les neutrons).
\end{savaistu}

\courssub{Aspect énergétique : $E = \Delta m\,c^2$}

\begin{propriete}[Équivalence masse-énergie (Einstein, 1905) -- Exigible au Bac]
Toute variation de masse $\Delta m$ d'un système isolé correspond à une énergie $E$ échangée avec l'extérieur :
\[ E = \Delta m \cdot c^2 \]
avec $c = \SI{3,00e8}{m.s^{-1}}$.
\begin{itemize}
\item Si $\Delta m = m_{\text{réactifs}} - m_{\text{produits}} > 0$ (défaut de masse), l'énergie est \textbf{libérée} (exothermique : fission, fusion).
\item Si $\Delta m < 0$, l'énergie doit être \textbf{fournie} (endothermique).
\end{itemize}
\end{propriete}

\begin{aretenir}[Unités pratiques en physique nucléaire]
\begin{itemize}
\item Masse atomique unifiée : $\SI{1}{u} = \SI{1,6605e-27}{kg} \approx \SI{931,5}{MeV/c^2}$.
\item Énergie : $\SI{1}{eV} = \SI{1,602e-19}{J}$ ; $\SI{1}{MeV} = \SI{1,602e-13}{J}$.
\item Puissance : $\SI{1}{W} = \SI{1}{J.s^{-1}}$.
\end{itemize}
\end{aretenir}

\begin{methode}[Calculer l'énergie libérée par une réaction nucléaire]
\begin{enumerate}
\item Écrire l'équation équilibrée (lois de Soddy).
\item Rechercher les masses atomiques des réactifs et des produits (en $u$ ou en $\si{MeV/c^2}$).
\item Calculer le défaut de masse : $\Delta m = \sum m_{\text{réactifs}} - \sum m_{\text{produits}}$.
\item Appliquer $E = \Delta m \cdot c^2$ (en utilisant $\SI{1}{u} = \SI{931,5}{MeV}$).
\item Convertir en joules si nécessaire ($1\,\text{MeV} = 1,60\times 10^{-13}\,\text{J}$).
\end{enumerate}
\end{methode}

\begin{exemplebox}[Comparaison fission / combustion]
\begin{itemize}
\item Fission d'\textbf{1 noyau} $\ce{^{235}U}$ : $\approx \SI{200}{MeV} = \SI{3,2e-11}{J}$.
\item Combustion d'\textbf{1 atome} de carbone (charbon/pétrole) : $\approx \SI{4}{eV} = \SI{6,4e-19}{J}$.
\item Rapport $\approx 5 \times 10^7$ : l'énergie nucléaire est \textbf{50 millions de fois} plus concentrée que l'énergie chimique.
\end{itemize}
\end{exemplebox}

\courssub{Effets biologiques des rayonnements et dosimétrie}

\begin{definition}[Irradiation vs Contamination -- Distinction cruciale]
\begin{itemize}
\item \textbf{Irradiation} : exposition au rayonnement d'une source \emph{externe}. Elle \textbf{cèsse} dès qu'on s'éloigne, qu'on éteint la source ou qu'on s'écranne. \emph{Exemple : patient sous le faisceau du linéaire de radiothérapie (hôpital de Yaoundé).}
\item \textbf{Contamination} : présence de substances radioactives \emph{sur} le corps (peau, vêtements) ou \emph{dans} l'organisme (inhalation, ingestion, plaie). Elle \textbf{persiste} tant que l'élément radioactif n'est pas éliminé (désincorporation, décroissance radioactive), même loin de toute source. \emph{Exemple : technicien manipulant une source non scellée sans gants.}
\end{itemize}
\end{definition}

\begin{propriete}[Grandeurs dosimétriques -- Notions exigibles]
\begin{itemize}
\item \textbf{Dose absorbée} $D$ : énergie déposée par unité de masse.
  \[ D = \frac{E}{m} \quad \text{unité : \textbf{gray (Gy)}, } \SI{1}{Gy} = \SI{1}{J.kg^{-1}}. \]
\item \textbf{Dose équivalente} $H$ : pondère $D$ par la dangerosité biologique du rayonnement (facteur de qualité $Q$).
  \[ H = Q \times D \quad \text{unité : \textbf{sievert (Sv)}}. \]
  \begin{center}
  \begin{tabular}{|c|c|}\hline
  Rayonnement & $Q$ \\ \hline
  Photons ($\gamma$, X), électrons ($\beta$) & 1 \\ \hline
  Neutrons (énergie thermique à rapide) & 5 à 20 \\ \hline
  Particules $\alpha$, fragments de fission & 20 \\ \hline
  \end{tabular}
  \end{center}
\item \textbf{Activité} $A$ : nombre de désintégrations par seconde. Unité : \textbf{becquerel (Bq)}, $\SI{1}{Bq} = \SI{1}{s^{-1}}$.
\end{itemize}
\end{propriete}

\begin{attention}[Ne pas confondre les unités !]
\begin{itemize}
\item \textbf{Bq} : activité de la \emph{source} (combien de désintégrations/s).
\item \textbf{Gy} : énergie absorbée par la \emph{cible} (dose physique).
\item \textbf{Sv} : risque biologique pour l'\emph{organisme} (dose équivalente).
\end{itemize}
Au Bac, on attend la formule $H = Q \times D$ et la conversion $\SI{1}{Sv} = \SI{1000}{mSv}$.
\end{attention}

\begin{savaistu}[Limites réglementaires (radioprotection)]
\begin{itemize}
\item Public : $\SI{1}{mSv.an^{-1}}$ (dose équivalente efficace).
\item Travailleurs exposés (catégorie A) : $\SI{20}{mSv.an^{-1}}$ (sur 12 mois glissants), $\SI{100}{mSv}$ sur 5 ans.
\item Ces limites visent à maintenir le risque stochastique (cancer, leucémie) à un niveau "aussi bas que raisonnablement possible" (principe ALARA).
\end{itemize}
\end{savaistu}

\courssub{Radioprotection et atténuation d'un faisceau de photons}

\begin{propriete}[Trois principes fondamentaux de radioprotection]
\begin{enumerate}
\item \textbf{La distance} : pour une source ponctuelle, le débit de dose suit une loi en $1/d^2$. Doubler la distance divise la dose par 4.
\item \textbf{Le temps} : la dose reçue est proportionnelle à la durée d'exposition $t$. Réduire $t$ au strict nécessaire.
\item \textbf{L'écran} : interposer un matériau absorbant (plomb, béton, acier, eau) entre la source et l'opérateur.
\end{enumerate}
\end{propriete}

\begin{propriete}[Loi d'atténuation d'un faisceau de photons -- Exigible au Bac]
Pour un faisceau de photons monoénergétiques traversant un écran d'épaisseur $x$ :
\[ N(x) = N_0 \, e^{-\mu x} \]
où :
\begin{itemize}
\item $N_0$ : nombre de photons incidents ;
\item $N(x)$ : nombre de photons transmis ;
\item $\mu$ : \textbf{coefficient d'atténuation linéique} (en $\si{m^{-1}}$ ou $\si{cm^{-1}}$), dépend du matériau et de l'énergie des photons.
\end{itemize}
\end{propriete}

\begin{propriete}[Couche de demi-atténuation $x_{1/2}$ -- Exigible au Bac]
Épaisseur qui divise par deux le flux de photons :
\[ N(x_{1/2}) = \frac{N_0}{2} \;\Longrightarrow\; x_{1/2} = \frac{\ln 2}{\mu} \approx \frac{0,693}{\mu}. \]
\end{propriete}

\begin{methode}[Calculer l'épaisseur d'écran pour une atténuation donnée]
\begin{enumerate}
\item Identifier $\mu$ (donné ou à chercher dans un tableau) et le facteur de transmission recherché $T = N(x)/N_0$.
\item Utiliser $T = e^{-\mu x} \;\Longrightarrow\; x = -\dfrac{\ln T}{\mu} = \dfrac{\ln(1/T)}{\mu}$.
\item Ou raisonner en nombre de couches de demi-atténuation $n = \log_2(1/T)$, puis $x = n \times x_{1/2}$.
\item Vérifier l'homogénéité des unités ($\mu$ en $\si{cm^{-1}} \Rightarrow x$ en $\si{cm}$).
\end{enumerate}
\end{methode}

\begin{exemplebox}[Écran de plomb pour radiologie]
Pour des photons de $\SI{100}{keV}$, $\mu_{\text{Pb}} \approx \SI{60}{cm^{-1}}$.
$x_{1/2} = 0,693 / 60 \approx \SI{0,012}{cm} = \SI{0,12}{mm}$.
Un tablier plombé de $\SI{0,5}{mm}$ d'équivalent plomb correspond à $\approx 4$ couches de demi-atténuation $\rightarrow$ transmission $1/16 \approx 6\,\%$.
Pour le béton ($\mu \approx \SI{0,2}{cm^{-1}}$), $x_{1/2} \approx \SI{3,5}{cm}$ : il faut des murs épais (salles de radiothérapie, bunkers).
\end{exemplebox}

\begin{attention}[Conditions d'application de la loi $N=N_0 e^{-\mu x}$]
\begin{itemize}
\item Faisceau \textbf{monoénergétique} (ou $\mu$ efficace pour un spectre).
\item Géométrie \textbf{étroite} (pas de diffusion latérale détectée, "good geometry"). En géométrie large, le rayonnement diffusé augmente le flux détecté $\rightarrow$ facteur de buildup $B > 1$ : $N = B N_0 e^{-\mu x}$.
\item $\mu$ dépend fortement de l'énergie : plus les photons sont énergétiques, plus $\mu$ est faible (sauf pics d'absorption).
\end{itemize}
\end{attention}

% ============================================================
% EXERCICES ET CORRIGÉS
% ============================================================

\courssub{Je vérifie mes connaissances}

\begin{corrige}[Exercice 1 — QCM : réactions nucléaires]
\begin{enumerate}
\item \textbf{Réponse c.} L'identification du noyau $X$ exige les \emph{deux} lois de Soddy : la conservation du nombre de nucléons $235 + 1 = 90 + 143 + 3 = 236$ fixe $A = 143$, et la conservation de la charge $92 = 38 + 54$ fixe $Z = 54$ (xénon). Une seule loi ne suffit pas : connaître $A$ seul laisserait $Z$ indéterminé, et inversement.
\item \textbf{Réponses b, c et d.} La fusion concerne des noyaux \emph{légers} (hydrogène, deutérium, tritium, hélium) : la réponse a est fausse, ce sont les noyaux lourds qui subissent la fission. La fusion est bien la source d'énergie des étoiles (réponse c) et exige des températures de plusieurs millions de kelvins pour que l'agitation thermique vainque la répulsion électrostatique entre noyaux positifs (réponse d).
\item \textbf{Réponse b.} La dose absorbée $D = \dfrac{E}{m}$ s'exprime en gray : $\SI{1}{Gy} = \SI{1}{J.kg^{-1}}$. Le becquerel (Bq) mesure l'activité d'une source (une désintégration par seconde) ; le sievert (Sv) mesure la dose \emph{équivalente} $H = Q \times D$, qui tient compte de la dangerosité biologique du rayonnement.
\item \textbf{Réponse b.} Chaque couche de demi-atténuation divise le nombre de photons par deux. Or $12,5\,\% = \dfrac{1}{8} = \left(\dfrac{1}{2}\right)^3$ : il faut donc $3$ couches de demi-atténuation, soit $x = 3 \times x_{1/2} = 3 \times 2 = \SI{6}{cm}$.
\end{enumerate}
\end{corrige}

\begin{corrige}[Exercice 2 — Vrai ou faux ?]
\begin{enumerate}
\item \textbf{Vrai.} Ce sont les lois de conservation de Soddy : dans toute réaction nucléaire (spontanée ou provoquée), le nombre total de nucléons $A$ et la charge électrique totale $Z$ se conservent. Ce sont elles qui permettent d'écrire et d'équilibrer toute équation nucléaire.
\item \textbf{Faux.} La fission de l'uranium-235 est une réaction \emph{provoquée} : elle est déclenchée par l'absorption d'un neutron lent (neutron thermique) par le noyau. Ce n'est pas une désintégration spontanée comme la radioactivité $\alpha$, $\beta$ ou $\gamma$.
\item \textbf{Faux.} C'est l'inverse : la masse des produits est \emph{inférieure} à la masse des réactifs. Le défaut de masse $\Delta m = m_{\text{réactifs}} - m_{\text{produits}} > 0$ est converti en énergie libérée selon la relation d'Einstein $E = \Delta m\,c^2$. C'est précisément ce défaut de masse qui explique l'énorme dégagement d'énergie.
\item \textbf{Vrai.} L'irradiation est une exposition à une source \emph{externe} : elle cesse dès que la source est éloignée, éteinte ou écran­tée. La contamination correspond à de la matière radioactive déposée \emph{sur} le corps (peau, vêtements) ou \emph{dans} l'organisme (inhalation, ingestion, plaie) : elle persiste tant que la substance radioactive n'a pas été éliminée, même loin de toute source externe.
\item \textbf{Vrai.} D'après la loi d'atténuation $N(x) = N_0\,e^{-\mu x}$, à épaisseur $x$ égale, plus le coefficient d'atténuation linéique $\mu$ est grand, plus l'exponentielle décroît vite et plus le nombre de photons transmis est faible. Le plomb ($\mu$ élevé) est donc un écran plus efficace que le béton à épaisseur égale.
\end{enumerate}
\end{corrige}

\begin{corrige}[Exercice 3 — Définitions et vocabulaire]
\begin{enumerate}
\item \begin{itemize}
\item \cle{Transmutation} : transformation d'un noyau en un noyau d'un \emph{autre élément chimique} (le numéro atomique $Z$ change). Elle peut être spontanée (radioactivité) ou provoquée (bombardement par une particule).
\item \cle{Fission} : rupture d'un noyau \emph{lourd} en deux noyaux plus légers, provoquée par l'absorption d'un neutron, accompagnée de l'émission de deux ou trois neutrons et d'une libération d'énergie.
\item \cle{Fusion} : union de deux noyaux \emph{légers} pour former un noyau plus lourd, avec libération d'énergie ; elle nécessite des températures très élevées.
\item \cle{Réaction en chaîne} : les neutrons émis lors d'une fission provoquent à leur tour de nouvelles fissions, qui émettent de nouveaux neutrons, et ainsi de suite. Contrôlée dans un réacteur, elle est explosive dans une bombe A.
\end{itemize}
\item La \cle{dose absorbée} $D$ est l'énergie absorbée par unité de masse de matière irradiée :
\[ D = \frac{E}{m} \quad \text{en gray (Gy), avec } \SI{1}{Gy} = \SI{1}{J.kg^{-1}}. \]
La \cle{dose équivalente} $H$ pondère la dose absorbée par la dangerosité biologique du rayonnement :
\[ H = Q \times D \quad \text{en sievert (Sv)}, \]
où $Q$ est le facteur de qualité du rayonnement ($Q = 1$ pour les photons et électrons, $Q = 20$ pour les particules $\alpha$).
\item Le \cle{coefficient d'atténuation linéique} $\mu$ (en $\si{m^{-1}}$ ou $\si{cm^{-1}}$) caractérise l'aptitude d'un matériau à atténuer un faisceau de photons : plus $\mu$ est grand, plus l'atténuation est rapide avec l'épaisseur. La \cle{couche de demi-atténuation} $x_{1/2}$ est l'épaisseur de matériau qui divise par deux le nombre de photons transmis :
\[ x_{1/2} = \frac{\ln 2}{\mu}. \]
\end{enumerate}
\end{corrige}

\begin{corrige}[Exercice 4 — Calculs directs]
\begin{enumerate}
\item Énergie libérée :
\[ E = \Delta m \cdot c^2 = 0,20 \times 931,5 = \SI{186,3}{MeV} \approx \SI{1,9e2}{MeV}. \]
Conversion en joules :
\[ E = 186,3 \times 1,60\times 10^{-13} \approx \SI{3,0e-11}{J}. \]
C'est une énergie minuscule à notre échelle, mais considérable pour un seul noyau.
\item Dose équivalente :
\[ H = Q \times D = 20 \times 0,02 = \SI{0,40}{Sv}. \]
\item Fraction transmise :
\[ \frac{N(x)}{N_0} = e^{-\mu x} = e^{-0,50 \times 4,0} = e^{-2} \approx 0,135. \]
Environ $13,5\,\%$ des photons incidents traversent l'écran de $\SI{4,0}{cm}$.
\end{enumerate}
\end{corrige}

\courssub{Je m'entraîne}

\begin{corrige}[Exercice 5 — Équilibrer et identifier]
\textbf{Méthode} : on applique les deux lois de Soddy — conservation du nombre de nucléons $A$ et conservation de la charge $Z$ — entre réactifs et produits.
\begin{enumerate}
\item Conservation de $A$ : $235 + 1 = 90 + A + 3 \Rightarrow A = 236 - 93 = 143$.\\
Conservation de $Z$ : $92 + 0 = 38 + Z + 0 \Rightarrow Z = 54$, qui correspond au xénon (Xe).\\
Donc $X = \ce{^{143}_{54}Xe}$. Un noyau lourd se scinde en deux noyaux plus légers : c'est une \textbf{fission}.
\item Conservation de $A$ : $2 + 3 = 4 + A \Rightarrow A = 1$.\\
Conservation de $Z$ : $1 + 1 = 2 + Z \Rightarrow Z = 0$.\\
Donc $X = \ce{^{1}_{0}n}$ (un neutron). Deux noyaux légers s'unissent : c'est une \textbf{fusion}.
\item Conservation de $Z$ : $92 = 54 + 38 = 92$ : la charge est déjà équilibrée, $X$ ne porte aucune charge, ce sont bien des neutrons.\\
Conservation de $A$ : $235 + 1 = 139 + 95 + n \Rightarrow n = 236 - 234 = 2$.\\
Donc $X = 2\,\ce{^{1}_{0}n}$. C'est une \textbf{fission}.
\end{enumerate}
\begin{attention}[Point de vigilance]
Le neutron $\ce{^{1}_{0}n}$ compte pour $1$ dans la conservation de $A$ et pour $0$ dans celle de $Z$ : ne pas l'oublier dans les bilans, c'est l'erreur la plus fréquente.
\end{attention}
\end{corrige}

\begin{corrige}[Exercice 6 — Énergie d'une fission et comparaison avec le pétrole]
\begin{enumerate}
\item Énergie libérée par la fission d'un noyau :
\[ E_1 = \Delta m \cdot c^2 = 0,193 \times 931,5 \approx \SI{180}{MeV}. \]
En joules :
\[ E_1 = 179,8 \times 1,60\times 10^{-13} \approx \SI{2,88e-11}{J}. \]
\item Nombre de noyaux dans $m = \SI{1000}{g}$ d'uranium-235 :
\[ N = \frac{m}{M}\,N_A = \frac{1000}{235} \times 6,02\times 10^{23} \approx \SI{2,56e24}{} \text{ noyaux}. \]
Énergie totale libérée :
\[ E = N \times E_1 = 2,56\times 10^{24} \times 2,88\times 10^{-11} \approx \SI{7,4e13}{J}. \]
\item Masse de pétrole équivalente :
\[ m_p = \frac{E}{4\times 10^{7}} = \frac{7,4\times 10^{13}}{4\times 10^{7}} \approx \SI{1,8e6}{kg}, \]
soit environ $\SI{1800}{tonnes}$ de pétrole.\\
\textbf{Conclusion} : $\SI{1}{kg}$ d'uranium-235 libère autant d'énergie qu'environ deux millions de kilogrammes de pétrole. Le rapport est de l'ordre de $10^6$ : l'énergie nucléaire est extraordinairement concentrée, ce qui explique à la fois son intérêt (peu de combustible) et sa dangerosité (réaction en chaîne, déchets radioactifs).
\end{enumerate}
\end{corrige}

\begin{corrige}[Exercice 7 — La masse perdue par le Soleil]
\begin{enumerate}
\item Relation d'équivalence masse-énergie d'Einstein :
\[ E = \Delta m\,c^2, \]
toute variation de masse $\Delta m$ d'un système correspond à une énergie $E$ échangée avec l'extérieur, et réciproquement.
\item Énergie rayonnée en une seconde (par définition de la puissance, $P = \dfrac{E}{\Delta t}$) :
\[ E = P \times \Delta t = 3,9\times 10^{26} \times 1 = \SI{3,9e26}{J}. \]
\item Masse perdue par seconde :
\[ \Delta m = \frac{E}{c^2} = \frac{3,9\times 10^{26}}{(3,00\times 10^{8})^2} = \frac{3,9\times 10^{26}}{9,0\times 10^{16}} \approx \SI{4,3e9}{kg}, \]
soit environ $4$ millions de tonnes chaque seconde.\\
Masse perdue par an :
\[ \Delta m_{an} = 4,3\times 10^{9} \times 3,16\times 10^{7} \approx \SI{1,4e17}{kg}. \]
\textbf{Commentaire} : cette perte est colossale à l'échelle humaine, mais elle reste infime devant la masse du Soleil ($\approx \SI{2e30}{kg}$) : il faudrait des milliards d'années pour qu'elle devienne sensible. Le Soleil peut donc briller ainsi encore très longtemps.
\end{enumerate}
\end{corrige}

\begin{corrige}[Exercice 8 — Écran de plomb et couches de demi-atténuation]
\begin{enumerate}
\item \textbf{Démonstration} (exigible) : par définition de la couche de demi-atténuation, $N(x_{1/2}) = \dfrac{N_0}{2}$. En appliquant la loi d'atténuation :
\[ N_0\,e^{-\mu x_{1/2}} = \frac{N_0}{2} \;\Longrightarrow\; e^{-\mu x_{1/2}} = \frac{1}{2} \;\Longrightarrow\; -\mu x_{1/2} = \ln\frac{1}{2} = -\ln 2, \]
d'où
\[ \boxed{x_{1/2} = \frac{\ln 2}{\mu}}. \]
\item Application numérique pour le plomb :
\[ x_{1/2} = \frac{0,693}{1,2} \approx \SI{0,58}{cm}. \]
\item On veut $\dfrac{N(x)}{N_0} = 0,01$, donc $e^{-\mu x} = 0,01$, soit :
\[ x = \frac{\ln 100}{\mu} = \frac{4,61}{1,2} \approx \SI{3,8}{cm}. \]
\item Avec les couches de demi-atténuation : transmettre $1\,\%$ exige $6,64$ divisions par deux, donc :
\[ x = 6,64 \times x_{1/2} = 6,64 \times 0,58 \approx \SI{3,8}{cm}. \]
On retrouve le même résultat : les deux méthodes sont équivalentes.
\item Pour le béton :
\[ x_{1/2} = \frac{0,693}{0,20} \approx \SI{3,5}{cm}, \]
soit environ $6$ fois plus que pour le plomb. À protection égale, il faut un mur de béton beaucoup plus épais qu'une plaque de plomb : le plomb est un écran bien plus efficace, d'où son usage systématique en radiologie et radiothérapie (tabliers plombés, portes plombées des salles de radiographie).
\end{enumerate}
\end{corrige}

\courssub{Je me prépare au Bac}

\begin{corrige}[Exercice 9 — Irradiation, contamination et dosimétrie en milieu hospitalier]
\begin{enumerate}
\item L'\cle{irradiation} est l'exposition du corps au rayonnement émis par une source \emph{externe} ; elle cesse dès que la source est éloignée, éteinte ou écran­tée. \emph{Exemple} : un patient recevant le faisceau de photons lors d'une séance de radiothérapie.\\
La \cle{contamination} est la présence de substance radioactive \emph{sur} le corps (peau, vêtements) ou \emph{dans} l'organisme (inhalation, ingestion, plaie) ; elle persiste tant que la substance n'est pas éliminée. \emph{Exemple} : des poussières radioactives déposées sur la blouse d'un technicien.
\item Les trois moyens fondamentaux de radioprotection :
\begin{itemize}
\item \textbf{la distance} : l'exposition diminue fortement lorsqu'on s'éloigne de la source ;
\item \textbf{le temps} : réduire la durée d'exposition réduit proportionnellement la dose reçue ;
\item \textbf{l'écran} : un matériau absorbant (plomb, béton) atténue le faisceau selon la loi $N(x) = N_0\,e^{-\mu x}$.
\end{itemize}
\item Dose équivalente :
\[ H = Q \times D = 20 \times 0,02 = \SI{0,40}{Sv} = \SI{400}{mSv}. \]
Cette dose est $400$ fois supérieure à la limite annuelle réglementaire pour le public ($\SI{1}{mSv}$) : l'exposition est grave, les procédures de protection ont manifestement été défaillantes et un suivi médical s'impose.
\item On exige $\dfrac{N(x)}{N_0} \leq 10^{-3}$, soit $e^{-\mu x} \leq 10^{-3}$, donc :
\[ x \geq \frac{\ln 1000}{\mu} = \frac{6,91}{0,80} \approx \SI{8,6}{cm}. \]
L'écran doit avoir une épaisseur minimale d'environ $\SI{8,6}{cm}$.
\end{enumerate}
\textbf{Barème indicatif (sur 5 pts)} : définitions et exemples (1,5 pt) ; trois moyens de protection avec principes (1 pt) ; calcul de $H$ et commentaire chiffré (1,5 pt) ; calcul de l'épaisseur avec passage au logarithme (1 pt).\\
\textbf{Points de vigilance} : citer le facteur de qualité $Q$ dans le calcul de $H$ ; convertir $\SI{0,40}{Sv}$ en mSv pour comparer à la limite réglementaire ; justifier le sens de l'inégalité (l'atténuation croît avec $x$).
\end{corrige}

\begin{corrige}[Exercice 10 — Étude d'une centrale nucléaire]
\begin{enumerate}
\item Exemple d'équation de fission :
\[ \ce{^{235}_{92}U + ^{1}_{0}n -> ^{90}_{38}Sr + ^{143}_{54}Xe + 3^{1}_{0}n} \]
Vérification par les lois de Soddy : conservation du nombre de nucléons $235 + 1 = 90 + 143 + 3 = 236$ ; conservation de la charge $92 = 38 + 54$.
\item Le rendement global est $\eta = \dfrac{P_e}{P_{th}}$, donc :
\[ P_{th} = \frac{P_e}{\eta} = \frac{900}{0,33} \approx \SI{2,73e3}{MW} = \SI{2,73e9}{W}. \]
\item Énergie thermique produite en un an de fonctionnement continu :
\[ E_{th} = P_{th} \times \Delta t = 2,73\times 10^{9} \times 3,16\times 10^{7} \approx \SI{8,6e16}{J}. \]
\item Énergie libérée par fission :
\[ E_1 = 200 \times 1,60\times 10^{-13} = \SI{3,2e-11}{J}. \]
Nombre de fissions par seconde :
\[ n = \frac{P_{th}}{E_1} = \frac{2,73\times 10^{9}}{3,2\times 10^{-11}} \approx \SI{8,5e19}{} \text{ fissions par seconde}. \]
Nombre de noyaux fissionnés en un an :
\[ N = n \times \Delta t = 8,5\times 10^{19} \times 3,16\times 10^{7} \approx \SI{2,7e27}{} \text{ noyaux}. \]
\item Masse d'uranium-235 consommée par an :
\[ m = \frac{N}{N_A}\times M = \frac{2,7\times 10^{27}}{6,02\times 10^{23}} \times 235 \approx \SI{1,05e6}{g} \approx \SI{1,1}{t}. \]
Masse de pétrole équivalente (même énergie thermique) :
\[ m_p = \frac{E_{th}}{4\times 10^{7}} = \frac{8,6\times 10^{16}}{4\times 10^{7}} \approx \SI{2,2e9}{kg}, \]
soit environ $2,2$ millions de tonnes de pétrole.\\
\textbf{Conclusion} : environ une tonne d'uranium-235 remplace plus de deux millions de tonnes de pétrole. La densité énergétique du combustible nucléaire est considérable, ce qui en fait une piste sérieuse pour la production massive d'électricité — une réflexion utile pour les grands projets énergétiques camerounais, en complément de l'hydroélectricité de Nachtigal, Lom Pangar ou Edéa.
\end{enumerate}
\textbf{Barème indicatif (sur 6 pts)} : équation équilibrée avec lois citées (1 pt) ; $P_{th}$ (1 pt) ; $E_{th}$ (0,5 pt) ; $n$ et $N$ (1,5 pt) ; masse d'uranium (1 pt) ; comparaison et conclusion argumentée (1 pt).\\
\textbf{Points de vigilance} : bien distinguer puissance \emph{électrique} et puissance \emph{thermique} (le rendement divise, ne multiplie pas) ; convertir les MeV en joules avant tout calcul en SI ; convertir la masse molaire en $\si{g.mol^{-1}}$ cohérent avec la masse cherchée ; conclure par une phrase rédigée, pas seulement par un nombre.
\end{corrige}

\begin{corrige}[Exercice 11 — Fusion dans les étoiles et projet ITER]
\begin{enumerate}
\item Vérification des lois de Soddy :
\begin{itemize}
\item conservation du nombre de nucléons : $2 + 3 = 5$ et $4 + 1 = 5$ ;
\item conservation de la charge : $1 + 1 = 2$ et $2 + 0 = 2$.
\end{itemize}
L'équation est correctement équilibrée.
\item Défaut de masse (masse des réactifs moins masse des produits) :
\begin{align*}
\Delta m &= m(\ce{^{2}_{1}H}) + m(\ce{^{3}_{1}H}) - m(\ce{^{4}_{2}He}) - m_n \\
&= 2,0136 + 3,0155 - 4,0015 - 1,0087 \\
&= \SI{0,0189}{u}.
\end{align*}
\item Énergie libérée par une fusion :
\[ E = \Delta m \cdot c^2 = 0,0189 \times 931,5 \approx \SI{17,6}{MeV}, \]
soit en joules :
\[ E = 17,6 \times 1,60\times 10^{-13} \approx \SI{2,8e-12}{J}. \]
\item Les noyaux de deutérium et de tritium sont tous deux chargés positivement : ils se repoussent par interaction électrostatique (force de Coulomb). Pour qu'ils s'approchent suffisamment et que l'interaction nucléaire forte, attractive à très courte distance, prenne le relais et les lie, il faut une agitation thermique extrême, d'où des températures de l'ordre de $10^8$ K. À ces températures, la matière est un \emph{plasma} (gaz totalement ionisé) qu'aucune paroi matérielle ne pourrait contenir sans être détruite : dans le tokamak d'ITER, de puissants champs magnétiques canalisent les particules chargées et maintiennent le plasma confiné loin des parois.
\item Comparaison qualitative :
\begin{itemize}
\item la fusion ne produit \emph{pas de réaction en chaîne} : aucun risque d'emballement incontrôlé du réacteur ;
\item son produit direct, l'hélium $\ce{^{4}_{2}He}$, n'est \emph{pas radioactif} ;
\item les déchets à vie longue sont bien moindres qu'avec la fission, dont les produits (strontium-90, césium-137, etc.) sont très radioactifs et exigent un stockage surveillé pendant des siècles.
\end{itemize}
La fusion est donc, en perspective, une source d'énergie plus sûre et plus propre du point de vue de la radioprotection — si l'on parvient à la maîtriser industriellement.
\end{enumerate}
\textbf{Barème indicatif (sur 6 pts)} : vérification des deux conservations (1 pt) ; défaut de masse avec bonne convention de signe (1 pt) ; énergie en MeV et en joules (1,5 pt) ; explication de la température et du confinement magnétique (1,5 pt) ; comparaison fission/fusion argumentée (1 pt).\\
\textbf{Points de vigilance} : le défaut de masse se calcule toujours $\Delta m = m_{\text{réactifs}} - m_{\text{produits}}$ et doit être \emph{positif} pour une réaction qui libère de l'énergie ; ne pas oublier la masse du neutron dans les produits ; pour la question 4, exiger la mention de la répulsion électrostatique et du rôle de l'interaction forte.
\end{corrige}

\newpage

\coursec{Fiche bilan}

\begin{popfigure}
\begin{tikzpicture}[mindmap, grow cyclic, every node/.style={concept, text=white, font=\small\bfseries},
root concept/.append style={concept color=popdark, font=\large\bfseries},
level 1/.append style={level distance=3.6cm, sibling angle=51.4},
level 1 concept/.append style={font=\footnotesize\bfseries}]
\node[root concept]{Réactions nucléaires provoquées et radioprotection}
child[concept color=popblue]{ node[align=center]{Transmutation\\ $\ce{^{A}_{Z}X} + \text{proj.} \to \dots$} }
child[concept color=poporange]{ node[align=center]{Fission\\ $\ce{^{235}_{92}U} + \ce{^{1}_{0}n}$\\ réaction en chaîne} }
child[concept color=poppurple]{ node[align=center]{Fusion\\ $\ce{^{2}_{1}H} + \ce{^{3}_{1}H}$\\ étoiles, ITER} }
child[concept color=popgreen]{ node[align=center]{Énergie\\ $E = \Delta m\, c^{2}$} }
child[concept color=poppink]{ node[align=center]{Effets biologiques\\ irradiation / contamination\\ doses} }
child[concept color=popgold]{ node[align=center]{Radioprotection\\ distance, temps, écran} }
child[concept color=popmuted]{ node[align=center]{Atténuation\\ $N(x) = N_{0}e^{-\mu x}$\\ $x_{1/2} = \ln 2 / \mu$} };
\end{tikzpicture}
\legende{Carte mentale de la leçon : les sept idées forces à maîtriser.}
\end{popfigure}

\begin{bilan}
\textbf{1. Définitions clés}
\begin{itemize}
\item \cle{Transmutation} : transformation d'un noyau en un autre noyau, provoquée par le bombardement d'un noyau cible par une particule (proton, neutron, particule $\alpha$).
\item \cle{Fission} : cassure d'un noyau lourd ($\ce{^{235}_{92}U}$, $\ce{^{239}_{94}Pu}$) sous l'impact d'un neutron, en deux noyaux plus légers, avec émission de 2 ou 3 neutrons et libération d'énergie.
\item \cle{Réaction en chaîne} : les neutrons émis provoquent de nouvelles fissions ; elle est contrôlée dans une centrale (modérateur, barres de contrôle), explosive dans une bombe A.
\item \cle{Fusion} : union de deux noyaux légers ($\ce{^{2}_{1}H}$, $\ce{^{3}_{1}H}$) en un noyau plus lourd ; elle nécessite une très haute température ($\approx 10^{8}$ K) ; source d'énergie des étoiles, projet ITER.
\item \cle{Irradiation} : exposition à un rayonnement provenant d'une source extérieure ; elle cesse dès que la source est éloignée. \cle{Contamination} : présence de substance radioactive sur ou dans l'organisme ; elle persiste tant que la substance n'est pas éliminée.
\end{itemize}

\textbf{2. Lois et formules à connaître par cœur}
\begin{center}
\begin{tabularx}{\linewidth}{|l|X|}
\hline
\rowcolor{popblueL} \textbf{Loi / formule} & \textbf{Contenu} \\
\hline
Conservation du nombre de nucléons (Soddy) & $\sum A$ identique avant et après la réaction \\
\hline
Conservation de la charge (Soddy) & $\sum Z$ identique avant et après la réaction \\
\hline
Défaut de masse & $\Delta m = \left|m_{\text{produits}} - m_{\text{réactifs}}\right| > 0$ \\
\hline
Énergie libérée & $E_{\text{libérée}} = \Delta m \cdot c^{2}$, avec $c = \num{3,00}\times 10^{8}$ m.s$^{-1}$ \\
\hline
Atténuation d'un faisceau de photons & $N(x) = N_{0}\, e^{-\mu x}$ ; $\mu$ : coefficient d'atténuation linéique (m$^{-1}$) \\
\hline
Couche de demi-atténuation & $x_{1/2} = \dfrac{\ln 2}{\mu} \approx \dfrac{0,693}{\mu}$ \\
\hline
Dose absorbée & $D = \dfrac{E_{\text{absorbée}}}{m}$, en gray (Gy = J.kg$^{-1}$) \\
\hline
Dose équivalente & $H = D \times w_{R}$ (facteur de pondération), en sievert (Sv) \\
\hline
\end{tabularx}
\end{center}

\textbf{3. Méthodes express}
\begin{itemize}
\item \textbf{Équilibrer une réaction nucléaire} : écrire les noyaux avec $A$ et $Z$ ; poser deux équations (somme des $A$, somme des $Z$) ; identifier la particule inconnue ($\ce{^{1}_{0}n}$, $\ce{^{4}_{2}He}$, $\ce{^{0}_{-1}e}$, $\ce{^{0}_{1}e}$, $\ce{^{1}_{1}H}$).
\item \textbf{Calculer une énergie libérée} : calculer $\Delta m$ en kg (1 u $= \num{1,6605}\times 10^{-27}$ kg) ; appliquer $E = \Delta m\, c^{2}$ ; convertir éventuellement en MeV (1 MeV $= \num{1,602}\times 10^{-13}$ J).
\item \textbf{Exploiter l'atténuation} : pour diviser le flux par 2, il faut une épaisseur $x_{1/2}$ ; pour le diviser par $2^{n}$, il faut $n$ couches de demi-atténuation.
\item \textbf{Se protéger} : diminuer le \cle{temps} d'exposition, augmenter la \cle{distance} à la source, interposer un \cle{écran} adapté (papier pour $\alpha$, aluminium pour $\beta$, plomb ou béton pour $\gamma$ et neutrons).
\end{itemize}
\end{bilan}

\begin{aretenir}[Checklist avant le Bac]
\begin{itemize}
\item[$\square$] Je sais définir transmutation, fission et fusion, et donner un exemple d'équation pour chacune.
\item[$\square$] J'applique les deux lois de Soddy (conservation de $A$ et de $Z$) pour équilibrer toute équation nucléaire.
\item[$\square$] Je sais identifier les particules usuelles : $\ce{^{1}_{0}n}$, $\ce{^{4}_{2}He}$, $\ce{^{0}_{-1}e}$, $\ce{^{0}_{1}e}$, $\ce{^{1}_{1}p}$.
\item[$\square$] Je calcule le défaut de masse $\Delta m$ et l'énergie libérée $E = \Delta m\, c^{2}$ avec les bonnes unités (kg, J).
\item[$\square$] Je sais convertir entre joules, MeV et unités de masse atomique.
\item[$\square$] J'explique le principe de la réaction en chaîne et son contrôle dans une centrale nucléaire.
\item[$\square$] Je sais pourquoi la fusion exige des températures très élevées et je cite ITER et les étoiles.
\item[$\square$] Je distingue clairement irradiation et contamination, avec un exemple concret pour chacune.
\item[$\square$] Je connais les unités de dosimétrie : gray (Gy) pour la dose absorbée, sievert (Sv) pour la dose équivalente.
\item[$\square$] Je cite les trois moyens de radioprotection : temps, distance, écran, et je sais quel écran pour quel rayonnement.
\item[$\square$] J'exploite $N(x) = N_{0} e^{-\mu x}$ et $x_{1/2} = \ln 2 / \mu$ pour calculer une épaisseur de protection.
\item[$\square$] Je vérifie toujours l'homogénéité de mes résultats et je donne un nombre raisonnable de chiffres significatifs.
\end{itemize}
\end{aretenir}

\findecours

\end{document}
